\documentclass[11pt,twoside]{report}

% ── Geometry & layout ──
\usepackage[margin=1in]{geometry}
\usepackage{parskip}
\setlength{\parindent}{0pt}

% ── Fonts & encoding ──
\usepackage[T1]{fontenc}
\usepackage{lmodern}
\usepackage{microtype}

% ── Math ──
\usepackage{amsmath,amssymb,amsthm,mathtools,bm}

% ── Algorithms ──
\usepackage[ruled,vlined,linesnumbered]{algorithm2e}

% ── Graphics ──
\usepackage{tikz,pgfplots}
\pgfplotsset{compat=1.18}
\usetikzlibrary{arrows.meta,calc,patterns,decorations.pathreplacing,positioning,shapes.geometric}

% ── Tables & lists ──
\usepackage{booktabs,array,enumitem,longtable}

% ── Boxes ──
\usepackage[most]{tcolorbox}
\newtcolorbox{definitionbox}[1][]{colback=blue!5!white,colframe=blue!50!black,title={#1},fonttitle=\bfseries}
\newtcolorbox{examplebox}[1][]{colback=green!5!white,colframe=green!40!black,title={#1},fonttitle=\bfseries}
\newtcolorbox{warningbox}[1][]{colback=red!5!white,colframe=red!50!black,title={#1},fonttitle=\bfseries}
\newtcolorbox{insightbox}[1][]{colback=orange!5!white,colframe=orange!50!black,title={#1},fonttitle=\bfseries}
\newtcolorbox{stepbox}[1][]{colback=violet!5!white,colframe=violet!50!black,title={#1},breakable,fonttitle=\bfseries}

% ── Theorems ──
\theoremstyle{definition}
\newtheorem{definition}{Definition}[chapter]
\newtheorem{example}{Example}[chapter]
\theoremstyle{plain}
\newtheorem{theorem}{Theorem}[chapter]
\newtheorem{proposition}{Proposition}[chapter]
\theoremstyle{remark}
\newtheorem{remark}{Remark}[chapter]

% ── Hyperlinks ──
\usepackage[colorlinks=true,linkcolor=blue!60!black,citecolor=green!50!black,urlcolor=blue!70!black]{hyperref}

% ── Macros ──
\newcommand{\R}{\mathbb{R}}
\newcommand{\N}{\mathbb{N}}
\newcommand{\x}{\bm{x}}
\newcommand{\uu}{\bm{u}}
\newcommand{\f}{\bm{f}}
\newcommand{\q}{\bm{q}}
\newcommand{\pp}{\bm{p}}
\newcommand{\lam}{\bm{\lambda}}
\newcommand{\rr}{\bm{r}}
\newcommand{\0}{\bm{0}}
\newcommand{\I}{\bm{I}}
\newcommand{\M}{\bm{M}}
\newcommand{\J}{\bm{J}}
\newcommand{\Kc}{\mathcal{K}}
\DeclareMathOperator*{\argmin}{arg\,min}
\DeclareMathOperator*{\argmax}{arg\,max}
\DeclareMathOperator{\diag}{diag}
\DeclareMathOperator{\proj}{proj}
\DeclareMathOperator{\dist}{dist}
\DeclareMathOperator{\tr}{tr}
\DeclareMathOperator{\subjectto}{subject\;to}

\title{\Huge\bfseries Trajectory Optimization and Contact Dynamics\\[6pt]
\Large A Comprehensive Tutorial\\[3pt]
\large From Mathematical Foundations to MuJoCo MPC}
\author{Robotics Control Tutorial Series}
\date{2026}

\begin{document}
\maketitle
\tableofcontents

%=============================================================================
% PART I: MATHEMATICAL FOUNDATIONS
%=============================================================================
\part{Mathematical Foundations}

%-----------------------------------------------------------------------------
\chapter{Optimization Fundamentals}
\label{ch:opt-fundamentals}

Before we can discuss trajectory optimization or contact dynamics,
we must establish the mathematical language of optimisation.
Every planner, every contact solver, and every learning algorithm
ultimately solves (or approximates the solution of) an optimisation
problem.

%-------------------------------------------------------------
\section{The General Optimisation Problem}

\begin{definitionbox}[General Optimisation Problem]
Given a \emph{decision variable} $\bm{z}\in\R^{n}$,
an \emph{objective function} $J:\R^{n}\to\R$, and constraint
functions $\bm{g}:\R^{n}\to\R^{m}$, $\bm{h}:\R^{n}\to\R^{p}$,
\begin{equation}\label{eq:general-opt}
\min_{\bm{z}}\; J(\bm{z})
\quad\text{subject to}\quad
\bm{g}(\bm{z})\le \0,\quad
\bm{h}(\bm{z})=\0,\quad
\bm{z}_{\text{lo}}\le\bm{z}\le\bm{z}_{\text{hi}}.
\end{equation}
\end{definitionbox}

\textbf{Terminology.}
A point $\bm{z}$ that satisfies all constraints is called \emph{feasible}.
The set of all feasible points is the \emph{feasible set} $\mathcal{F}$.
A feasible $\bm{z}^{\star}$ is \emph{(locally) optimal} if
$J(\bm{z}^{\star})\le J(\bm{z})$
for all feasible $\bm{z}$ in some neighbourhood.
If this holds for \emph{all} feasible $\bm{z}$, the minimum is \emph{global}.

\paragraph{How to read the standard form.}
Equation~\eqref{eq:general-opt} is worth internalising, because
\emph{every} problem in this tutorial will be forced into this shape.
Three reading conventions:
\begin{enumerate}
\item \textbf{Minimisation is without loss of generality.}
      Maximising $J$ is the same as minimising $-J$.
      We always minimise; a ``reward'' in reinforcement learning is just
      a negated cost.
\item \textbf{Inequalities point one way.} Any inequality can be written
      as $g(\bm{z})\le 0$: the constraint $z_{1}+z_{2}\ge 1$ becomes
      $1-z_{1}-z_{2}\le 0$. Getting the sign right matters later,
      because the sign convention determines the sign of the
      Lagrange multipliers.
\item \textbf{Equalities are two inequalities}, but we keep them separate
      because they are handled differently by algorithms
      (an equality is \emph{always} active; an inequality may or may not be).
\end{enumerate}

\paragraph{The two questions every algorithm must answer.}
All optimisation algorithms, from gradient descent to interior-point
solvers to the contact solvers of Chapter~\ref{ch:contact},
are iterative procedures that repeatedly answer:
\begin{enumerate}
\item \emph{Which direction should I move?} (a \emph{search direction}
      $\bm{d}_{k}$, usually derived from a local model of $J$), and
\item \emph{How far should I move?} (a \emph{step size} $\alpha_{k}$,
      usually chosen by a \emph{line search} or trust region).
\end{enumerate}
producing iterates $\bm{z}_{k+1}=\bm{z}_{k}+\alpha_{k}\bm{d}_{k}$.
The local model is almost always a truncated Taylor expansion;
the entire algorithmic zoo (gradient descent, Newton, Gauss--Newton,
DDP, iLQR) differs only in \emph{how many Taylor terms} are kept and
\emph{how the resulting model is minimised}.
Keep this two-question framing in mind: it converts a bewildering
collection of named methods into small variations on one idea.

\begin{insightbox}[The Taylor expansion is the engine of everything]
For a smooth $J$, the expansion around a point $\bm{z}_{k}$ is
\begin{equation*}
J(\bm{z}_{k}+\bm{\delta})
\;\approx\;
\underbrace{J(\bm{z}_{k})}_{\text{constant}}
+\underbrace{\nabla J(\bm{z}_{k})^{\!\top}\bm{\delta}}_{\text{linear model}}
+\underbrace{\tfrac12\,\bm{\delta}^{\!\top}\nabla^{2}J(\bm{z}_{k})\,\bm{\delta}}_{\text{quadratic model}}.
\end{equation*}
Truncating after the linear term and stepping opposite the gradient
gives \emph{gradient descent}.
Truncating after the quadratic term and jumping to the minimiser of the
quadratic gives \emph{Newton's method}.
Doing this \emph{along an entire trajectory}, time step by time step,
gives \emph{DDP/iLQR} (Chapter~\ref{ch:ddp}).
When you see a new algorithm, your first question should always be:
\emph{what local model is it building, and how is it minimising that model?}
\end{insightbox}

%-------------------------------------------------------------
\section{Unconstrained Optimisation and Optimality Conditions}
\label{sec:unconstrained}

Consider first the unconstrained case: $\min_{\bm{z}} J(\bm{z})$ with no
constraints.

\begin{definitionbox}[First-Order Necessary Condition]
If $\bm{z}^{\star}$ is a local minimum of $J$, and $J$ is differentiable at $\bm{z}^{\star}$, then the gradient vanishes:
\begin{equation}\label{eq:fonc}
\nabla J(\bm{z}^{\star}) = \0.
\end{equation}
Any point satisfying \eqref{eq:fonc} is called a \emph{stationary point}.
It could be a minimum, maximum, or saddle point.
\end{definitionbox}

\begin{definitionbox}[Second-Order Sufficient Condition]
If additionally the Hessian $\bm{H}=\nabla^{2}J(\bm{z}^{\star})$ is
\emph{positive definite} ($\bm{H}\succ 0$), then $\bm{z}^{\star}$ is a
strict local minimum.
\end{definitionbox}

\paragraph{Why these conditions are true.}
Both conditions fall straight out of the Taylor expansion, and the
argument is worth seeing once because it is the template for every
optimality condition in this tutorial.
Suppose $\nabla J(\bm{z}^{\star})\ne\0$.
Take a small step $\bm{\delta}=-\epsilon\,\nabla J(\bm{z}^{\star})$
in the negative gradient direction. Then
\begin{equation*}
J(\bm{z}^{\star}+\bm{\delta})
\approx J(\bm{z}^{\star})
-\epsilon\,\|\nabla J(\bm{z}^{\star})\|^{2}
< J(\bm{z}^{\star})
\quad\text{for small }\epsilon>0,
\end{equation*}
so $\bm{z}^{\star}$ cannot be a minimum---we just found a strictly
better point next door. Hence at a minimum the gradient \emph{must}
vanish. With the gradient gone, the next Taylor term dominates:
$J(\bm{z}^{\star}+\bm{\delta})\approx J(\bm{z}^{\star})
+\tfrac12\bm{\delta}^{\!\top}\bm{H}\bm{\delta}$.
If $\bm{H}\succ 0$, every direction $\bm{\delta}$ increases the cost
(the quadratic form is a bowl); if $\bm{H}$ has a negative eigenvalue,
the corresponding eigenvector is a downhill escape direction and the
point is a saddle. The Hessian's eigenvalues thus classify stationary
points: all positive $\Rightarrow$ minimum, all negative $\Rightarrow$
maximum, mixed signs $\Rightarrow$ saddle.

\paragraph{Why the gradient is the steepest ascent direction.}
The \emph{directional derivative} of $J$ along a unit vector $\bm{d}$ is
$\nabla J^{\!\top}\bm{d}=\|\nabla J\|\cos\theta$, where $\theta$ is the
angle between $\bm{d}$ and $\nabla J$. This is maximised at
$\theta=0$ ($\bm{d}$ aligned with the gradient) and minimised at
$\theta=\pi$ ($\bm{d}=-\nabla J/\|\nabla J\|$).
So ``move against the gradient'' is not a heuristic---it is literally
the locally fastest way downhill, \emph{measured in the Euclidean norm}.
That last qualifier matters: if we measure step lengths in a different
metric $\|\bm{\delta}\|_{\bm{H}}$, the steepest-descent direction becomes
$-\bm{H}^{-1}\nabla J$, which is exactly the Newton direction.
Newton's method is gradient descent in the geometry defined by the Hessian.

\begin{examplebox}[Scalar quadratic]
Let $J(z)=\tfrac12 a\,z^{2}+b\,z+c$ with $a>0$.
Then $J'(z)=az+b=0\;\Rightarrow\;z^{\star}=-b/a$.
And $J''(z^{\star})=a>0$, confirming a minimum.
\end{examplebox}

\begin{examplebox}[Worked Example 1.1: Minimising a non-convex scalar function]
Find the minimum of $J(z)=z^{4}-4z^{2}+z+4$.

\textbf{Step 1.} Compute the gradient:
$J'(z)=4z^{3}-8z+1.$

\textbf{Step 2.} Set $J'(z)=0$. This cubic has three real roots: $z\approx -1.47,\; 0.13,\; 1.35$.

\textbf{Step 3.} Check the Hessian: $J''(z)=12z^{2}-8$.
\begin{itemize}
\item At $z\approx -1.47$: $J''=12(1.47)^{2}-8\approx 17.9>0$ \checkmark\ local min.
\item At $z\approx 0.13$: $J''=12(0.13)^{2}-8\approx -7.8<0$ \texttimes\ local max.
\item At $z\approx 1.35$: $J''=12(1.35)^{2}-8\approx 13.9>0$ \checkmark\ local min.
\end{itemize}
Evaluating: $J(-1.47)\approx 0.26$, $J(1.35)\approx 1.28$.  Global minimum at $z^{\star}\approx -1.47$.
\end{examplebox}

\begin{center}
\begin{tikzpicture}
\begin{axis}[
  width=10cm, height=6cm,
  xlabel={$z$}, ylabel={$J(z)$},
  domain=-2.5:2.5, samples=100,
  grid=major, grid style={gray!30},
  title={$J(z)=z^4-4z^2+z+4$}
]
\addplot[blue, thick]{x^4-4*x^2+x+4};
\addplot[only marks, mark=*, mark size=3pt, red] coordinates {(-1.47,0.26) (1.35,1.28)};
\addplot[only marks, mark=triangle*, mark size=3pt, orange] coordinates {(0.13,3.93)};
\node[red, right] at (axis cs:-1.47,0.26) {\footnotesize global min};
\node[red, left] at (axis cs:1.35,1.28) {\footnotesize local min};
\node[orange, above] at (axis cs:0.13,3.93) {\footnotesize local max};
\end{axis}
\end{tikzpicture}
\end{center}

\paragraph{Gradient descent.}
The simplest iterative algorithm: starting from an initial guess $\bm{z}_{0}$,
\begin{equation}
\bm{z}_{k+1}=\bm{z}_{k}-\alpha_{k}\,\nabla J(\bm{z}_{k}),
\end{equation}
where $\alpha_{k}>0$ is a step size (learning rate).
Each step moves in the direction of steepest descent.

\begin{examplebox}[Worked Example 1.2: Gradient descent on $J(z)=z^{4}-4z^{2}+z+4$]
Starting from $z_{0}=2$ with step size $\alpha=0.05$:
\begin{center}
\begin{tabular}{ccccc}
\toprule
$k$ & $z_{k}$ & $J(z_{k})$ & $J'(z_{k})$ & $z_{k+1}=z_{k}-\alpha J'$\\
\midrule
0 & 2.000 & 6.000 & 17.000 & 1.150\\
1 & 1.150 & 0.948 & $-2.908$ & 1.295\\
2 & 1.295 & 1.133 & $-1.315$ & 1.361\\
3 & 1.361 & 1.278 & 0.089 & 1.356\\
\bottomrule
\end{tabular}
\end{center}
\begin{warningbox}[Initialization matters!]
Gradient descent found the \emph{local} minimum near $z=1.35$, not the global one at $z\approx -1.47$.
Starting from $z_{0}=-2$ instead would converge to the global minimum.
This is a fundamental issue with non-linear optimisation that carries over to trajectory optimisation:
the quality of the initial guess determines which solution you find.
\end{warningbox}
\end{examplebox}

\paragraph{Choosing the step size: backtracking line search.}
A fixed $\alpha$ is fragile: too large and the iterates overshoot and
diverge; too small and progress crawls. The standard fix is
\emph{backtracking line search} with the \emph{Armijo (sufficient decrease)
condition}: start from a candidate $\alpha=1$ (or some $\alpha_{\max}$) and
halve it until
\begin{equation}\label{eq:armijo}
J(\bm{z}_{k}+\alpha\,\bm{d}_{k})
\;\le\;
J(\bm{z}_{k})+c_{1}\,\alpha\,\nabla J(\bm{z}_{k})^{\!\top}\bm{d}_{k},
\qquad c_{1}\in(0,1),\ \text{typically }10^{-4},
\end{equation}
where $\bm{d}_{k}$ is the search direction.
The right-hand side is the \emph{linear model's prediction} of the new cost,
scaled down by $c_{1}$: we accept any step that achieves at least a fraction
$c_{1}$ of the improvement the linear model promised.
This simple test is everywhere in this tutorial:
the iLQR forward pass (Section~\ref{sec:forward-pass}) is precisely a
parallelised line search with an Armijo-style acceptance test, and
MuJoCo MPC's gradient planner uses a parallel line search over a fixed
set of $\alpha$ values.

\begin{examplebox}[Worked Example: Backtracking on the quartic]
At $z_{0}=2$ for $J(z)=z^{4}-4z^{2}+z+4$: $J(2)=6$, $J'(2)=17$, descent
direction $d=-17$. Try $\alpha=0.1$: $z=2-1.7=0.3$, $J(0.3)\approx 3.95$.
Armijo requires $J\le 6+10^{-4}(0.1)(-17^{2})\approx 5.97$. Since
$3.95\le 5.97$: \emph{accept}.
Had we tried $\alpha=0.2$: $z=2-3.4=-1.4$, $J(-1.4)\approx 0.60$
---also accepted, and luckily lands in the global basin.
With $\alpha=0.3$: $z=-3.1$, $J(-3.1)\approx 54.8>6$: \emph{reject},
halve, retry. The test costs one function evaluation per trial and
guards against catastrophic overshoot.
\end{examplebox}

\paragraph{Convergence rates: how fast is fast?}
Iterative methods are classified by how quickly the error
$e_{k}=\|\bm{z}_{k}-\bm{z}^{\star}\|$ shrinks:
\begin{itemize}
\item \textbf{Linear:} $e_{k+1}\le\rho\,e_{k}$ with $\rho<1$.
      Each iteration adds a roughly \emph{constant number of correct digits}.
      Gradient descent is linear, with
      $\rho\approx(\kappa-1)/(\kappa+1)$ where
      $\kappa=\lambda_{\max}(\bm{H})/\lambda_{\min}(\bm{H})$ is the
      \emph{condition number} of the Hessian.
      For $\kappa=100$, $\rho\approx 0.98$: hundreds of iterations
      per digit. This is why naive gradient descent is rarely used for
      trajectory optimisation.
\item \textbf{Superlinear:} $e_{k+1}/e_{k}\to 0$. Quasi-Newton methods
      (BFGS) and iLQR achieve this.
\item \textbf{Quadratic:} $e_{k+1}\le C\,e_{k}^{2}$.
      The number of correct digits \emph{doubles} each iteration:
      errors of $10^{-1}\to 10^{-2}\to 10^{-4}\to 10^{-8}$.
      Newton's method and full DDP are quadratic near the solution.
\end{itemize}

\begin{insightbox}[Conditioning, or why gradient descent zig-zags]
Picture the contours of a quadratic $J=\tfrac12\bm{z}^{\!\top}\bm{H}\bm{z}$
with $\kappa\gg 1$: a long, narrow valley (an elongated ellipse).
The gradient is perpendicular to the contour lines, so from most points
it points \emph{across} the valley rather than \emph{along} it.
Gradient descent therefore bounces from wall to wall, making tiny progress
along the valley floor. Newton's method fixes this by rescaling each
direction by the inverse curvature: it sees the valley as a perfect circle.
In trajectory optimisation the Hessian is almost always badly conditioned
(early controls influence the whole trajectory; late controls barely
matter), which is the practical reason second-order methods like
iLQR dominate first-order ones.
\end{insightbox}

\paragraph{Newton's method.}
Uses second-order information for faster convergence:
\begin{equation}\label{eq:newton}
\bm{z}_{k+1}=\bm{z}_{k}-\bigl[\nabla^{2}J(\bm{z}_{k})\bigr]^{-1}\nabla J(\bm{z}_{k}).
\end{equation}

\paragraph{Where the formula comes from.}
Build the quadratic Taylor model at $\bm{z}_{k}$:
$m(\bm{\delta})=J(\bm{z}_{k})+\nabla J^{\!\top}\bm{\delta}
+\tfrac12\bm{\delta}^{\!\top}\bm{H}\bm{\delta}$.
Minimise the model exactly by setting its gradient to zero:
$\nabla_{\bm{\delta}}m=\nabla J+\bm{H}\bm{\delta}=\0
\;\Rightarrow\;\bm{\delta}^{\star}=-\bm{H}^{-1}\nabla J$.
Newton's method jumps to the bottom of the local quadratic bowl in one
step. If $J$ \emph{is} quadratic (as in LQR), one step solves the problem
exactly---a fact we will exploit repeatedly.
Near a minimum where $\nabla^{2}J\succ 0$, convergence is quadratic.

\begin{examplebox}[Worked Example 1.2b: Newton on the quartic]
$J(z)=z^{4}-4z^{2}+z+4$, so $J'(z)=4z^{3}-8z+1$ and $J''(z)=12z^{2}-8$.
Starting from $z_{0}=2$:
\begin{center}
\begin{tabular}{ccccc}
\toprule
$k$ & $z_{k}$ & $J'(z_{k})$ & $J''(z_{k})$ & $z_{k+1}=z_{k}-J'/J''$\\
\midrule
0 & 2.0000 & 17.000 & 40.00 & 1.5750\\
1 & 1.5750 & 4.0252 & 21.77 & 1.3901\\
2 & 1.3901 & 0.6258 & 15.19 & 1.3489\\
3 & 1.3489 & 0.0264 & 13.83 & 1.3470\\
4 & 1.3470 & $5\!\times\!10^{-5}$ & 13.77 & 1.3470\\
\bottomrule
\end{tabular}
\end{center}
Note the quadratic convergence signature in $|J'|$:
$4.0\to 0.63\to 0.026\to 5\!\times\!10^{-5}$---the exponent roughly
doubles each step. Compare with the gradient descent table in Worked
Example~1.2, which needed comparable iterations just to get one digit.

\textbf{Failure mode.} Start instead from $z_{0}=0$: there
$J''(0)=-8<0$ (we are near the local \emph{maximum}), and the Newton step
$z_{1}=0-1/(-8)=0.125$ moves \emph{toward} the stationary point at
$z\approx 0.13$---which is a maximum! Newton's method finds
\emph{stationary points}, not minima; with negative curvature it happily
climbs hills. Practical fixes: (i)~check that the Newton direction is a
descent direction ($\nabla J^{\!\top}\bm{\delta}<0$) and fall back to
gradient descent otherwise, or (ii)~\emph{regularise}: replace $\bm{H}$
with $\bm{H}+\mu\I$ for $\mu$ large enough to make it positive definite
(Levenberg--Marquardt). As $\mu\to\infty$ the step direction blends
toward plain gradient descent. This exact trick reappears as the
regularisation of $Q_{\uu\uu}$ in iLQR (Section~\ref{sec:ilqr}).
\end{examplebox}

\paragraph{Gauss--Newton: exploiting least-squares structure.}
A huge fraction of robotics costs are \emph{sums of squares}:
\begin{equation}\label{eq:least-squares}
J(\bm{z})=\tfrac12\|\rr(\bm{z})\|^{2}
=\tfrac12\sum_{i}r_{i}(\bm{z})^{2},
\end{equation}
where $\rr:\R^{n}\to\R^{m}$ is a vector of \emph{residuals}---quantities
that should be small when the task is solved (tracking errors, sensor
mismatches; compare the MJPC cost in Chapter~\ref{ch:mjpc}).
Differentiating with the chain rule (let $\J_{r}=\partial\rr/\partial\bm{z}$):
\begin{align}
\nabla J&=\J_{r}^{\!\top}\rr,\\
\nabla^{2}J&=\underbrace{\J_{r}^{\!\top}\J_{r}}_{\text{Gauss--Newton term}}
+\underbrace{\textstyle\sum_{i}r_{i}\,\nabla^{2}r_{i}}_{\text{second-order term}}.
\end{align}
The \textbf{Gauss--Newton approximation} keeps only the first term:
$\nabla^{2}J\approx\J_{r}^{\!\top}\J_{r}$. Two things make this a great deal:
\begin{enumerate}
\item $\J_{r}^{\!\top}\J_{r}$ is \emph{always} positive semidefinite
      (it is a Gram matrix: $\bm{\delta}^{\!\top}\J_{r}^{\!\top}\J_{r}\bm{\delta}
      =\|\J_{r}\bm{\delta}\|^{2}\ge 0$), so the Newton-with-GN-Hessian step
      is always well-posed and always a descent direction.
      No negative-curvature pathology.
\item It needs only \emph{first} derivatives of the residuals---no
      Hessians of $\rr$ required.
\end{enumerate}
The dropped term $\sum_{i}r_{i}\nabla^{2}r_{i}$ is weighted by the
residuals themselves, so near a good solution (small $\rr$) it is small
and Gauss--Newton converges almost as fast as full Newton.
\textbf{This is precisely the relationship between DDP and iLQR}
(Section~\ref{sec:ilqr}): iLQR is Gauss--Newton applied along a
trajectory, dropping the dynamics second derivatives, whose contribution
is weighted by the value-function gradient $V_{\x}'$ (which plays the role
of the residual: it vanishes at optimality).
Once you have seen this pattern in $\R^{n}$, the iLQR derivation will
feel inevitable rather than mysterious.

%-------------------------------------------------------------
\section{Constrained Optimisation and the KKT Conditions}

When constraints are present, the gradient need not vanish at the optimum.
Instead we have the celebrated \emph{Karush--Kuhn--Tucker} (KKT) conditions.

Consider the equality-constrained problem
\begin{equation}
\min_{\bm{z}}\; J(\bm{z}) \quad\text{s.t.}\quad \bm{h}(\bm{z})=\0.
\end{equation}
We introduce \emph{Lagrange multipliers} $\bm{\nu}\in\R^{p}$ and form the
\emph{Lagrangian}:
\begin{equation}\label{eq:lagrangian}
\mathcal{L}(\bm{z},\bm{\nu})=J(\bm{z})+\bm{\nu}^{\!\top}\bm{h}(\bm{z}).
\end{equation}
A necessary condition for optimality is that both partial derivatives vanish:
\begin{equation}
\nabla_{\bm{z}}\mathcal{L}=\nabla J+\bm{\nu}^{\!\top}\nabla\bm{h}=\0,
\qquad
\nabla_{\bm{\nu}}\mathcal{L}=\bm{h}(\bm{z})=\0.
\end{equation}

\paragraph{The geometric picture (read this twice).}
Why should this strange recipe---add the constraints to the objective
with unknown coefficients, then differentiate---find the optimum?
Picture a single constraint surface $h(\bm{z})=0$ and the contour lines
of $J$. Walking along the surface, you can keep decreasing $J$ as long
as the surface has a direction with a downhill component, i.e.\ as long
as $\nabla J$ has a nonzero component \emph{tangent} to the surface.
You can only be stuck when $\nabla J$ is entirely \emph{perpendicular}
to the surface. But the perpendicular direction to the surface
$h(\bm{z})=0$ is spanned by $\nabla h$. Hence, at a constrained optimum,
\begin{equation*}
\nabla J(\bm{z}^{\star})=-\nu\,\nabla h(\bm{z}^{\star})
\quad\text{for some scalar }\nu,
\end{equation*}
which is exactly the stationarity condition above. The Lagrange
multiplier $\nu$ is simply the unknown \emph{proportionality constant}
between the two gradients. Equivalently: at the optimum, the contour
line of $J$ is \emph{tangent} to the constraint surface.
With $p$ constraints, $\nabla J$ must lie in the span of the $p$
constraint gradients (this requires the constraint gradients to be
linearly independent at $\bm{z}^{\star}$---the standard
\emph{constraint qualification}, LICQ, which holds for all problems
in this tutorial).

\paragraph{Multipliers are prices.}
The multipliers also carry quantitative meaning: if the constraint is
perturbed to $h(\bm{z})=\epsilon$, the optimal cost changes as
$\partial J^{\star}/\partial\epsilon=-\nu^{\star}$.
Each multiplier is the \emph{sensitivity of the optimal cost to relaxing
its constraint}---economists call it a shadow price.
This interpretation will become physical in Part~III:
\textbf{contact forces are exactly the Lagrange multipliers of the
non-penetration constraints}. A large normal force at a contact means
the non-penetration constraint there is ``expensive''---the dynamics
are pushing hard against it. The whole machinery of contact solvers is
machinery for computing Lagrange multipliers.

For inequality constraints $g_{i}(\bm{z})\le 0$ with multipliers $\mu_{i}\ge 0$,
the KKT conditions add \emph{complementary slackness}:
\begin{equation}\label{eq:comp-slack}
\mu_{i}\,g_{i}(\bm{z}^{\star})=0,\quad\forall\,i.
\end{equation}
This says: either the constraint is \emph{active} ($g_{i}=0$) or its multiplier is zero.
Complementarity conditions will reappear throughout this tutorial, most prominently in contact dynamics (Chapter~\ref{ch:contact}).

\paragraph{Why inequality multipliers must be non-negative.}
For an equality constraint the multiplier can have either sign---the
surface pushes back in whichever direction is needed. For an inequality
$g\le 0$, the feasible side is one-sided: the constraint can only
``push'' the solution \emph{into} the feasible region, never pull it out.
Formally, at an active constraint, $-\nabla J=\mu\,\nabla g$ must point
\emph{out of} the feasible set (otherwise we could move into the
interior and decrease $J$, contradiction), forcing $\mu\ge 0$.
The physical analogue is immediate: a floor can push up on a block but
cannot pull down; normal contact forces satisfy $\lambda_{N}\ge 0$
\emph{because they are inequality multipliers}.

\begin{definitionbox}[The complete KKT conditions]
For $\min J(\bm{z})$ s.t.\ $\bm{g}(\bm{z})\le\0$, $\bm{h}(\bm{z})=\0$,
a point $\bm{z}^{\star}$ with multipliers $(\bm{\mu}^{\star},\bm{\nu}^{\star})$
is a KKT point if:
\begin{align*}
&\text{(stationarity)} & \nabla J+\textstyle\sum_{i}\mu_{i}^{\star}\nabla g_{i}
+\sum_{j}\nu_{j}^{\star}\nabla h_{j}&=\0,\\
&\text{(primal feasibility)} & \bm{g}(\bm{z}^{\star})\le\0,\quad
\bm{h}(\bm{z}^{\star})&=\0,\\
&\text{(dual feasibility)} & \bm{\mu}^{\star}&\ge\0,\\
&\text{(complementary slackness)} & \mu_{i}^{\star}\,g_{i}(\bm{z}^{\star})&=0
\quad\forall i.
\end{align*}
Under constraint qualification these are \emph{necessary} for a local
minimum; for \emph{convex} problems they are also \emph{sufficient}
for a global minimum. Memorise the four names---every solver log and
every contact-dynamics paper refers to them.
\end{definitionbox}

\begin{insightbox}[Active sets and why inequality constraints are hard]
If an oracle told us \emph{which} inequality constraints are active at
the solution, we could treat those as equalities, discard the rest,
and solve a much easier equality-constrained problem.
The difficulty of inequality-constrained optimisation is entirely the
\emph{combinatorial} problem of discovering the active set: with $m$
inequalities there are $2^{m}$ possible active sets.
Contact dynamics inherits this combinatorics directly---with $n_{c}$
contacts each in \{separating, sticking, sliding\} there are
$3^{n_{c}}$ contact modes, and a contact solver is implicitly an
active-set search. The three main solution strategies
(active-set methods, interior-point methods, and projection/proximal
methods) are surveyed in Section~\ref{sec:solving-constrained},
and each has a direct descendant among the contact solvers of
Chapter~\ref{ch:contact}.
\end{insightbox}

\begin{examplebox}[Worked Example 1.3: 2D constrained optimisation]
\begin{equation*}
\min_{z_1,z_2}\; z_{1}^{2}+z_{2}^{2}
\quad\text{s.t.}\quad z_{1}+z_{2}\ge 1.
\end{equation*}

\textbf{Step 1.} Rewrite: $g(z_{1},z_{2})=1-z_{1}-z_{2}\le 0$.

\textbf{Step 2.} Lagrangian: $\mathcal{L}=z_{1}^{2}+z_{2}^{2}+\mu(1-z_{1}-z_{2})$.

\textbf{Step 3.} KKT:
$\partial\mathcal{L}/\partial z_{1}=2z_{1}-\mu=0\;\Rightarrow\;z_{1}=\mu/2$,\quad
$\partial\mathcal{L}/\partial z_{2}=2z_{2}-\mu=0\;\Rightarrow\;z_{2}=\mu/2$.

\textbf{Step 4.} Complementary slackness: $\mu(1-z_{1}-z_{2})=0$, $\mu\ge 0$.
If $\mu=0$: $z_{1}=z_{2}=0$, but then $z_{1}+z_{2}=0<1$ violates the constraint.
So the constraint is \emph{active}: $z_{1}+z_{2}=1$, giving $\mu/2+\mu/2=1$, thus $\mu=1$.

\textbf{Solution:} $z_{1}^{\star}=z_{2}^{\star}=1/2$, $J^{\star}=1/2$, $\mu^{\star}=1$.

\begin{warningbox}[Connection to contact]
This has \emph{exactly} the same structure as the Signorini contact condition:
either the ``gap'' $g=0$ (constraint active / in contact) or the ``force'' $\mu=0$ (no contact force).
The complementary slackness $\mu\cdot g=0$ is a 1D complementarity condition.
\end{warningbox}
\end{examplebox}

\begin{center}
\begin{tikzpicture}
\begin{axis}[
  width=8cm, height=8cm,
  xlabel={$z_1$}, ylabel={$z_2$},
  domain=-0.5:1.5, samples=30,
  view={0}{90}, colormap/viridis,
  title={Contours of $z_1^2+z_2^2$ with constraint $z_1+z_2\ge 1$}
]
\addplot3[contour prepared={labels=false}, domain=-0.5:1.5, y domain=-0.5:1.5, samples=30]
  {x^2+y^2};
\addplot[red, very thick, domain=-0.5:1.5]{1-x};
\addplot[only marks, mark=*, mark size=4pt, black] coordinates {(0.5,0.5)};
\node[black, above right] at (axis cs:0.5,0.5) {\footnotesize $z^{\star}=(0.5,0.5)$};
\node[red, rotate=-45] at (axis cs:1.1,0.1) {\footnotesize $z_1+z_2=1$};
\end{axis}
\end{tikzpicture}
\end{center}

%-------------------------------------------------------------
\section{How Constrained Problems Are Actually Solved}
\label{sec:solving-constrained}

The KKT conditions characterise the solution but do not provide an
algorithm. In practice, almost every constrained solver converts the
constrained problem into a sequence of \emph{unconstrained} (or simpler)
problems. Four conversion strategies cover essentially everything used
in robotics, and each one foreshadows a contact solver.

\subsection{Penalty Methods}
Replace the constraint with a cost that punishes violation:
\begin{equation}\label{eq:penalty}
\min_{\bm{z}}\;J(\bm{z})+\frac{\rho}{2}\,\|\bm{h}(\bm{z})\|^{2}
+\frac{\rho}{2}\,\|\max(\bm{g}(\bm{z}),\0)\|^{2},
\end{equation}
and solve for increasing penalty weight $\rho\to\infty$.
Simple and robust, but: the solution of the penalised problem only
satisfies the constraints \emph{exactly} in the limit $\rho\to\infty$,
and large $\rho$ makes the Hessian ill-conditioned
($\kappa\propto\rho$), so Newton steps become numerically fragile.
You face a dilemma: accurate constraints or solvable subproblems.

\begin{examplebox}[Worked Example: Quadratic penalty on Example 1.3]
For $\min z_{1}^{2}+z_{2}^{2}$ s.t.\ $z_{1}+z_{2}\ge 1$, the penalised
problem is $\min z_{1}^{2}+z_{2}^{2}+\tfrac{\rho}{2}\max(1-z_{1}-z_{2},0)^{2}$.
By symmetry $z_{1}=z_{2}=z$; in the violated region the stationarity
condition is $2z-\rho(1-2z)=0$, giving $z=\rho/(2+2\rho)$.
\begin{center}
\begin{tabular}{cccc}
\toprule
$\rho$ & $z^{\star}_{\rho}$ & constraint violation $1-2z$ & exact?\\
\midrule
1 & 0.250 & 0.500 & no\\
10 & 0.455 & 0.091 & no\\
100 & 0.495 & 0.010 & no\\
$\infty$ & 0.500 & 0 & limit only\\
\bottomrule
\end{tabular}
\end{center}
The iterate approaches the true solution $z=1/2$ from \emph{inside the
violated region}: penalty methods systematically under-enforce
constraints. The violation $\approx 1/\rho$ is the signature of a
quadratic penalty. MuJoCo's soft contact is exactly this phenomenon:
penetration plays the role of constraint violation, and contact
stiffness plays the role of $\rho$. Finite stiffness $\Leftrightarrow$
finite penetration.
\end{examplebox}

\subsection{The Augmented Lagrangian (Method of Multipliers)}
The fix for the penalty dilemma: keep $\rho$ \emph{moderate}, and
instead carry an explicit running estimate of the multipliers:
\begin{equation}\label{eq:auglag}
\mathcal{L}_{\rho}(\bm{z},\bm{\nu})
=J(\bm{z})+\bm{\nu}^{\!\top}\bm{h}(\bm{z})
+\frac{\rho}{2}\|\bm{h}(\bm{z})\|^{2},
\end{equation}
alternating: (i) minimise $\mathcal{L}_{\rho}$ over $\bm{z}$ with
$\bm{\nu}$ fixed; (ii) update the multiplier estimate
$\bm{\nu}\leftarrow\bm{\nu}+\rho\,\bm{h}(\bm{z})$.
The multiplier update has a beautiful interpretation: whatever residual
constraint violation the penalty failed to remove gets absorbed into the
linear term, so the next subproblem is ``pre-tilted'' to land closer to
feasibility. The method converges with \emph{finite} $\rho$---no
ill-conditioning blow-up. ADMM (Section~\ref{sec:ccp}), the workhorse of
modern contact solvers, is an augmented-Lagrangian method that
additionally splits the variables into blocks solved in alternation.

\subsection{Barrier / Interior-Point Methods}
Penalties punish violation \emph{after} it occurs; barriers make
violation \emph{impossible} by adding a term that blows up at the
boundary of the feasible region:
\begin{equation}\label{eq:barrier}
\min_{\bm{z}}\;J(\bm{z})-\frac{1}{t}\sum_{i}\ln\bigl(-g_{i}(\bm{z})\bigr),
\qquad t>0.
\end{equation}
The log term is finite inside the feasible set and $\to+\infty$ as any
$g_{i}\to 0^{-}$, so the minimiser stays strictly interior---hence
``interior point.'' As $t\to\infty$ the barrier weakens and the solution
traces a smooth \emph{central path} converging to $\bm{z}^{\star}$.
One can show the barrier subproblem's optimality condition reproduces
the KKT system with complementarity relaxed to
$\mu_{i}\,(-g_{i})=1/t$: \textbf{interior-point methods are
complementarity-smoothing methods}. Instead of the hard ``either force
or gap is zero'' they enforce ``force $\times$ gap $=$ small constant.''
IPOPT, the solver behind much contact-implicit trajectory optimisation
(Chapter~\ref{ch:trajopt-contact}), works this way, and MuJoCo's relaxed
contact complementarity has the same flavour: small simultaneous force
and gap, vanishing as the relaxation tightens.

\begin{examplebox}[Worked Example: Log barrier on Example 1.3]
Barrier subproblem:
$\min\;z_{1}^{2}+z_{2}^{2}-\tfrac{1}{t}\ln(z_{1}+z_{2}-1)$.
By symmetry $z_{1}=z_{2}=z$ with $z>1/2$ (strictly feasible).
Stationarity: $2z-\tfrac{1}{t}\cdot\tfrac{1}{2z-1}=0$, i.e.\
$2z(2z-1)=1/t$, giving
$z(t)=\tfrac14\bigl(1+\sqrt{1+4/t}\,\bigr)$.
\begin{center}
\begin{tabular}{ccc}
\toprule
$t$ & $z(t)$ & distance to optimum\\
\midrule
1 & 0.809 & 0.309\\
10 & 0.546 & 0.046\\
100 & 0.5050 & 0.0050\\
1000 & 0.50050 & 0.00050\\
\bottomrule
\end{tabular}
\end{center}
The iterates approach $z^{\star}=1/2$ from \emph{inside} the feasible
region (compare the penalty method, which approached from outside), with
error $\approx 1/(4t)$. The implied multiplier estimate
$\mu(t)=\tfrac{1}{t(2z-1)}\to 1=\mu^{\star}$ recovers the true
multiplier from Worked Example 1.3.
\end{examplebox}

\subsection{Projection Methods}
When the feasible set $\mathcal{C}$ is \emph{simple}---a box, a ball, or
a cone---the cheapest strategy is \emph{projected gradient descent}:
take an unconstrained step, then snap back to the feasible set:
\begin{equation}\label{eq:proj-grad}
\bm{z}_{k+1}=\proj_{\mathcal{C}}\bigl(\bm{z}_{k}-\alpha\nabla J(\bm{z}_{k})\bigr),
\qquad
\proj_{\mathcal{C}}(\bm{y})=\argmin_{\bm{z}\in\mathcal{C}}\|\bm{z}-\bm{y}\|.
\end{equation}
Closed-form projections worth memorising:
\begin{itemize}
\item \textbf{Box} $[\bm{z}_{\text{lo}},\bm{z}_{\text{hi}}]$:
      clamp each coordinate,
      $\proj(\bm{y})_{i}=\min(\max(y_{i},z_{\text{lo},i}),z_{\text{hi},i})$.
\item \textbf{Non-negative orthant} $\{\bm{z}\ge\0\}$:
      $\proj(\bm{y})=\max(\bm{y},\0)$ componentwise. The ubiquitous
      $\max(\cdot,0)$ in contact formulas (e.g.\ the ComFree force
      \eqref{eq:beta-closed}) is literally this projection.
\item \textbf{Second-order cone} $\Kc_{\mu}$: a three-case formula
      (inside: identity; inside the polar cone: zero; otherwise: a
      radial shrink onto the cone surface)---given explicitly in
      Section~\ref{sec:cones}. This is the projection inside the
      ADMM contact solver (Algorithm~\ref{alg:admm}).
\end{itemize}
Projected gradient and its proximal generalisations are the conceptual
core of PGS and ADMM contact solvers: \emph{contact solvers are
projection methods onto friction cones}.

\begin{insightbox}[One table to keep]
\begin{center}
\begin{tabular}{lll}
\toprule
\textbf{Strategy} & \textbf{Constraint handling} & \textbf{Contact-world descendant}\\
\midrule
Penalty & violated, $\to 0$ as $\rho\to\infty$ & soft/compliant contact (MuJoCo spring)\\
Augmented Lagrangian & exact at finite $\rho$ & ADMM contact solvers\\
Barrier / interior point & strictly feasible path & IPOPT contact-implicit trajopt\\
Projection & feasible after each step & PGS, cone projections\\
\bottomrule
\end{tabular}
\end{center}
\end{insightbox}

%-------------------------------------------------------------
\section{Lagrangian Duality}
\label{sec:duality}

One more concept from convex optimisation pays for itself many times
over in Part~III: every minimisation problem has a shadow
\emph{maximisation} problem defined over its multipliers.

Define the \emph{dual function} by minimising the Lagrangian over the
primal variable:
\begin{equation}
d(\bm{\mu},\bm{\nu})=\min_{\bm{z}}\;\mathcal{L}(\bm{z},\bm{\mu},\bm{\nu}),
\end{equation}
and the \emph{dual problem}:
\begin{equation}\label{eq:dual-problem}
\max_{\bm{\mu}\ge\0,\,\bm{\nu}}\;d(\bm{\mu},\bm{\nu}).
\end{equation}

\textbf{Weak duality} (always true): $d(\bm{\mu},\bm{\nu})\le J^{\star}$
for any feasible multipliers---the dual gives lower bounds on the optimal
cost. \textbf{Strong duality} (true for convex problems satisfying a mild
condition such as Slater's): the bound is tight, $d^{\star}=J^{\star}$,
and the dual optimisers are exactly the KKT multipliers.

\begin{examplebox}[Worked Example: The dual of Example 1.3]
$\mathcal{L}=z_{1}^{2}+z_{2}^{2}+\mu(1-z_{1}-z_{2})$.
Minimising over $\bm{z}$ (set gradients to zero):
$z_{1}=z_{2}=\mu/2$, so
\begin{equation*}
d(\mu)=2\cdot\frac{\mu^{2}}{4}+\mu\Bigl(1-\mu\Bigr)
=\mu-\frac{\mu^{2}}{2}.
\end{equation*}
Maximising the concave parabola over $\mu\ge 0$:
$d'(\mu)=1-\mu=0\Rightarrow\mu^{\star}=1$, $d^{\star}=1/2$.
This matches $J^{\star}=1/2$ and $\mu^{\star}=1$ from the primal KKT
solution exactly---strong duality in action. Notice the dual problem is
an \emph{unconstrained-except-for-sign} maximisation of a concave
function: often dramatically simpler than the primal.
\end{examplebox}

\begin{insightbox}[Why duality matters for contact]
The contact problem can be posed in the \emph{primal} (find the
post-contact velocities) or the \emph{dual} (find the contact
impulses/multipliers). MuJoCo solves a primal formulation
\eqref{eq:mujoco-primal}; ADMM and the complementarity-free model of
Jin work on the dual \eqref{eq:ccp-dual-recap}, where the decision
variables \emph{are the contact forces} and the constraint is simply
$\bm{\beta}\ge\0$---an orthant, with a trivial projection.
When you reach Section~\ref{sec:comfree-derivation}, the phrase
``the CCP dual'' will mean precisely the construction above: form the
Lagrangian of the velocity problem, minimise out the velocity
analytically, and maximise what remains over the non-negative impulses.
\end{insightbox}

%-------------------------------------------------------------
\section{Classes of Optimisation Problems}
\label{sec:opt-classes}

Different problem structures admit different algorithms.  In robotics we encounter all of the following.

\subsection{Linear Programs (LP)}
Objective and constraints are all affine:
\begin{equation}
\min_{\bm{z}}\;\bm{c}^{\!\top}\bm{z}\quad\text{s.t.}\quad
\bm{A}\bm{z}\le\bm{b}.
\end{equation}
Solvable in polynomial time via interior-point methods or the simplex method.

\subsection{Quadratic Programs (QP)}
Quadratic objective, linear constraints:
\begin{equation}\label{eq:qp}
\min_{\bm{z}}\;\tfrac12\bm{z}^{\!\top}\bm{Q}\bm{z}+\bm{c}^{\!\top}\bm{z}
\quad\text{s.t.}\quad\bm{A}\bm{z}\le\bm{b}.
\end{equation}
If $\bm{Q}\succeq 0$ (positive semidefinite), the problem is convex and can
be solved efficiently.  QPs are the workhorse of model-predictive control
and arise as sub-problems in many trajectory optimisation algorithms.

\subsection{Second-Order Cone Programs (SOCP)}
\begin{equation}
\min_{\bm{z}}\;\bm{c}^{\!\top}\bm{z}\quad\text{s.t.}\quad
\|\bm{A}_{i}\bm{z}+\bm{b}_{i}\|_{2}\le\bm{d}_{i}^{\!\top}\bm{z}+e_{i},\;
\forall\,i.
\end{equation}
These generalise QPs and LPs.  The friction cone in Coulomb's law is a second-order cone, making SOCPs directly relevant to contact (Section~\ref{sec:coulomb}).

\subsection{Non-linear Programs (NLP)}
The general case~\eqref{eq:general-opt}.  No special structure is assumed.
NLPs are solved by interior-point methods (IPOPT) or sequential quadratic programming (SNOPT).  Trajectory optimisation via direct collocation or shooting produces NLPs.

\begin{insightbox}[Hierarchy of difficulty]
LP $\subset$ QP $\subset$ SOCP $\subset$ SDP $\subset$ Convex $\subset$ NLP.
Moving right, problems become harder.
A central theme of contact modelling (Chapter~\ref{ch:contact}) is that the
``true'' contact problem (an NCP) is non-convex, and various simulators
make different relaxations to recover tractable convex sub-problems.
\end{insightbox}

%-------------------------------------------------------------
\section{Convexity}

Convexity is the dividing line between optimisation problems we can
solve reliably and those we cannot. It comes in two flavours that work
together: convex \emph{sets} (for the feasible region) and convex
\emph{functions} (for the objective).

\begin{definitionbox}[Convex set]
A set $\mathcal{C}\subseteq\R^{n}$ is \emph{convex} if the line segment
between any two of its points stays inside:
$\bm{z}_{1},\bm{z}_{2}\in\mathcal{C}\Rightarrow
\theta\bm{z}_{1}+(1{-}\theta)\bm{z}_{2}\in\mathcal{C}$
for all $\theta\in[0,1]$.
\end{definitionbox}

Examples: boxes, balls, hyperplanes and half-spaces, and---crucially for
us---\emph{cones} like the Coulomb friction cone, and intersections of
any of these. Non-examples: any set with a hole, a notch, or two
disconnected pieces. The set of \emph{collision-free configurations} of
a robot is almost never convex (an obstacle punches a hole in it), which
is one deep reason motion planning is hard.

\begin{definitionbox}[Convex function]
$J:\R^{n}\to\R$ is \emph{convex} if for all $\bm{z}_{1},\bm{z}_{2}$ and
$\theta\in[0,1]$:
\begin{equation}
J\bigl(\theta\bm{z}_{1}+(1{-}\theta)\bm{z}_{2}\bigr)
\le\theta\,J(\bm{z}_{1})+(1{-}\theta)\,J(\bm{z}_{2}).
\end{equation}
Equivalently, $\nabla^{2}J\succeq 0$ everywhere.
\end{definitionbox}

In words: \emph{the chord lies above the graph}. Two further
characterisations are worth knowing because each powers a different
proof technique:
\begin{itemize}
\item \textbf{First-order:} a differentiable $J$ is convex iff its
      tangent planes are global underestimators,
      \begin{equation}\label{eq:convex-first-order}
      J(\bm{y})\ge J(\bm{z})+\nabla J(\bm{z})^{\!\top}(\bm{y}-\bm{z})
      \quad\forall\,\bm{y},\bm{z}.
      \end{equation}
      This single inequality is the source of all the global guarantees:
      if $\nabla J(\bm{z}^{\star})=\0$, then
      \eqref{eq:convex-first-order} immediately gives
      $J(\bm{y})\ge J(\bm{z}^{\star})$ for \emph{every} $\bm{y}$---a
      stationary point of a convex function is a global minimum, with a
      one-line proof.
\item \textbf{Second-order:} $\nabla^{2}J\succeq 0$ everywhere.
      Useful for checking: $J(z)=z^{4}$ is convex ($J''=12z^{2}\ge 0$);
      our running quartic $z^{4}-4z^{2}+z+4$ is not
      ($J''=12z^{2}-8<0$ near the origin---which is exactly where its
      saddle-shaped local max lives).
\end{itemize}

\paragraph{Recognition rules (the practical skill).}
Convexity is usually established not from the definition but
compositionally. Convexity is preserved by: non-negative weighted sums;
composition with an affine map ($J(\bm{A}\bm{z}+\bm{b})$ is convex if
$J$ is); pointwise maximum ($\max_{i}J_{i}$ is convex if each $J_{i}$
is); and certain monotone compositions. So
$\|\bm{A}\bm{z}-\bm{b}\|^{2}+\lambda\|\bm{z}\|_{1}$ is convex by
inspection (squared norm of affine, plus a norm). Note that the
pointwise \emph{max} rule explains why the $\max(\cdot,0)$ appearing
throughout contact modelling preserves convex structure, while a
pointwise \emph{min}---like the mode-switching min over contact
configurations---destroys it.

The key property: \textbf{every local minimum of a convex function over a
convex set is a global minimum}.  This is why we care about convex
relaxations of the contact problem.

\begin{insightbox}[What convexity buys you, operationally]
For convex problems: (i)~any stationary point found by any descent
method is \emph{the} answer---no initialisation anxiety;
(ii)~strong duality holds, so dual formulations (Section~\ref{sec:duality})
are exact, not just bounds; (iii)~polynomial-time interior-point solvers
exist with certificates of optimality. For non-convex problems all three
fail, and we settle for local minima from good initial guesses.
Trajectory optimisation through contact is non-convex twice over:
the dynamics constraint $\x_{t+1}=\f(\x_{t},\uu_{t})$ is a non-convex
equality whenever $\f$ is non-linear, and the contact complementarity
is non-convex even for linear dynamics. The art of Part~III is choosing
\emph{which} non-convexity to relax.
\end{insightbox}

%-------------------------------------------------------------
\section{Complementarity Problems}
\label{sec:complementarity}

Complementarity conditions are central to contact modelling.
They formalise the idea that two non-negative quantities cannot
both be positive simultaneously.

\begin{definitionbox}[Linear Complementarity Problem (LCP)]
Given $\bm{G}\in\R^{n\times n}$ and $\bm{g}\in\R^{n}$, find $\bm{\lambda}\in\R^{n}$ such that:
\begin{equation}\label{eq:lcp}
\bm{\lambda}\ge\0,\quad
\bm{G}\bm{\lambda}+\bm{g}\ge\0,\quad
\bm{\lambda}^{\!\top}(\bm{G}\bm{\lambda}+\bm{g})=0.
\end{equation}
\end{definitionbox}

Condition~\eqref{eq:lcp} is written in shorthand as
$\0\le\bm{\lambda}\perp(\bm{G}\bm{\lambda}+\bm{g})\ge\0$.
The ``$\perp$'' symbol denotes complementarity: for each component $i$,
either $\lambda_{i}=0$ or $(G\lambda+g)_{i}=0$ (or both).

When $\bm{G}$ is positive definite, the LCP has a unique solution and
is equivalent to a QP.  When $\bm{G}$ is only positive semidefinite
(which occurs for over-determined contact), multiple solutions may exist.

\paragraph{How LCPs are solved: case enumeration.}
The defining feature of an LCP is its hidden combinatorial structure:
each component $i$ is in one of two ``modes''
($\lambda_{i}=0$, constraint inactive; or $(\bm{G}\bm{\lambda}+\bm{g})_{i}=0$,
constraint active). With $n$ components there are $2^{n}$ mode
combinations; \emph{guessing a mode combination reduces the LCP to a
linear system}. Pivoting algorithms (Lemke, Dantzig) are clever ways of
searching mode combinations without trying all $2^{n}$.

\begin{examplebox}[Worked Example 1.4a: A 2D LCP by enumeration]
Solve $\0\le\bm{\lambda}\perp\bm{G}\bm{\lambda}+\bm{g}\ge\0$ with
\begin{equation*}
\bm{G}=\begin{bmatrix}2&1\\1&2\end{bmatrix},\qquad
\bm{g}=\begin{bmatrix}-1\\-3\end{bmatrix}.
\end{equation*}
(Think of two contact points whose forces couple through the
off-diagonal $1$'s---pressing down at one contact reduces the load
needed at the other.) Enumerate the four mode combinations:

\textbf{Mode (inactive, inactive):} $\lambda_{1}=\lambda_{2}=0$.
Need $\bm{g}\ge\0$: but $\bm{g}=(-1,-3)<\0$. \texttimes

\textbf{Mode (active, inactive):} $\lambda_{2}=0$, solve
$(G\lambda+g)_{1}=2\lambda_{1}-1=0\Rightarrow\lambda_{1}=1/2$.
Check the inactive row: $(G\lambda+g)_{2}=1/2-3=-2.5<0$. \texttimes

\textbf{Mode (inactive, active):} $\lambda_{1}=0$, solve
$2\lambda_{2}-3=0\Rightarrow\lambda_{2}=3/2$.
Check row 1: $(G\lambda+g)_{1}=3/2-1=1/2\ge 0$. \checkmark\
And $\lambda_{2}=3/2\ge 0$. \checkmark

\textbf{Mode (active, active):} solve
$\bm{G}\bm{\lambda}=-\bm{g}$:
$\lambda_{1}=-1/3<0$. \texttimes\ (Would require an adhesive force.)

\textbf{Solution:} $\bm{\lambda}^{\star}=(0,\,3/2)$. The second contact
carries all the load; the first separates with a residual gap velocity
of $1/2$. Notice that three of four modes had to be \emph{tried and
rejected}---this is the combinatorics that contact solvers spend their
time on, and that grows as $2^{n}$ (or $3^{n_c}$ with friction modes).
\end{examplebox}

\begin{definitionbox}[Non-linear Complementarity Problem (NCP)]
The NCP generalises the LCP by replacing the linear mapping with a
non-linear one.  The frictional contact problem in its full generality
is an NCP, because Coulomb friction involves a non-linear (conic)
constraint.
\end{definitionbox}

\begin{examplebox}[Worked Example 1.4: 1D LCP --- A ball on a floor]
A ball of mass $m=1$\,kg at height $\phi$ above a floor, with gravity $g=10$\,m/s$^{2}$ and time step $\Delta t=0.01$\,s.
Free velocity (gravity only): $v^{f}=v^{t}+(-g)\Delta t=v^{t}-0.1$.

The Signorini condition at the velocity level:
$0\le\lambda_{N}\perp(v^{+})\ge 0$,
where $v^{+}=v^{f}+\lambda_{N}/m$.
This is the simplest possible LCP, with $G=1/m=1$ and $g=v^{f}$.

\textbf{Case 1: Ball in air} ($\phi>0$). No contact: $\lambda_{N}=0$, $v^{+}=v^{f}$.

\textbf{Case 2: Ball on floor} ($\phi=0$).
\begin{itemize}
\item If $v^{f}\ge 0$ (moving up): $\lambda_{N}=0$, $v^{+}=v^{f}\ge 0$.
      Ball lifts off. \checkmark
\item If $v^{f}<0$ (falling into floor): we need $v^{+}\ge 0$, so
      $\lambda_{N}=-m\,v^{f}>0$, giving $v^{+}=0$ (perfectly inelastic).
      The floor pushes up just enough to prevent penetration. \checkmark
\end{itemize}

Note that both cases satisfy complementarity: when $\lambda_{N}>0$ we have $v^{+}=0$, and when $v^{+}>0$ we have $\lambda_{N}=0$.
\end{examplebox}

\begin{center}
\begin{tikzpicture}[scale=0.9]
% Hard Signorini graph
\begin{scope}
\draw[-{Stealth}] (-0.5,0) -- (3,0) node[right] {$\phi$ (gap)};
\draw[-{Stealth}] (0,-0.5) -- (0,3) node[above] {$\lambda_N$ (force)};
\draw[blue, ultra thick] (0,0) -- (0,2.5);
\draw[blue, ultra thick] (0,0) -- (2.5,0);
\fill[blue] (0,0) circle (3pt);
\node[above right] at (1.5,2) {\footnotesize Hard Signorini};
\node[blue, right] at (0.1,1.5) {\footnotesize contact};
\node[blue, above] at (1.5,0.1) {\footnotesize free};
\end{scope}
% Relaxed (spring) version
\begin{scope}[xshift=7cm]
\draw[-{Stealth}] (-1.5,0) -- (3,0) node[right] {$\phi$};
\draw[-{Stealth}] (0,-0.5) -- (0,3) node[above] {$\lambda_N$};
\draw[red, ultra thick, domain=-1.2:0, samples=50] plot (\x, {-2.5*\x});
\draw[red, ultra thick] (0,0) -- (2.5,0);
\node[above right] at (0.5,2) {\footnotesize Compliant (spring)};
\node[red, left] at (-0.8,2) {\footnotesize $\lambda_N=-k\phi$};
\end{scope}
\end{tikzpicture}
\end{center}

The left graph is the true Signorini condition---\emph{set-valued} at $\phi=0$ (infinitely steep, not a function). The right graph shows the compliant relaxation used by MuJoCo: a virtual spring replaces the hard complementarity with a smooth function $\lambda_{N}=\max(0,-k\phi)$. This smoothing is what enables gradient-based methods to work through contact, at the cost of physical fidelity.


%=============================================================================
\chapter{Linear Algebra Essentials}
\label{ch:linalg}

%-------------------------------------------------------------
\section{Vectors, Matrices, and Norms}

We work in $\R^{n}$.  Column vectors are bold lowercase ($\x$),
matrices bold uppercase ($\M$).

\textbf{Inner product:}
$\langle\x,\bm{y}\rangle=\x^{\!\top}\bm{y}=\sum_{i}x_{i}y_{i}$.

\textbf{Norms.}
The $p$-norm of $\x\in\R^{n}$:
\begin{equation}
\|\x\|_{p}=\Bigl(\sum_{i=1}^{n}|x_{i}|^{p}\Bigr)^{1/p}.
\end{equation}
Common cases: $p=2$ (Euclidean), $p=1$ (Manhattan), $p=\infty$ (max).

\textbf{Weighted norm:}
$\|\x\|_{\bm{W}}=\sqrt{\x^{\!\top}\bm{W}\x}$ for
$\bm{W}\succ 0$.
This appears throughout trajectory optimisation in cost functions.
The notation deserves a moment: $\|\bm{v}-\bm{v}^{f}\|_{\M}^{2}$ in
MuJoCo's contact problem \eqref{eq:mujoco-primal} measures velocity
change \emph{in the metric of the inertia matrix}---i.e.\ it measures
\emph{kinetic energy} of the velocity change, which is why that
formulation is called a ``least-action'' principle.

%-------------------------------------------------------------
\section{Quadratic Forms, Eigenvalues, and Conditioning}
\label{sec:quadratic-forms}

The function $J(\bm{z})=\tfrac12\bm{z}^{\!\top}\bm{A}\bm{z}$ (with
$\bm{A}$ symmetric) is the model problem of all of optimisation:
every smooth function looks like a quadratic form near a minimum.
Understanding its geometry is therefore understanding the local
geometry of \emph{every} optimisation problem.

The \emph{spectral theorem} says a symmetric $\bm{A}$ has an
orthonormal basis of eigenvectors:
$\bm{A}=\bm{V}\bm{\Lambda}\bm{V}^{\!\top}$ with
$\bm{\Lambda}=\diag(\lambda_{1},\ldots,\lambda_{n})$ and
$\bm{V}^{\!\top}\bm{V}=\I$. In eigenvector coordinates
$\bm{y}=\bm{V}^{\!\top}\bm{z}$ the quadratic decouples completely:
\begin{equation}
J=\tfrac12\sum_{i}\lambda_{i}\,y_{i}^{2}.
\end{equation}
Each eigenvalue is the \emph{curvature} along its eigenvector.
The level sets $J=\text{const}$ are ellipsoids whose axis lengths are
proportional to $1/\sqrt{\lambda_{i}}$: large eigenvalue $\Rightarrow$
steep, short axis; small eigenvalue $\Rightarrow$ shallow, long axis.

\begin{definitionbox}[Condition number]
For $\bm{A}\succ 0$, the condition number is
$\kappa(\bm{A})=\lambda_{\max}/\lambda_{\min}$.
$\kappa=1$ means spherical level sets (gradient descent converges in
one step); $\kappa\gg 1$ means a long narrow valley
(gradient descent zig-zags, see the insight box in
Section~\ref{sec:unconstrained}); $\kappa=\infty$ means a singular
matrix---a flat direction along which the minimiser is non-unique.
\end{definitionbox}

Conditioning is not an abstraction in this tutorial; it is the
recurring villain:
\begin{itemize}
\item The trajectory-optimisation Hessian is ill-conditioned because
      early controls affect everything and late controls almost nothing.
\item The Delassus matrix $\bm{G}=\J\M^{-1}\J^{\!\top}$ of a
      \emph{hyperstatic} contact configuration (more contact constraints
      than degrees of freedom, e.g.\ a four-legged table) is singular:
      $\kappa=\infty$, multiple force distributions produce identical
      motion. MuJoCo's compliance $\bm{R}$, ADMM's proximal $\rho\I$, and
      iLQR's Levenberg--Marquardt $\mu\I$ are all the same medicine:
      add a small multiple of the identity, shifting every eigenvalue
      up by the same amount, bounding $\kappa$ at the price of a
      slightly perturbed answer.
\end{itemize}

%-------------------------------------------------------------
\section{Positive Definiteness and the Cholesky Decomposition}

A symmetric matrix $\bm{A}\in\R^{n\times n}$ is \emph{positive definite}
($\bm{A}\succ 0$) if $\x^{\!\top}\bm{A}\x>0$ for all $\x\ne\0$.
Equivalently, all eigenvalues are positive.

The \emph{Cholesky decomposition} factors $\bm{A}=\bm{L}\bm{L}^{\!\top}$
where $\bm{L}$ is lower-triangular with positive diagonal entries.
This is numerically stable and roughly twice as fast as LU decomposition.
It is used extensively in contact solvers (the ADMM algorithm in
Chapter~\ref{ch:contact} requires solving
$(\bm{G}+\rho\bm{I})\bm{\lambda}=\bm{b}$,
which is done via Cholesky).

\paragraph{Why factor instead of invert.}
A rule that separates production code from homework code:
\emph{never form a matrix inverse to solve a linear system}.
To solve $\bm{A}\bm{z}=\bm{b}$, factor $\bm{A}=\bm{L}\bm{L}^{\!\top}$
once ($\tfrac{n^{3}}{3}$ flops), then solve two triangular systems
by substitution ($n^{2}$ flops each):
$\bm{L}\bm{w}=\bm{b}$ (forward), $\bm{L}^{\!\top}\bm{z}=\bm{w}$ (back).
This is faster, more accurate, and---key for ADMM, where the same
$\bm{G}+\rho\I$ is reused every iteration---the expensive factorisation
is paid once while each subsequent solve costs only $O(n^{2})$.
When you read $\bm{k}_{t}=-Q_{\uu\uu}^{-1}Q_{\uu}$ in the iLQR
equations, the implementation is a Cholesky factor-and-substitute,
never an explicit inverse. As a bonus, Cholesky doubles as the standard
\emph{test} for positive definiteness: the factorisation succeeds iff
$\bm{A}\succ 0$, and iLQR implementations use exactly this failure
signal to trigger regularisation of $Q_{\uu\uu}$.

\begin{examplebox}[Worked Example: A $2\times 2$ Cholesky solve]
Solve $\bm{A}\bm{z}=\bm{b}$ with
$\bm{A}=\begin{bmatrix}4&2\\2&3\end{bmatrix}$,
$\bm{b}=\begin{bmatrix}2\\1\end{bmatrix}$.

\textbf{Factor.} Match entries of $\bm{L}\bm{L}^{\!\top}$:
$L_{11}=\sqrt{4}=2$;\;
$L_{21}=2/L_{11}=1$;\;
$L_{22}=\sqrt{3-L_{21}^{2}}=\sqrt{2}$. So
$\bm{L}=\begin{bmatrix}2&0\\1&\sqrt{2}\end{bmatrix}$.

\textbf{Forward solve} $\bm{L}\bm{w}=\bm{b}$:
$2w_{1}=2\Rightarrow w_{1}=1$;\;
$w_{1}+\sqrt{2}\,w_{2}=1\Rightarrow w_{2}=0$.

\textbf{Back solve} $\bm{L}^{\!\top}\bm{z}=\bm{w}$:
$\sqrt{2}\,z_{2}=0\Rightarrow z_{2}=0$;\;
$2z_{1}+z_{2}=1\Rightarrow z_{1}=1/2$.

Check: $\bm{A}[1/2,\,0]^{\!\top}=[2,\,1]^{\!\top}$. \checkmark
\end{examplebox}

%-------------------------------------------------------------
\section{Least Squares and the Normal Equations}
\label{sec:least-squares}

Given an overdetermined system $\bm{A}\bm{z}\approx\bm{b}$
($\bm{A}\in\R^{m\times n}$, $m>n$), the least-squares solution
minimises $J(\bm{z})=\tfrac12\|\bm{A}\bm{z}-\bm{b}\|^{2}$.
Setting the gradient to zero,
$\nabla J=\bm{A}^{\!\top}(\bm{A}\bm{z}-\bm{b})=\0$, gives the
\emph{normal equations}:
\begin{equation}\label{eq:normal-eqs}
\bm{A}^{\!\top}\bm{A}\,\bm{z}=\bm{A}^{\!\top}\bm{b}
\quad\Longrightarrow\quad
\bm{z}^{\star}=(\bm{A}^{\!\top}\bm{A})^{-1}\bm{A}^{\!\top}\bm{b}
=\bm{A}^{+}\bm{b},
\end{equation}
where $\bm{A}^{+}$ is the (Moore--Penrose) \emph{pseudo-inverse}.
Geometrically, $\bm{A}\bm{z}^{\star}$ is the orthogonal projection of
$\bm{b}$ onto the column space of $\bm{A}$, and the residual
$\bm{b}-\bm{A}\bm{z}^{\star}$ is perpendicular to it.
Three places this exact computation appears in robotics:
solving differential inverse kinematics $\J\dot{\q}=\dot{\x}_{\text{des}}$;
fitting dynamics models to data; and the Gauss--Newton step of
Section~\ref{sec:unconstrained}, which is nothing but the normal
equations for the \emph{linearised} residual:
$\J_{r}^{\!\top}\J_{r}\,\bm{\delta}=-\J_{r}^{\!\top}\rr$.
If $\bm{A}^{\!\top}\bm{A}$ is ill-conditioned, one adds Tikhonov
regularisation $(\bm{A}^{\!\top}\bm{A}+\mu\I)$---the same
identity-shift medicine as everywhere else, here also interpretable as
a penalty $\mu\|\bm{z}\|^{2}$ on the solution norm (in IK: damped
least squares, which trades tracking accuracy for bounded joint
velocities near singularities).

%-------------------------------------------------------------
\section{Block Matrices and the Schur Complement}
\label{sec:schur}

Partition a symmetric matrix and consider minimising a quadratic over
only the second block of variables:
\begin{equation}
\min_{\bm{u}}\;
\frac12
\begin{bmatrix}\bm{x}\\\bm{u}\end{bmatrix}^{\!\top}
\begin{bmatrix}\bm{A}&\bm{B}\\\bm{B}^{\!\top}&\bm{C}\end{bmatrix}
\begin{bmatrix}\bm{x}\\\bm{u}\end{bmatrix}
+
\begin{bmatrix}\bm{a}\\\bm{c}\end{bmatrix}^{\!\top}
\begin{bmatrix}\bm{x}\\\bm{u}\end{bmatrix}.
\end{equation}
Setting the gradient w.r.t.\ $\bm{u}$ to zero:
$\bm{u}^{\star}=-\bm{C}^{-1}(\bm{c}+\bm{B}^{\!\top}\bm{x})$.
Substituting back yields a quadratic in $\bm{x}$ alone with matrix
\begin{equation}\label{eq:schur}
\bm{A}-\bm{B}\bm{C}^{-1}\bm{B}^{\!\top}
\qquad\text{(the \emph{Schur complement} of }\bm{C}\text{)}.
\end{equation}
\emph{Minimising a quadratic over a subset of variables produces a
Schur complement in the remaining variables.}
Compare this with the iLQR value-function update
\eqref{eq:Vxx-update}:
$V_{\x\x}=Q_{\x\x}-Q_{\x\uu}Q_{\uu\uu}^{-1}Q_{\uu\x}$
---it is precisely \eqref{eq:schur} with
$(\bm{A},\bm{B},\bm{C})=(Q_{\x\x},Q_{\x\uu},Q_{\uu\uu})$.
The entire DDP backward pass is one Schur complement per time step:
``minimise out the control, leave a quadratic in the state.''
A useful fact follows for free: the Schur complement of a positive
definite matrix is positive definite, which is why the backward pass
preserves $V_{\x\x}\succ 0$ when $Q_{\uu\uu}\succ 0$.

%-------------------------------------------------------------
\section{Matrix Calculus}

For a scalar-valued function $J(\bm{z})$, the \emph{gradient} is
$\nabla J\in\R^{n}$ with $(\nabla J)_{i}=\partial J/\partial z_{i}$.
The \emph{Hessian} is $\nabla^{2}J\in\R^{n\times n}$ with
$(\nabla^{2}J)_{ij}=\partial^{2}J/\partial z_{i}\partial z_{j}$.

For a vector-valued function $\f(\x)\in\R^{m}$, the \emph{Jacobian} is:
\begin{equation}
\J=\frac{\partial\f}{\partial\x}\in\R^{m\times n},\quad
J_{ij}=\frac{\partial f_{i}}{\partial x_{j}}.
\end{equation}

\paragraph{Chain rule for Jacobians.}
If $\bm{z}=\f(\x)$ and $\bm{w}=\bm{g}(\bm{z})$, then
$\frac{\partial\bm{w}}{\partial\x}=\frac{\partial\bm{g}}{\partial\bm{z}}\frac{\partial\f}{\partial\x}$.
This is how sensitivities propagate through dynamics in DDP/iLQR.

%-------------------------------------------------------------
\section{Cones}
\label{sec:cones}

\begin{definitionbox}[Cone]
A set $\Kc\subset\R^{n}$ is a \emph{cone} if for every $\x\in\Kc$ and
$\alpha\ge 0$, we have $\alpha\x\in\Kc$.
\end{definitionbox}

\begin{definitionbox}[Second-order (Lorentz / ice-cream) cone]
\begin{equation}\label{eq:soc}
\Kc_{\mu}=\bigl\{\bm{\lambda}=(\lambda_{N},\bm{\lambda}_{T})
\in\R\times\R^{2}\;\big|\;\lambda_{N}\ge 0,\;
\|\bm{\lambda}_{T}\|_{2}\le\mu\,\lambda_{N}\bigr\}.
\end{equation}
\end{definitionbox}

This is exactly the Coulomb friction cone (Section~\ref{sec:coulomb}).
The \emph{dual cone} $\Kc^{*}$ satisfies $\langle\bm{\lambda},\bm{c}\rangle\ge 0$
for all $\bm{\lambda}\in\Kc$ and $\bm{c}\in\Kc^{*}$.
For the second-order cone, $\Kc_{\mu}^{*}=\Kc_{1/\mu}$.

\paragraph{Dual cones, slowly.}
The dual cone collects every vector that makes a non-obtuse angle
($\le 90^{\circ}$) with the \emph{entire} primal cone:
\begin{equation}
\Kc^{*}=\{\bm{c}\;|\;\bm{c}^{\!\top}\bm{\lambda}\ge 0
\;\;\forall\,\bm{\lambda}\in\Kc\}.
\end{equation}
Intuition by example:
\begin{itemize}
\item The non-negative orthant is \emph{self-dual}:
      $(\R^{n}_{+})^{*}=\R^{n}_{+}$. (Only vectors with all
      non-negative entries have non-negative inner product with every
      non-negative vector.) This self-duality is why the LCP
      $\0\le\bm{\lambda}\perp\bm{w}\ge\0$ has the same cone on both
      sides.
\item A \emph{narrow} second-order cone has a \emph{wide} dual and
      vice versa: $\Kc_{\mu}^{*}=\Kc_{1/\mu}$. A high-friction cone
      ($\mu$ large, wide) has a narrow dual; a frictionless cone
      ($\mu=0$, a ray) has a half-space as its dual.
      Check the formula at $\mu=1$: $\Kc_{1}$ is self-dual,
      and indeed the $45^{\circ}$ cone is its own dual.
\item Why the dual cone appears in contact: the contact-point velocity
      in the CCP \eqref{eq:ccp} is constrained to the \emph{dual} of
      the friction cone. Forces live in $\Kc_{\mu}$; the
      complementarity partner (a velocity-like quantity) lives in
      $\Kc_{\mu}^{*}$. This generalises the 1D pairing
      ``$\lambda_{N}\ge 0$ vs.\ $v_{N}\ge 0$,'' where both sides used
      the self-dual cone $\R_{+}$, to the frictional setting where
      the two cones differ.
\end{itemize}

\paragraph{Projection onto the second-order cone.}
Section~\ref{sec:solving-constrained} promised the formula; here it is.
Write a point as $(\lambda_{N},\bm{\lambda}_{T})$ with
$s=\|\bm{\lambda}_{T}\|$. Then
\begin{equation}\label{eq:soc-projection}
\proj_{\Kc_{\mu}}(\lambda_{N},\bm{\lambda}_{T})=
\begin{cases}
(\lambda_{N},\bm{\lambda}_{T}) & s\le\mu\lambda_{N}
\quad\text{(already inside: keep)},\\[2pt]
(0,\0) & \mu s\le-\lambda_{N}
\quad\text{(inside the polar cone: project to tip)},\\[2pt]
\dfrac{\lambda_{N}+\mu s}{1+\mu^{2}}\,
\Bigl(1,\;\mu\,\dfrac{\bm{\lambda}_{T}}{s}\Bigr)
& \text{otherwise (project radially onto the surface)}.
\end{cases}
\end{equation}
The three cases have physical readings for a tentative contact force:
(i)~force already satisfies Coulomb: leave it;
(ii)~force is so adhesive/backwards that the closest admissible force is
zero: release the contact;
(iii)~force exceeds the friction limit: scale the tangential part down
to the cone surface---the contact \emph{slides}.
The $\bm{z}$-update of the ADMM contact solver
(Algorithm~\ref{alg:admm}) evaluates exactly
\eqref{eq:soc-projection} once per contact per iteration; stick, slip,
and separation emerge from which branch fires, with no explicit mode
enumeration.

\begin{examplebox}[Visualising the friction cone in 3D]
In 3D, each contact produces forces $\bm{\lambda}=(\lambda_{N},\lambda_{T_{x}},\lambda_{T_{y}})$.
The Coulomb cone $\|\bm{\lambda}_{T}\|_{2}\le\mu\lambda_{N}$ has a circular cross-section.
The LCP linearisation replaces it with a \emph{pyramid}: $\|\bm{\lambda}_{T}\|_{\infty}\le\mu\lambda_{N}$ (4 facets).
\end{examplebox}

\begin{center}
\begin{tikzpicture}[scale=0.9]
% Cross-section view
\draw[-{Stealth}] (-2.5,0) -- (2.5,0) node[right] {$\lambda_{T_{x}}$};
\draw[-{Stealth}] (0,-2.5) -- (0,2.5) node[above] {$\lambda_{T_{y}}$};
\draw[blue, very thick] (0,0) circle (1.8);
\node[blue, above right] at (1.2,1.2) {\footnotesize Coulomb (circle)};
\draw[orange, very thick] (-1.8,0) -- (0,1.8) -- (1.8,0) -- (0,-1.8) -- cycle;
\node[orange, below left] at (-1.5,1.5) {\footnotesize LCP (diamond)};
\draw[-{Stealth}, red, thick] (0,0) -- (1.27,1.27);
\node[red, right] at (1.3,1.3) {\footnotesize friction bias};
\node[below] at (0,-3) {\footnotesize Cross-section at fixed $\lambda_{N}$};
\end{tikzpicture}
\end{center}

The red arrow shows the MDP-induced bias: with the pyramid approximation, friction forces are drawn toward the corners.  A cube pushed diagonally on a flat plane will \emph{curve} instead of travelling straight---a purely numerical artifact of the LCP model.


%=============================================================================
% PART II: TRAJECTORY OPTIMISATION
%=============================================================================
\part{Trajectory Optimisation}

\chapter{Fundamentals of Trajectory Optimisation}
\label{ch:trajopt-fund}

%-------------------------------------------------------------
\section{The Optimal Control Problem}

We consider a discrete-time\footnote{Continuous-time formulations are
generally equivalent; we use discrete time for direct connection to
implementation and to MuJoCo's time-stepping.} dynamical system
with state $\x_{t}\in\R^{n}$ and control $\uu_{t}\in\R^{m}$:
\begin{equation}\label{eq:dynamics-discrete}
\x_{t+1}=\f(\x_{t},\uu_{t}),\quad t=0,1,\ldots,T{-}1.
\end{equation}

\begin{definitionbox}[Finite-Horizon Optimal Control]
Given initial state $\x_{0}$, find the control sequence
$\uu_{0:T-1}=(\uu_{0},\ldots,\uu_{T-1})$ that minimises the total cost:
\begin{equation}\label{eq:oc-objective}
\boxed{
J(\x_{0},\uu_{0:T-1})
=\sum_{t=0}^{T-1}c(\x_{t},\uu_{t})+c_{T}(\x_{T})
}
\end{equation}
subject to the dynamics~\eqref{eq:dynamics-discrete} and possibly
path constraints $\bm{h}(\x_{t},\uu_{t})\le\0$
and control bounds $\uu_{\min}\le\uu_{t}\le\uu_{\max}$.
\end{definitionbox}

The running cost $c(\x,\uu)$ penalises deviations from desired behaviour at each time step.
The terminal cost $c_{T}(\x_{T})$ penalises the final state.

\paragraph{From open-loop plans to MPC.}
Solving \eqref{eq:oc-objective} once produces an \emph{open-loop} plan:
a control tape that is correct only if the model is perfect and nothing
disturbs the system---neither of which is ever true with contact.
\emph{Model-predictive control} (MPC) converts the optimiser into a
feedback controller by re-planning constantly:
\begin{enumerate}
\item Measure (or estimate) the current state $\x_{0}$.
\item Solve (or merely \emph{improve}) the finite-horizon problem from
      $\x_{0}$.
\item Apply only the first portion of the plan to the real system.
\item Shift the previous solution forward in time to \emph{warm-start}
      the next solve; go to 1.
\end{enumerate}
The horizon recedes ahead of the system like a headlight beam---hence
``receding-horizon control.'' Feedback arises not from explicit gains
but from the re-measurement in step 1: disturbances and model errors
are absorbed at the next replan. Everything in
Chapters~\ref{ch:ddp}--\ref{ch:mjpc} should be read with this loop in
mind; in particular, the planner's job inside MPC is much easier than
solving \eqref{eq:oc-objective} to convergence, a point the next
section makes precise.

\begin{insightbox}[MuJoCo MPC notation]
In MuJoCo MPC, the running cost takes the specific form
\begin{equation}\label{eq:mjpc-cost}
c(\x,\uu)=\rho\!\left(\sum_{i=0}^{M}w_{i}\cdot\mathsf{n}_{i}\bigl(\rr_{i}(\x,\uu)\bigr);\;R\right),
\end{equation}
where $\rr_{i}$ are \emph{residuals} (sensor-based quantities that are
``small when the task is solved''), $\mathsf{n}_{i}$ are norms, $w_{i}$
are weights, and $\rho(\cdot;R)$ is a risk-sensitive exponential
transformation.  We will unpack this in Chapter~\ref{ch:mjpc}.
\end{insightbox}

%-------------------------------------------------------------
\section{Direct vs.\ Indirect Methods}

\textbf{Indirect methods} first derive analytical necessary conditions
(the Pontryagin Maximum Principle / Euler--Lagrange equations) and
then discretise these conditions for numerical solution.

\textbf{Direct methods} discretise the trajectory first, converting
the infinite-dimensional optimal control problem into a finite-dimensional
NLP.  Direct methods are broadly divided into:
\begin{itemize}
\item \textbf{Collocation (simultaneous) methods:} Both states and controls
at collocation points become decision variables.  The dynamics are
enforced as equality constraints.
\item \textbf{Shooting methods:} Only the controls (and possibly the initial
state) are decision variables.  States are determined by forward
simulation.  This is the approach used by MuJoCo MPC.
\item \textbf{Sampling-based methods:} Controls are sampled from a distribution,
rolled out, evaluated, and the distribution is updated.  No gradient
information is needed.  This includes Predictive Sampling, MPPI,
CEM, and DIAL-MPC.
\end{itemize}

\begin{insightbox}[Why shooting for MPC?]
In \emph{predictive control} (model-predictive control / receding-horizon
control), the optimisation runs in real time as the system evolves.
The key observation from Howell et al.\ (MJPC) is:
\begin{enumerate}
\item The optimiser does not need to \emph{converge}---it only needs to
      \emph{improve}.
\item \emph{Warm-starting} from the previous plan means only small
      corrections are needed.
\item Faster computation directly improves performance, because
      more optimisation steps occur per unit of real time.
\end{enumerate}
Shooting methods are natural for this setting because they produce
dynamically feasible trajectories at every iteration (unlike
collocation, which only satisfies dynamics at convergence).
\end{insightbox}

%-------------------------------------------------------------
\section{The Bellman Equation and Dynamic Programming}
\label{sec:bellman}

Define the \emph{cost-to-go} (value function) from state $\x$ at time $t$:
\begin{equation}\label{eq:value-fn}
V_{t}(\x)=\min_{\uu_{t:T-1}}\left[\sum_{\tau=t}^{T-1}c(\x_{\tau},\uu_{\tau})+c_{T}(\x_{T})\right],
\quad V_{T}(\x)=c_{T}(\x).
\end{equation}
The \emph{Bellman equation} (principle of optimality) states:
\begin{equation}\label{eq:bellman}
\boxed{
V_{t}(\x)=\min_{\uu}\bigl[c(\x,\uu)+V_{t+1}\bigl(\f(\x,\uu)\bigr)\bigr].
}
\end{equation}

This recursive structure is the foundation of DDP and iLQR:
we approximate $V_{t}$ locally as a quadratic function and solve the
minimisation~\eqref{eq:bellman} analytically at each time step.

\paragraph{Why the Bellman equation is true.}
The principle of optimality says: \emph{if a trajectory from $\x_{0}$ to
the goal is optimal, then its tail---from any intermediate state
onward---must itself be optimal} (otherwise we could splice in a better
tail and improve the whole, a contradiction).
The Bellman equation is this principle written as algebra:
the best cost from $\x$ at time $t$ equals the best over the
\emph{first action only} of (immediate cost) $+$ (best cost from
wherever that action lands you). The minimisation over a length-$(T{-}t)$
control sequence collapses into a sequence of minimisations over a
\emph{single} control each---this collapse is the entire magic of
dynamic programming.

\paragraph{Exact DP and the curse of dimensionality.}
If the state space is small and discrete, the Bellman recursion can be
evaluated \emph{exactly}: tabulate $V_{T}=c_{T}$ on a grid, then sweep
backward, evaluating the min over $\uu$ at each grid cell. This yields
the globally optimal feedback policy---no local minima, no
initialisation. The catch is cost: a grid with $g$ points per dimension
in an $n$-dimensional state space has $g^{n}$ cells. For our 9-DOF hand
with $n\approx 36$ (positions + velocities) and a coarse $g=10$:
$10^{36}$ cells. Hopeless. Every practical method in this tutorial is a
way of \emph{cheating} the curse:
\begin{itemize}
\item \textbf{DDP/iLQR} represents $V_{t}$ not on a grid but as a local
      quadratic around a single trajectory---exact DP on a quadratic
      approximation, valid in a tube around the nominal.
\item \textbf{Sampling MPC} skips $V$ entirely and just evaluates
      candidate trajectories.
\item \textbf{RL} fits $V$ (or $Q$) with a function approximator
      trained from rollouts.
\end{itemize}
Keeping ``local quadratic model of the value function'' in mind makes
the DDP equations of Chapter~\ref{ch:ddp} read as bookkeeping rather
than as revelation.

%-------------------------------------------------------------
\section{The Linear-Quadratic Regulator (LQR)}
\label{sec:lqr}

Before tackling non-linear problems, solve the one case where dynamic
programming is \emph{exact and finite-dimensional}: linear dynamics with
quadratic cost. This is not just a warm-up---iLQR is literally
``solve an LQR problem around the current trajectory, repeat,'' so
every formula here returns in Chapter~\ref{ch:ddp} wearing a $Q$.

\begin{definitionbox}[Discrete-time finite-horizon LQR]
\begin{equation}
\min_{\uu_{0:T-1}}\;
\sum_{t=0}^{T-1}\Bigl(\tfrac12\x_{t}^{\!\top}\bm{Q}_{c}\x_{t}
+\tfrac12\uu_{t}^{\!\top}\bm{R}_{c}\uu_{t}\Bigr)
+\tfrac12\x_{T}^{\!\top}\bm{Q}_{f}\x_{T}
\quad\text{s.t.}\quad
\x_{t+1}=\bm{A}\x_{t}+\bm{B}\uu_{t},
\end{equation}
with $\bm{Q}_{c},\bm{Q}_{f}\succeq 0$ and $\bm{R}_{c}\succ 0$.
\end{definitionbox}

\textbf{Claim:} the value function is exactly quadratic at every time,
$V_{t}(\x)=\tfrac12\x^{\!\top}\bm{P}_{t}\x$.
\textbf{Proof by backward induction}, which doubles as the algorithm.
The base case is $\bm{P}_{T}=\bm{Q}_{f}$. For the inductive step,
substitute into the Bellman equation:
\begin{equation*}
V_{t}(\x)=\min_{\uu}\Bigl[
\tfrac12\x^{\!\top}\bm{Q}_{c}\x+\tfrac12\uu^{\!\top}\bm{R}_{c}\uu
+\tfrac12(\bm{A}\x+\bm{B}\uu)^{\!\top}\bm{P}_{t+1}(\bm{A}\x+\bm{B}\uu)
\Bigr].
\end{equation*}
The bracket is jointly quadratic in $(\x,\uu)$. Setting its
$\uu$-gradient to zero:
\begin{equation*}
\bm{R}_{c}\uu+\bm{B}^{\!\top}\bm{P}_{t+1}(\bm{A}\x+\bm{B}\uu)=\0
\;\Longrightarrow\;
\uu^{\star}=-\underbrace{(\bm{R}_{c}+\bm{B}^{\!\top}\bm{P}_{t+1}\bm{B})^{-1}
\bm{B}^{\!\top}\bm{P}_{t+1}\bm{A}}_{\displaystyle\bm{K}_{t}}\;\x.
\end{equation*}
The optimal control is a \emph{linear feedback law}
$\uu^{\star}=-\bm{K}_{t}\x$---not a precomputed sequence but a policy
that reacts to wherever the state actually is.
Substituting $\uu^{\star}$ back (a Schur complement,
Section~\ref{sec:schur}) gives the \emph{Riccati recursion}:
\begin{equation}\label{eq:riccati}
\boxed{
\bm{P}_{t}=\bm{Q}_{c}+\bm{A}^{\!\top}\bm{P}_{t+1}\bm{A}
-\bm{A}^{\!\top}\bm{P}_{t+1}\bm{B}
(\bm{R}_{c}+\bm{B}^{\!\top}\bm{P}_{t+1}\bm{B})^{-1}
\bm{B}^{\!\top}\bm{P}_{t+1}\bm{A}.
}
\end{equation}
The quadratic form survives the induction, so the claim holds:
LQR is solved \emph{exactly} by one backward sweep
(compute $\bm{P}_{t},\bm{K}_{t}$ for $t=T{-}1,\ldots,0$) and one forward
rollout. Compare with Algorithm~\ref{alg:ilqr} and you will recognise
the skeleton: iLQR's backward/forward structure \emph{is} LQR's, applied
to the Taylor-expanded problem, with the additions of a feedforward term
$\bm{k}_{t}$ (because the expansion point is not the origin) and a line
search (because the expansion is only locally valid).

\begin{examplebox}[Worked Example: scalar LQR by hand]
Take $x_{t+1}=x_{t}+u_{t}$ ($A=B=1$), cost
$\sum_{t}\tfrac12(x_{t}^{2}+u_{t}^{2})$, horizon $T=3$,
$Q_{c}=R_{c}=1$, $Q_{f}=1$. Riccati backward
($P_{t}=1+P'-P'^{2}/(1+P')$ with $P'=P_{t+1}$, and
$K_{t}=P'/(1+P')$):
\begin{center}
\begin{tabular}{cccc}
\toprule
$t$ & $P_{t+1}$ & $K_{t}$ & $P_{t}$\\
\midrule
2 & 1.000 & 0.500 & 1.500\\
1 & 1.500 & 0.600 & 1.600\\
0 & 1.600 & 0.615 & 1.615\\
\bottomrule
\end{tabular}
\end{center}
From $x_{0}=1$: $u_{0}=-0.615$, $x_{1}=0.385$;
$u_{1}=-0.231$, $x_{2}=0.154$; $u_{2}=-0.077$, $x_{3}=0.077$.
Total cost $=\tfrac12 P_{0}x_{0}^{2}=0.808$---the value function
predicts the achieved cost \emph{exactly}, a luxury unique to the
linear-quadratic world. Note the gains converging
($0.500\to 0.600\to 0.615$) toward the infinite-horizon value
$K_{\infty}=0.618$ (the golden-ratio conjugate---the fixed point of
this Riccati equation); for long horizons LQR controllers become
effectively stationary.
\end{examplebox}

%-------------------------------------------------------------
\section{A Running Example: The Sliding Block}
\label{sec:block-example}

\begin{examplebox}[1D Block on a Frictionless Surface]
\textbf{State:} $\x=[x,\,v]^{\!\top}$ (position and velocity).\\
\textbf{Control:} $u$ (force, scalar).\\
\textbf{Dynamics} (unit mass, Euler integration with step $h$):
\begin{equation}\label{eq:block-dynamics}
\x_{t+1}=\begin{bmatrix}x_{t}+h\,v_{t}\\v_{t}+h\,u_{t}\end{bmatrix}
=\underbrace{\begin{bmatrix}1&h\\0&1\end{bmatrix}}_{\bm{A}}\x_{t}
+\underbrace{\begin{bmatrix}0\\h\end{bmatrix}}_{\bm{B}}\,u_{t}.
\end{equation}
\textbf{Boundary conditions:} $\x_{0}=[0,0]^{\!\top}$,
$\x_{T}=[1,0]^{\!\top}$.\\
\textbf{Objective (minimum effort):}
$J=\sum_{t=0}^{T-1}u_{t}^{2}$.\\[4pt]
Because the dynamics are linear and the cost quadratic, this is a
\emph{linear-quadratic} problem with an analytic solution.
The optimal force is $u^{\star}(t)=6-12t$ (on the unit interval).
\end{examplebox}

We will use this example repeatedly as we develop each algorithm.

\begin{center}
\begin{tikzpicture}[scale=1]
\fill[gray!20] (-1,-0.3) rectangle (6,-0.1);
\draw[thick] (-1,-0.1) -- (6,-0.1);
\draw[fill=blue!30, thick] (0,0) rectangle (0.8,0.6);
\node[below] at (0.4,-0.3) {\footnotesize $t=0$};
\draw[-{Stealth}, red, very thick] (0.8,0.3) -- (1.8,0.3) node[above] {$u$};
\draw[fill=green!30, thick, dashed] (4,0) rectangle (4.8,0.6);
\node[below] at (4.4,-0.3) {\footnotesize $t=T$};
\node[above] at (0.4,0.7) {\footnotesize $x=0,\;v=0$};
\node[above] at (4.4,0.7) {\footnotesize $x=1,\;v=0$};
\draw[<->, thick] (0,-0.8) -- node[below] {1 unit} (4,-0.8);
\end{tikzpicture}
\end{center}

\begin{examplebox}[Worked Example: Numerical solution with $T=4$, $h=0.25$]
We have 4 time steps and 4 controls $u_{0},u_{1},u_{2},u_{3}$. The boundary conditions $x_{4}=1$, $v_{4}=0$ give two linear equations in the four unknowns.  Combined with the quadratic cost $J=\sum u_{t}^{2}$, this is a small QP whose solution is:
\begin{center}
\begin{tabular}{ccccccc}
\toprule
$t$ & 0 & 1 & 2 & 3 & 4\\
\midrule
$u_{t}^{\star}$ & 4.80 & 1.60 & $-1.60$ & $-4.80$ & ---\\
$x_{t}$ & 0.00 & 0.00 & 0.30 & 0.70 & 1.00\\
$v_{t}$ & 0.00 & 1.20 & 1.60 & 1.20 & 0.00\\
\bottomrule
\end{tabular}
\end{center}
The force profile is anti-symmetric: accelerate then decelerate.
The velocity is bell-shaped: peak in the middle, zero at both ends.
As $T\to\infty$ with $h\to 0$, this converges to $u^{\star}(t)=6-12t$.
\end{examplebox}

\begin{center}
\begin{tikzpicture}
\begin{axis}[
  width=10cm, height=5cm, xlabel={$t$}, ylabel={},
  legend pos=north west, grid=major, grid style={gray!30},
  xtick={0,0.25,0.5,0.75,1}
]
\addplot[blue, thick, mark=*] coordinates {(0,0)(0.25,0)(0.5,0.30)(0.75,0.70)(1,1)};
\addlegendentry{$x(t)$ position}
\addplot[red, thick, mark=square*] coordinates {(0,0)(0.25,1.2)(0.5,1.6)(0.75,1.2)(1,0)};
\addlegendentry{$v(t)$ velocity}
\addplot[green!60!black, thick, mark=triangle*] coordinates {(0,4.8)(0.25,1.6)(0.5,-1.6)(0.75,-4.8)};
\addlegendentry{$u(t)$ force}
\end{axis}
\end{tikzpicture}
\end{center}


%=============================================================================
\chapter{Differential Dynamic Programming and iLQR}
\label{ch:ddp}

Differential Dynamic Programming (DDP) and its Gauss--Newton
approximation, the iterative Linear-Quadratic Regulator (iLQR),
are the most important derivative-based shooting methods for
trajectory optimisation.  MuJoCo MPC implements iLQG---which is
effectively iLQR with a slightly different historical name.

%-------------------------------------------------------------
\section{Setup and Notation}

We have the discrete-time system~\eqref{eq:dynamics-discrete} with
cost~\eqref{eq:oc-objective}.  We denote a \emph{nominal trajectory}
$(\bar\x_{0:T},\bar\uu_{0:T-1})$ satisfying the dynamics.
We seek perturbations $\delta\x_{t}=\x_{t}-\bar\x_{t}$,
$\delta\uu_{t}=\uu_{t}-\bar\uu_{t}$ that reduce the total cost.

%-------------------------------------------------------------
\section{The Q-Function}
\label{sec:q-function}

Define the \emph{state-action value function} (Q-function):
\begin{equation}\label{eq:Q-function}
Q_{t}(\x,\uu)=c(\x,\uu)+V_{t+1}\bigl(\f(\x,\uu)\bigr).
\end{equation}

The Bellman equation~\eqref{eq:bellman} becomes $V_{t}(\x)=\min_{\uu}Q_{t}(\x,\uu)$.

Now expand $Q_{t}$ to second order around the nominal $(\bar\x_{t},\bar\uu_{t})$.
Writing $\delta\x=\x-\bar\x$, $\delta\uu=\uu-\bar\uu$:
\begin{equation}\label{eq:Q-expansion}
Q_{t}\approx Q_{t}^{(0)}
+\begin{bmatrix}Q_{\x}\\Q_{\uu}\end{bmatrix}^{\!\top}
\begin{bmatrix}\delta\x\\\delta\uu\end{bmatrix}
+\frac{1}{2}
\begin{bmatrix}\delta\x\\\delta\uu\end{bmatrix}^{\!\top}
\begin{bmatrix}Q_{\x\x}&Q_{\x\uu}\\Q_{\uu\x}&Q_{\uu\uu}\end{bmatrix}
\begin{bmatrix}\delta\x\\\delta\uu\end{bmatrix}.
\end{equation}

The partial derivatives are (suppressing time subscripts):
\begin{align}
Q_{\x}&=c_{\x}+\f_{\x}^{\!\top}V_{\x}',\label{eq:Qx}\\
Q_{\uu}&=c_{\uu}+\f_{\uu}^{\!\top}V_{\x}',\label{eq:Qu}\\
Q_{\x\x}&=c_{\x\x}+\f_{\x}^{\!\top}V_{\x\x}'\,\f_{\x}
+V_{\x}'\cdot\f_{\x\x},\label{eq:Qxx}\\
Q_{\uu\uu}&=c_{\uu\uu}+\f_{\uu}^{\!\top}V_{\x\x}'\,\f_{\uu}
+V_{\x}'\cdot\f_{\uu\uu},\label{eq:Quu}\\
Q_{\x\uu}&=c_{\x\uu}+\f_{\x}^{\!\top}V_{\x\x}'\,\f_{\uu}
+V_{\x}'\cdot\f_{\x\uu},\label{eq:Qxu}
\end{align}
where primes denote evaluation at $t{+}1$, and
$\f_{\x}=\partial\f/\partial\x$,
$\f_{\uu}=\partial\f/\partial\uu$ are the dynamics Jacobians.
The terms $V_{\x}'\cdot\f_{\x\x}$, etc., involve contractions of the
value-function gradient with the dynamics Hessians (second derivatives of $\f$).

\paragraph{Where these five equations come from.}
Nothing here is new machinery---it is the chain rule applied to the
composition $Q(\x,\uu)=c(\x,\uu)+V'(\f(\x,\uu))$, and it is worth
deriving once by hand. For the gradient:
\begin{equation*}
\frac{\partial Q}{\partial\x}
=\frac{\partial c}{\partial\x}
+\Bigl(\frac{\partial\f}{\partial\x}\Bigr)^{\!\top}
\frac{\partial V'}{\partial\x'}
=c_{\x}+\f_{\x}^{\!\top}V_{\x}',
\end{equation*}
which is \eqref{eq:Qx}; equation \eqref{eq:Qu} is identical with
$\uu$ in place of $\x$. Differentiating once more produces two kinds of
terms by the product rule: differentiating $V_{\x}'$ (which depends on
$\x$ through $\f$) gives the \emph{Gauss--Newton} terms
$\f^{\!\top}V_{\x\x}'\f$, while differentiating the Jacobian
$\f_{\x}^{\!\top}$ itself gives the \emph{dynamics-curvature} terms
$V_{\x}'\cdot\f_{\x\x}$.
Pattern-match against Section~\ref{sec:unconstrained}: with
$V_{\x}'$ playing the role of the residual $\rr$ and $\f$ the role of
the residual map, this is exactly the full-Hessian-vs-Gauss--Newton
split of the least-squares Hessian
$\J_{r}^{\!\top}\J_{r}+\sum_{i}r_{i}\nabla^{2}r_{i}$.
That correspondence is the entire content of the DDP-vs-iLQR
distinction in Section~\ref{sec:ilqr}.

%-------------------------------------------------------------
\section{The DDP Backward Pass}
\label{sec:ddp-backward}

\textbf{Full DDP} retains all terms in \eqref{eq:Qxx}--\eqref{eq:Qxu},
including the dynamics Hessians $\f_{\x\x}$, $\f_{\uu\uu}$, $\f_{\x\uu}$.

Minimising the quadratic~\eqref{eq:Q-expansion} over $\delta\uu$:
\begin{equation}\label{eq:opt-du}
\delta\uu^{\star}=\argmin_{\delta\uu}Q_{t}
=-Q_{\uu\uu}^{-1}\bigl(Q_{\uu}+Q_{\uu\x}\,\delta\x\bigr)
=\bm{k}_{t}+\bm{K}_{t}\,\delta\x_{t},
\end{equation}
where
\begin{equation}\label{eq:gains}
\boxed{
\bm{k}_{t}=-Q_{\uu\uu}^{-1}Q_{\uu},\qquad
\bm{K}_{t}=-Q_{\uu\uu}^{-1}Q_{\uu\x}.
}
\end{equation}
Here $\bm{k}_{t}$ is the \emph{feedforward} correction and
$\bm{K}_{t}$ is the \emph{feedback gain} matrix.

\paragraph{Deriving the gains (do this once with a pencil).}
Hold $\delta\x$ fixed and minimise the quadratic
\eqref{eq:Q-expansion} over $\delta\uu$. Collect the
$\delta\uu$-dependent terms:
\begin{equation*}
Q_{\uu}^{\!\top}\delta\uu
+\delta\x^{\!\top}Q_{\x\uu}\,\delta\uu
+\tfrac12\,\delta\uu^{\!\top}Q_{\uu\uu}\,\delta\uu.
\end{equation*}
Set the $\delta\uu$-gradient to zero:
$Q_{\uu}+Q_{\uu\x}\delta\x+Q_{\uu\uu}\delta\uu=\0$, and solve.
The structure of the answer is meaningful:
the feedforward $\bm{k}_{t}=-Q_{\uu\uu}^{-1}Q_{\uu}$ is a
\emph{Newton step on the control} (gradient preconditioned by
curvature---compare \eqref{eq:newton}), telling us how to improve the
nominal control itself; the feedback
$\bm{K}_{t}=-Q_{\uu\uu}^{-1}Q_{\uu\x}$ tells us how to \emph{adjust}
that improvement when the state arrives off-nominal, because earlier
controls in the rollout will have changed.
Without the feedback term, an improvement at $t=0$ would invalidate the
corrections computed for $t=1,\ldots,T-1$; with it, downstream
corrections automatically adapt. This is what lets shooting methods
take large steps.

\paragraph{Deriving the value update.}
Substitute $\delta\uu^{\star}=\bm{k}+\bm{K}\delta\x$ back into the
quadratic model of $Q$ and collect terms by power of $\delta\x$.
The quadratic coefficient is
\begin{equation*}
V_{\x\x}=Q_{\x\x}+\bm{K}^{\!\top}Q_{\uu\uu}\bm{K}
+\bm{K}^{\!\top}Q_{\uu\x}+Q_{\x\uu}\bm{K}
=Q_{\x\x}-Q_{\x\uu}Q_{\uu\uu}^{-1}Q_{\uu\x},
\end{equation*}
where the second equality follows by inserting
$\bm{K}=-Q_{\uu\uu}^{-1}Q_{\uu\x}$ (two of the three correction terms
cancel against each other). The linear coefficient simplifies the same
way to \eqref{eq:Vx-update}. The result is precisely the
\emph{Schur complement} construction of Section~\ref{sec:schur}:
minimising a joint quadratic in $(\delta\x,\delta\uu)$ over
$\delta\uu$ leaves a quadratic in $\delta\x$ with the Schur-complement
matrix. Since the Schur complement of a positive definite block matrix
is positive definite, the backward pass propagates
$V_{\x\x}\succ 0$ as long as each $Q_{\uu\uu}\succ 0$---this is why
regularising $Q_{\uu\uu}$ (Section below) stabilises the whole
recursion and not just one step.

Substituting back into the value-function expansion:
\begin{align}
V_{\x,t}&=Q_{\x}-Q_{\x\uu}\,Q_{\uu\uu}^{-1}\,Q_{\uu},\label{eq:Vx-update}\\
V_{\x\x,t}&=Q_{\x\x}-Q_{\x\uu}\,Q_{\uu\uu}^{-1}\,Q_{\uu\x}.\label{eq:Vxx-update}
\end{align}
The backward pass sweeps from $t=T{-}1$ down to $t=0$, computing
$(\bm{k}_{t},\bm{K}_{t})$ and updating $(V_{\x},V_{\x\x})$ at each step.

%-------------------------------------------------------------
\section{From DDP to iLQR: Dropping the Hessian}
\label{sec:ilqr}

\begin{warningbox}[The key simplification]
Computing the dynamics Hessians $\f_{\x\x}$, $\f_{\uu\uu}$, $\f_{\x\uu}$
is expensive and can yield indefinite $Q_{\uu\uu}$ matrices.
The \textbf{Gauss--Newton approximation} drops these terms entirely:
\begin{align}
Q_{\x\x}&\approx c_{\x\x}+\f_{\x}^{\!\top}V_{\x\x}'\,\f_{\x},\label{eq:Qxx-ilqr}\\
Q_{\uu\uu}&\approx c_{\uu\uu}+\f_{\uu}^{\!\top}V_{\x\x}'\,\f_{\uu},\label{eq:Quu-ilqr}\\
Q_{\x\uu}&\approx c_{\x\uu}+\f_{\x}^{\!\top}V_{\x\x}'\,\f_{\uu}.\label{eq:Qxu-ilqr}
\end{align}
This is the \textbf{iLQR} (iterative Linear-Quadratic Regulator) algorithm.
\end{warningbox}

\textbf{Why this works well:}
\begin{enumerate}
\item The approximated $Q_{\uu\uu}$ is guaranteed positive semidefinite
      (since $c_{\uu\uu}\succeq 0$ and $\f_{\uu}^{\!\top}V_{\x\x}'\f_{\uu}\succeq 0$
      when $V_{\x\x}'\succeq 0$), ensuring the control update is
      well-defined.
\item Near the solution, the dropped terms are small
      (they are multiplied by $V_{\x}'$ which vanishes at optimality).
\item Convergence changes from quadratic (DDP) to superlinear (iLQR),
      but each iteration is much cheaper.
\end{enumerate}

%-------------------------------------------------------------
\section{The Forward Pass (Rollout with Line Search)}
\label{sec:forward-pass}

After the backward pass, we have $(\bm{k}_{t},\bm{K}_{t})$ for each $t$.
The forward pass applies these with a step-size parameter $\alpha\in[\alpha_{\min},1]$:
\begin{equation}\label{eq:forward-pass}
\boxed{
\uu_{t}=\bar\uu_{t}+\alpha\,\bm{k}_{t}+\bm{K}_{t}(\x_{t}-\bar\x_{t}),
\qquad
\x_{t+1}=\f(\x_{t},\uu_{t}).
}
\end{equation}
Multiple values of $\alpha$ are tried in parallel (line search), and the
trajectory with the lowest total cost becomes the new nominal.

\paragraph{Why scale only the feedforward.}
Note that $\alpha$ multiplies $\bm{k}_{t}$ but \emph{not} the feedback
term. The feedforward is the part that deliberately moves away from the
nominal (and is only trustworthy while the quadratic model is); the
feedback merely keeps the rollout near whatever trajectory the
feedforward is producing. Shrinking $\alpha$ therefore interpolates
between ``full Newton-like step'' ($\alpha=1$) and ``stay on the
nominal'' ($\alpha=0$), while keeping the rollout stabilised throughout.
For very non-linear systems, the $\alpha=1$ rollout can diverge
wildly---the feedback gains were computed for the old neighbourhood---so
small-$\alpha$ fallbacks are essential in practice.

\paragraph{Accepting a step: expected vs.\ actual improvement.}
The backward pass furnishes a prediction of how much a step of size
$\alpha$ should help. Summing the model decrease over time:
\begin{equation}\label{eq:expected-improvement}
\Delta J(\alpha)=
-\alpha\sum_{t}\bm{k}_{t}^{\!\top}Q_{\uu,t}
-\frac{\alpha^{2}}{2}\sum_{t}\bm{k}_{t}^{\!\top}Q_{\uu\uu,t}\bm{k}_{t}.
\end{equation}
A robust implementation (e.g.\ Tassa's) accepts the rollout only if
the \emph{actual} cost reduction is at least a fraction
$c_{1}\in(0,\tfrac12)$ of $\Delta J(\alpha)$---the Armijo condition
\eqref{eq:armijo} transplanted to trajectory space. The ratio
(actual)/(expected) is also the standard health signal:
ratios near 1 mean the quadratic model is trustworthy (decrease
regularisation, be aggressive); small or negative ratios mean the model
is over-reaching (increase regularisation $\mu$, shrink steps).
This model-trust feedback loop is what makes iLQR reliable on
hard problems, and it is exactly the trust-region logic of classical
nonlinear optimisation in trajectory clothing.

%-------------------------------------------------------------
\section{The Complete iLQR Algorithm}

\begin{algorithm}[H]
\caption{iLQR / iLQG}\label{alg:ilqr}
\KwIn{Initial state $\x_{0}$, nominal plan $\bar\uu_{0:T-1}$}
\KwOut{Improved plan $\uu_{0:T-1}^{\star}$}
\Repeat{convergence}{
  \textbf{Forward rollout:} Compute $\bar\x_{0:T}$ from $\x_{0}$ using $\bar\uu_{0:T-1}$\;
  \textbf{Backward pass:}\;
  $V_{\x,T}\leftarrow\nabla c_{T}(\bar\x_{T})$,\;
  $V_{\x\x,T}\leftarrow\nabla^{2}c_{T}(\bar\x_{T})$\;
  \For{$t=T{-}1$ \KwTo $0$}{
    Compute $\f_{\x},\f_{\uu}$ at $(\bar\x_{t},\bar\uu_{t})$\;
    Compute $Q_{\x},Q_{\uu},Q_{\x\x},Q_{\uu\uu},Q_{\x\uu}$
      via \eqref{eq:Qx}--\eqref{eq:Qu}, \eqref{eq:Qxx-ilqr}--\eqref{eq:Qxu-ilqr}\;
    $\bm{k}_{t}\leftarrow-Q_{\uu\uu}^{-1}Q_{\uu}$,\;
    $\bm{K}_{t}\leftarrow-Q_{\uu\uu}^{-1}Q_{\uu\x}$\;
    Update $V_{\x,t},V_{\x\x,t}$ via \eqref{eq:Vx-update}--\eqref{eq:Vxx-update}\;
  }
  \textbf{Forward pass with line search:}\;
  \For{each $\alpha\in\{\alpha_{1},\alpha_{2},\ldots\}$}{
    Roll out trajectory using \eqref{eq:forward-pass}\;
  }
  Set $\bar\uu\leftarrow$ best trajectory from line search\;
}
\end{algorithm}

%-------------------------------------------------------------
\section{Regularisation and Control Limits}

In practice, $Q_{\uu\uu}$ can become ill-conditioned.
Adding a regularisation term $\mu\I$ to $Q_{\uu\uu}$ before inversion
(\emph{Levenberg--Marquardt} style) ensures numerical stability:
\begin{equation}
\bm{k}_{t}=-(Q_{\uu\uu}+\mu\I)^{-1}Q_{\uu},\qquad
\bm{K}_{t}=-(Q_{\uu\uu}+\mu\I)^{-1}Q_{\uu\x}.
\end{equation}

Control limits $\uu_{\min}\le\uu\le\uu_{\max}$ are handled by the
\emph{control-limited} variant of DDP (Tassa et al., 2014), which
solves a box-constrained QP at each time step instead of the
unconstrained minimisation~\eqref{eq:opt-du}.

%-------------------------------------------------------------
\section{Worked Example: iLQR on the Sliding Block}
\label{sec:ilqr-block}

For the block \eqref{eq:block-dynamics} with $\bm{A}$, $\bm{B}$ constant:
\begin{equation}
\f_{\x}=\bm{A}=\begin{bmatrix}1&h\\0&1\end{bmatrix},\quad
\f_{\uu}=\bm{B}=\begin{bmatrix}0\\h\end{bmatrix}.
\end{equation}
Running cost $c=u^{2}$: $c_{\x}=\0$, $c_{\uu}=2u$, $c_{\uu\uu}=2$,
$c_{\x\x}=\0$, $c_{\x\uu}=\0$.
Terminal cost $c_{T}=w\|\x_{T}-\x_{\text{goal}}\|^{2}$:
$V_{\x,T}=2w(\bar\x_{T}-\x_{\text{goal}})$, $V_{\x\x,T}=2w\I$.

Since $\f_{\x\x}=\f_{\uu\uu}=\0$ (linear dynamics),
\textbf{iLQR and DDP are identical} for this problem.
The backward pass yields time-varying $\bm{k}_{t},\bm{K}_{t}$.
With large enough $w$, the solution converges to $u^{\star}(t)=6-12t$
in a single iteration (because the sub-problem is exactly quadratic).

\begin{stepbox}[Step-by-step: iLQR backward pass on the block ($T=3$, $h=1/3$, $w=100$)]

\textbf{Setup.}
$\bm{A}=\begin{bmatrix}1&1/3\\0&1\end{bmatrix}$,
$\bm{B}=\begin{bmatrix}0\\1/3\end{bmatrix}$,
$\x_{\text{goal}}=[1,0]^{\!\top}$.

Cost derivatives: $c_{\x}=\0$, $c_{\uu}=2u$, $c_{\uu\uu}=2$, $c_{\x\x}=\0$, $c_{\x\uu}=\0$.

\textbf{Initialisation.} $\bar{u}_{0}=\bar{u}_{1}=\bar{u}_{2}=0$.
Rollout: $\bar\x_{0}=\bar\x_{1}=\bar\x_{2}=\bar\x_{3}=[0,0]^{\!\top}$.

\textbf{Terminal:}
$V_{\x,3}=2w(\bar\x_{3}-\x_{\text{goal}})=200\bigl([0,0]-[1,0]\bigr)^{\!\top}=[-200,\;0]^{\!\top}$, \quad
$V_{\x\x,3}=2w\I=200\I$.

\textbf{Backward pass at $t=2$:}
\begin{align*}
Q_{\uu}&=c_{\uu}+\bm{B}^{\!\top}V_{\x,3}
=\underbrace{2\bar{u}_{2}}_{=0}
+\begin{bmatrix}0&\tfrac13\end{bmatrix}\!\begin{bmatrix}-200\\0\end{bmatrix}
=0.
\end{align*}
Pause on this: the gradient came out \emph{exactly zero}, even though we
are nowhere near the goal. Why? The terminal-cost gradient
$V_{\x,3}=[-200,\,0]^{\!\top}$ has a zero \emph{velocity} component
(the nominal already satisfies $v_{3}=0$), while
$\bm{B}=[0,\,h]^{\!\top}$ couples the control only to the velocity
component. The position error is enormous, but a force applied at the
\emph{last} step cannot change the final position (it only changes the
final velocity), so the model correctly reports ``no first-order
improvement available at $t=2$.'' The position error will instead flow
backward through the value function and create non-zero gradients at
$t=1$ and $t=0$, where forces \emph{do} have time to move the block.
This is the value function doing its job: routing credit to the time
steps that can actually act on it.

The curvature $Q_{\uu\uu}$, however, is non-zero:
\begin{align*}
Q_{\uu\uu}&=c_{\uu\uu}+\bm{B}^{\!\top}V_{\x\x,3}\bm{B}
=2+\begin{bmatrix}0&\tfrac13\end{bmatrix}\begin{bmatrix}200&0\\0&200\end{bmatrix}\begin{bmatrix}0\\\tfrac13\end{bmatrix}
=2+\tfrac{200}{9}=24.22.
\end{align*}
And the cross-term:
\begin{align*}
Q_{\uu\x}&=c_{\uu\x}+\bm{B}^{\!\top}V_{\x\x,3}\bm{A}
=\0+\begin{bmatrix}0&\tfrac13\end{bmatrix}\begin{bmatrix}200&0\\0&200\end{bmatrix}\begin{bmatrix}1&\tfrac13\\0&1\end{bmatrix}
=\begin{bmatrix}0&\tfrac{200}{3}\end{bmatrix}\begin{bmatrix}1&\tfrac13\\0&1\end{bmatrix}
=\begin{bmatrix}0&\tfrac{200}{3}\end{bmatrix}.
\end{align*}
(Watch the bookkeeping: $\bm{B}^{\!\top}V_{\x\x,3}=[0,\;200h]=[0,\;66.7]$
picks up \emph{one} factor of $h$, unlike the $Q_{\uu\uu}$ computation
above, where $\bm{B}$ appears on both sides and contributes $h^{2}$.)

So the gains at $t=2$ are:
\begin{align*}
\bm{k}_{2}&=-Q_{\uu\uu}^{-1}Q_{\uu}=-\frac{1}{24.22}\cdot 0=0,\\
\bm{K}_{2}&=-Q_{\uu\uu}^{-1}Q_{\uu\x}=-\frac{1}{24.22}\begin{bmatrix}0&66.7\end{bmatrix}=\begin{bmatrix}0&-2.752\end{bmatrix}.
\end{align*}

The feedforward $\bm{k}_{2}=0$ (no improvement at the nominal), but the feedback gain $\bm{K}_{2}$ is non-zero: if the block has positive velocity at $t=2$, the controller applies \emph{negative} force (braking).

Continuing backward to $t=1$ and $t=0$ builds up increasingly large feedforward terms as the value function propagates the position error backward.  After the forward pass with $\alpha=1$, the block moves substantially toward the goal.

\textbf{Key insight:} For this linear-quadratic problem, the backward pass computes the \emph{exact} optimal feedback policy. One forward pass gives the optimal trajectory. For non-linear problems, multiple iLQR iterations are needed, each re-linearising around the improved trajectory.
\end{stepbox}

\begin{center}
\begin{tikzpicture}
\begin{axis}[
  width=9cm, height=4.5cm,
  xlabel={iLQR iteration}, ylabel={Total cost $J$},
  grid=major, grid style={gray!30},
  ymode=log, title={iLQR convergence on block problem},
  xtick={0,1,2,3}
]
\addplot[blue, thick, mark=*] coordinates {(0,10000)(1,12.1)(2,12.0)(3,12.0)};
\node[right] at (axis cs:1.2,12.1) {\footnotesize converged};
\end{axis}
\end{tikzpicture}
\end{center}


%=============================================================================
\chapter{Sampling-Based Trajectory Optimisation}
\label{ch:sampling}

Sampling-based methods optimise trajectories without computing
gradients or Hessians.  They are well-suited to:
\begin{itemize}
\item Non-smooth dynamics (contacts, impacts).
\item Parallelisable rollouts (GPU-friendly).
\item Simple implementation---no need for dynamics derivatives.
\end{itemize}

%-------------------------------------------------------------
\section{The Common Framework}
\label{sec:sampling-framework}

All sampling-based shooting methods share the following structure:
\begin{enumerate}
\item Maintain a \emph{nominal} parameter vector $\bm{\theta}$
      parameterising the control trajectory (e.g.\ spline knots).
\item \textbf{Sample} $N{-}1$ perturbations
      $\bm{\theta}^{(i)}\sim p(\bm{\theta}\,|\,\bm{\theta}_{\text{nom}})$.
\item \textbf{Rollout} each sample (and the nominal) through the dynamics.
\item \textbf{Evaluate} the total cost $J^{(i)}$ for each sample.
\item \textbf{Update} $\bm{\theta}_{\text{nom}}$ based on the costs.
\end{enumerate}

The methods differ primarily in steps~2 (the sampling distribution)
and~5 (the update rule).

%-------------------------------------------------------------
\section{Predictive Sampling (MJPC Baseline)}
\label{sec:pred-sampling}

The simplest possible method.

\begin{algorithm}[H]
\caption{Predictive Sampling}\label{alg:pred-sampling}
\KwIn{State $\x_{0}$, nominal parameters $\bm{\theta}$, noise $\sigma$, $N$ rollouts}
Sample $\bm{\theta}^{(i)}=\bm{\theta}+\bm{\epsilon}^{(i)}$
  where $\bm{\epsilon}^{(0)}=\0$,
  $\bm{\epsilon}^{(i)}\sim\mathcal{N}(\0,\sigma^{2}\I)$ for $i\ge 1$\;
Roll out all $N$ trajectories from $\x_{0}$, compute costs $J^{(i)}$\;
$\bm{\theta}\leftarrow\bm{\theta}^{(i^{\star})}$ where
$i^{\star}=\argmin_{i}J^{(i)}$\;
\end{algorithm}

This is ``trivial random search''---pick the best sample.
Despite its simplicity, it works surprisingly well in the
predictive control setting because:
\begin{itemize}
\item Warm-starting means $\bm{\theta}$ is already near-optimal.
\item The optimiser only needs to \emph{stay in the basin of attraction}
      of the minimum, not converge to it precisely.
\item It runs faster per iteration than derivative-based methods,
      offsetting the lower per-iteration improvement.
\end{itemize}

%-------------------------------------------------------------
\section{Model Predictive Path Integral (MPPI)}
\label{sec:mppi}

MPPI (Williams et al., 2017) uses a \emph{soft} update instead of hard selection.

\begin{algorithm}[H]
\caption{MPPI}\label{alg:mppi}
\KwIn{State $\x_{0}$, nominal $\uu_{0:T-1}$, temperature $\lambda$}
\For{$i=1,\ldots,N$}{
  Sample $\bm{\epsilon}_{t}^{(i)}\sim\mathcal{N}(\0,\bm{\Sigma})$\;
  $\uu_{t}^{(i)}=\bar\uu_{t}+\bm{\epsilon}_{t}^{(i)}$\;
  Rollout, compute $J^{(i)}$\;
}
Compute weights:
$w^{(i)}=\dfrac{\exp\bigl(-J^{(i)}/\lambda\bigr)}%
{\sum_{j}\exp\bigl(-J^{(j)}/\lambda\bigr)}$\;
Update: $\bar\uu_{t}\leftarrow\bar\uu_{t}+\sum_{i}w^{(i)}\bm{\epsilon}_{t}^{(i)}$
  for all $t$\;
\end{algorithm}

\textbf{Key differences from Predictive Sampling:}
\begin{itemize}
\item \emph{Soft} (weighted average) update vs.\ \emph{hard} (best-of-$N$).
\item The temperature $\lambda$ controls exploration: $\lambda\to 0$ recovers
      hard selection; $\lambda\to\infty$ gives uniform weighting.
\item The information-theoretic derivation relates to minimising
      KL divergence between the sampling distribution and the
      optimal (Boltzmann) distribution over controls.
\end{itemize}

\paragraph{Where the exponential weights come from.}
The derivation is short enough to sketch, and it explains every design
choice. Treat the control sequence as a random variable and define the
``optimal'' distribution over control sequences as the Boltzmann
distribution $p^{\star}(U)\propto q(U)\,e^{-J(U)/\lambda}$, where $q$ is
the sampling distribution (Gaussian noise around the nominal):
low-cost trajectories get exponentially more probability mass, with
$\lambda$ setting how sharply.
We cannot sample from $p^{\star}$ directly, but we want our new nominal
to be the \emph{mean} of $p^{\star}$ (the distribution-level analogue of
``move toward good trajectories''). Estimate that mean by
\emph{importance sampling} with samples drawn from $q$:
\begin{equation*}
\mathbb{E}_{p^{\star}}[U]
=\mathbb{E}_{q}\!\left[U\,\frac{p^{\star}(U)}{q(U)}\right]
\approx\frac{\sum_{i}U^{(i)}\,e^{-J^{(i)}/\lambda}}
{\sum_{i}e^{-J^{(i)}/\lambda}}
=\sum_{i}w^{(i)}U^{(i)},
\end{equation*}
which is precisely the MPPI update. So the algorithm is: define a
target distribution that reweights your samples by exponentiated cost,
then move the nominal to its estimated mean. Two practical notes follow
directly. First, the weights are invariant to adding a constant to all
costs, so implementations subtract $\min_{i}J^{(i)}$ before
exponentiating---otherwise $e^{-J/\lambda}$ underflows to zero for
large costs. Second, the \emph{effective sample size}
$N_{\text{eff}}=1/\sum_{i}(w^{(i)})^{2}$ diagnoses the temperature:
$N_{\text{eff}}\approx 1$ means hard selection in disguise
($\lambda$ too small), $N_{\text{eff}}\approx N$ means the costs are
being ignored ($\lambda$ too large). A common heuristic adapts
$\lambda$ to hold $N_{\text{eff}}$ at a target fraction of $N$.

\begin{insightbox}[Predictive Sampling as MPPI with $\lambda\to 0$]
In the zero-temperature limit, only the best sample receives
non-zero weight, recovering Predictive Sampling.
Conversely, Predictive Sampling can be seen as
``MPPI with infinite temperature'' (from the MJPC paper's
perspective of CEM) or equivalently ``CEM with a non-adaptive
distribution.''
\end{insightbox}

\begin{examplebox}[Worked Example: MPPI hard vs.\ soft selection]
Suppose 5 samples produce costs $J^{(1:5)}=[10,\;3,\;7,\;2,\;8]$.

\textbf{Predictive Sampling} (hard selection): picks sample~4 ($J=2$).

\textbf{MPPI} with $\lambda=1$: weights $w^{(i)}\propto e^{-J^{(i)}/\lambda}$:
\begin{center}
\begin{tabular}{lccccc}
\toprule
Sample $i$ & 1 & 2 & 3 & 4 & 5\\
\midrule
$J^{(i)}$ & 10 & 3 & 7 & 2 & 8\\
$e^{-J/1}$ & $4.5\!\times\!10^{-5}$ & 0.050 & $9.1\!\times\!10^{-4}$ & 0.135 & $3.4\!\times\!10^{-4}$\\
$w^{(i)}$ (normalised) & 0.000 & 0.268 & 0.005 & \textbf{0.726} & 0.002\\
\bottomrule
\end{tabular}
\end{center}
MPPI mostly uses sample~4 but blends in some of sample~2.

\textbf{MPPI} with $\lambda=5$ (higher temperature, more exploration):
\begin{center}
\begin{tabular}{lccccc}
\toprule
Sample $i$ & 1 & 2 & 3 & 4 & 5\\
\midrule
$w^{(i)}$ & 0.075 & 0.304 & 0.137 & \textbf{0.372} & 0.112\\
\bottomrule
\end{tabular}
\end{center}
Much more uniform: all samples contribute. Higher temperature means more exploration but slower convergence.
\end{examplebox}

%-------------------------------------------------------------
\section{Cross-Entropy Method (CEM)}
\label{sec:cem}

CEM adapts both the \emph{mean} and \emph{covariance} of the sampling
distribution.

\begin{algorithm}[H]
\caption{Cross-Entropy Method}\label{alg:cem}
\KwIn{Initial distribution $\mathcal{N}(\bm{\mu},\bm{\Sigma})$, elite fraction $\rho$}
\Repeat{convergence}{
  Sample $\bm{\theta}^{(1)},\ldots,\bm{\theta}^{(N)}\sim\mathcal{N}(\bm{\mu},\bm{\Sigma})$\;
  Rollout each, compute $J^{(i)}$\;
  Sort by cost; let $\mathcal{E}$ = top-$\lfloor\rho N\rfloor$ samples\;
  $\bm{\mu}\leftarrow\frac{1}{|\mathcal{E}|}\sum_{i\in\mathcal{E}}\bm{\theta}^{(i)}$\;
  $\bm{\Sigma}\leftarrow\frac{1}{|\mathcal{E}|}\sum_{i\in\mathcal{E}}
    (\bm{\theta}^{(i)}-\bm{\mu})(\bm{\theta}^{(i)}-\bm{\mu})^{\!\top}$\;
}
\end{algorithm}

CEM fits a new Gaussian to the best (``elite'') samples.
This adapts the covariance, which can be advantageous when different
control dimensions have very different scales.

\begin{examplebox}[Worked Example: One CEM iteration in 1D]
Suppose we sample $N=10$ scalar parameters from
$\mathcal{N}(\mu=0,\sigma^{2}=1)$ and obtain (sorted by their rollout
cost, best first):
\begin{equation*}
\theta:\;\{0.9,\;1.1,\;0.7,\;1.4,\;|\;{-0.3},\;0.1,\;2.2,\;{-1.0},\;0.4,\;{-1.7}\},
\end{equation*}
with elite fraction $\rho=0.4$, so the elite set is the first four
samples. The update fits a Gaussian to the elites:
\begin{equation*}
\mu\leftarrow\tfrac14(0.9+1.1+0.7+1.4)=1.025,\qquad
\sigma^{2}\leftarrow\tfrac14\textstyle\sum(\theta_{i}-1.025)^{2}=0.0669
\;\;(\sigma\approx 0.26).
\end{equation*}
Two things happened at once: the mean \emph{moved} toward the
good region (0 $\to$ 1.025), and the distribution \emph{contracted}
($\sigma$: $1\to 0.26$) because the elites cluster.
This automatic contraction is CEM's signature---and its danger:
if the variance collapses before finding the right basin, the search
freezes (\emph{premature convergence}). Implementations counter this
with a variance floor $\sigma^{2}\ge\sigma_{\min}^{2}$ or by smoothing
the update, $\sigma^{2}\leftarrow(1-\beta)\sigma_{\text{old}}^{2}
+\beta\,\sigma_{\text{elite}}^{2}$.
Contrast with MPPI, whose covariance is \emph{fixed}: MPPI never
contracts (robust, but never refines), CEM contracts automatically
(refines, but can collapse). The annealed methods below get the best of
both by \emph{scheduling} the contraction explicitly.
\end{examplebox}

%-------------------------------------------------------------
\section{Annealed Sampling and DIAL-MPC}
\label{sec:dial-mpc}

Recent work (DIAL-MPC, LeCAR Lab) draws inspiration from
\emph{diffusion models} and \emph{simulated annealing} to construct
a sequence of sampling distributions that gradually concentrate
around the optimum.

The key idea: instead of a single round of sampling,
run $K$ rounds with \emph{decreasing} noise levels
$\sigma_{1}>\sigma_{2}>\cdots>\sigma_{K}$, re-centering the
distribution at the best (or weighted-average) sample after each round.

\begin{algorithm}[H]
\caption{Annealed / Diffusion-Inspired Sampling}\label{alg:annealed}
\KwIn{State $\x_{0}$, nominal $\bm{\theta}$, noise schedule $\sigma_{1:K}$}
\For{$k=1,\ldots,K$}{
  Sample $\bm{\theta}^{(i)}\sim\mathcal{N}(\bm{\theta},\sigma_{k}^{2}\I)$,
    $i=1,\ldots,N$\;
  Rollout each, compute $J^{(i)}$\;
  Update $\bm{\theta}$ (e.g.\ MPPI-style weighted average
    or best-of-$N$)\;
}
\end{algorithm}

This combines the broad exploration of large $\sigma$ in early rounds
with the fine-grained refinement of small $\sigma$ in later rounds,
often outperforming single-round methods for the same total number
of samples.

%-------------------------------------------------------------
\section{Spline Parameterisation}
\label{sec:splines}

All sampling methods in MuJoCo MPC do not sample the full
$m\times T$ control trajectory directly.  Instead, the control is
parameterised as a \emph{spline} with knot points
$(\tau_{0:P},\bm{\theta}_{0:P})$, and at each time step:
\begin{equation}
\uu_{t}=s\bigl(t;\,(\tau_{0:P},\bm{\theta}_{0:P})\bigr),
\end{equation}
where $s$ is the spline interpolant (zero-order hold, linear, or cubic).
This reduces the search space from $mT$ to $mP$ where $P\ll T$,
smooths the control trajectory, and makes the spline
parameters $\bm{\theta}$ the natural sampling target.

\paragraph{Why this matters: a concrete count.}
Take a 9-DOF hand ($m=9$) planned over a 1-second horizon at 100\,Hz
($T=100$): the raw control trajectory has $mT=900$ dimensions.
Random search in $\R^{900}$ is hopeless---almost every random
perturbation direction is useless, and the probability that Gaussian
noise happens to point along the narrow subspace of \emph{coordinated,
smooth} improvements is vanishing. With $P=5$ cubic spline knots, the
search space is $mP=45$ dimensions, a $20\times$ reduction, and---more
importantly---\emph{every} point in the reduced space is a smooth,
physically plausible control signal. The spline is acting as a strong
prior: it spends the sampling budget exclusively on candidates that a
robot could actually execute.
The trade-offs to keep in mind:
\begin{itemize}
\item \textbf{Expressiveness vs.\ sample efficiency.} Few knots
      $\Rightarrow$ cheap search but the optimal control may not be
      representable (e.g.\ a sharp contact-breaking force spike needs a
      knot nearby). MJPC exposes the knot count as a per-task tuning
      parameter for exactly this reason.
\item \textbf{Interpolation order.} Zero-order hold can represent
      bang-bang solutions but injects steps; cubic splines are smooth
      but round off sharp switches. For contact-rich tasks,
      lower-order interpolants often work better than intuition
      suggests, because contact events themselves provide the sharpness.
\item \textbf{Receding-horizon bookkeeping.} As the horizon slides
      forward in MPC, knot times slide with it and the oldest knot is
      dropped while a new one is appended (warm-starting,
      Section~\ref{sec:pred-sampling})---cheap precisely because the
      representation is low-dimensional.
\end{itemize}


%=============================================================================
\chapter{MuJoCo MPC: Architecture and Implementation}
\label{ch:mjpc}

%-------------------------------------------------------------
\section{Overview}

MuJoCo MPC (MJPC) is an open-source, interactive framework for
real-time predictive control built on MuJoCo physics.
It supports three planners:
\begin{enumerate}
\item \textbf{iLQG} (second-order, derivative-based),
\item \textbf{Gradient Descent} (first-order, derivative-based),
\item \textbf{Predictive Sampling} (zero-order, sampling-based).
\end{enumerate}

%-------------------------------------------------------------
\section{Cost Specification via Residuals and Norms}
\label{sec:mjpc-cost}

The running cost is:
\begin{equation}
l(\x,\uu)=\sum_{i=0}^{M}w_{i}\cdot\mathsf{n}_{i}\!\bigl(\rr_{i}(\x,\uu)\bigr),
\end{equation}
where each $\rr_{i}\in\R^{p_{i}}$ is a \emph{residual}---a vector-valued
quantity computed from MuJoCo \emph{sensors} that should be ``small when the
task is solved.''  The norm $\mathsf{n}_{i}:\R^{p_{i}}\to\R_{+}$ can be
quadratic, smooth-$\ell_{1}$, hyperbolic-cosine, etc.

The final running cost includes a risk-sensitive transformation:
\begin{equation}\label{eq:risk}
c(\x,\uu)=\rho\bigl(l(\x,\uu);R\bigr)=\frac{e^{R\cdot l(\x,\uu)}-1}{R},
\end{equation}
with $R=0$ (risk-neutral, identity), $R>0$ (risk-averse), $R<0$ (risk-seeking).

\paragraph{Why the norms matter.}
The choice of $\mathsf{n}_{i}$ is a behaviour-shaping decision, not a
formality. A quadratic norm has a gradient proportional to the residual:
large errors generate large corrective pressure, but near zero the
gradient vanishes---tasks converge quickly from far away and then stall
politely near the goal. A smooth-$\ell_{1}$ norm keeps a roughly
constant gradient at large error (robust to outliers; one bad residual
cannot dominate the cost) while remaining smooth at the origin so
derivative-based planners stay happy. Rule of thumb: quadratic for
regulation-style terms (velocities, postures), smooth-$\ell_{1}$ for
goal-reaching terms where the residual starts large.

\paragraph{Why the exponential, and a feel for $R$.}
The transformation \eqref{eq:risk} is the standard
\emph{risk-sensitive} (exponential-utility) construction. Two ways to
see what it does:
\begin{itemize}
\item \textbf{Convexity surgery.} $\rho(\cdot;R)$ is increasing, so it
      never changes \emph{which} state is better, only \emph{by how
      much}. For $R>0$ it is convex: high-cost moments get amplified
      disproportionately, so the planner works hardest on the
      \emph{worst} time steps of the trajectory---useful when a single
      bad instant (dropping the object) is catastrophic. For $R<0$ it is
      concave and saturating: beyond some badness, things can't get
      much worse, so the planner gambles---useful for escaping
      ``stuck'' configurations where every action briefly increases
      cost.
\item \textbf{Numbers.} Compare two trajectories of two steps each:
      trajectory A has per-step costs $(1,1)$; trajectory B has
      $(0,2)$---identical totals. With $R=0$: tie.
      With $R=1$: A scores $2(e^{1}-1)=3.44$, B scores
      $(e^{0}-1)+(e^{2}-1)=6.39$---the risk-averse planner strongly
      prefers the \emph{consistent} trajectory A.
      With $R=-1$: A scores $2(1-e^{-1})=1.26$, B scores
      $(1-e^{0})+(1-e^{-2})=0.86$---the risk-seeking planner prefers
      B's gamble of one perfect step.
      In stochastic settings this generalises: minimising the
      expectation of \eqref{eq:risk} is equivalent to minimising
      $\mathbb{E}[l]+\tfrac{R}{2}\mathrm{Var}[l]+\cdots$, i.e.\ $R$
      directly prices the \emph{variance} of the outcome.
\end{itemize}
For contact-rich tasks, mild risk aversion ($R>0$) is a common and
cheap robustness knob: contact uncertainty creates heavy-tailed cost
distributions, and penalising the tail keeps the planner away from
plans that are great on average but occasionally drop the object.

%-------------------------------------------------------------
\section{Why Sensor Derivatives Are ``Free''}
\label{sec:free-sensors}

\begin{insightbox}[The key efficiency trick in MuJoCo]
MuJoCo computes dynamics derivatives via \emph{finite differences} (FD).
For a system with state dimension $n$ and control dimension $m$,
the FD computation of $\partial\x_{t+1}/\partial\x_{t}$ and
$\partial\x_{t+1}/\partial\uu_{t}$ requires $n+m$ perturbed evaluations
of the step function.

Crucially, MuJoCo's step function simultaneously computes
the next state $\bm{y}$ \emph{and} all sensor values $\rr$:
\begin{equation}
(\bm{y},\rr)=\f(\x,\uu).
\end{equation}
The FD Jacobians
\begin{equation}
\frac{\partial\bm{y}}{\partial\x},\quad
\frac{\partial\bm{y}}{\partial\uu},\quad
\frac{\partial\rr}{\partial\x},\quad
\frac{\partial\rr}{\partial\uu}
\end{equation}
are all computed from the \emph{same} $n+m$ perturbed evaluations.
The cost of FD scales with the \emph{input} dimension ($n+m$),
not the output dimension.  Therefore, \textbf{adding more sensor
residuals $\rr$ does not increase the cost of derivative computation.}
\end{insightbox}

This means task designers can freely add rich sensor-based cost terms
(proximity sensors for obstacle avoidance, force sensors for contact
tasks, etc.) without any computational penalty in the derivative-based
planners.

The cost gradients are then computed via the chain rule:
\begin{equation}
\frac{\partial c}{\partial\x}=e^{Rl}\sum_{i}w_{i}\frac{\partial\mathsf{n}_{i}}{\partial\rr}\frac{\partial\rr_{i}}{\partial\x},\qquad
\frac{\partial c}{\partial\uu}=e^{Rl}\sum_{i}w_{i}\frac{\partial\mathsf{n}_{i}}{\partial\rr}\frac{\partial\rr_{i}}{\partial\uu},
\end{equation}
where the norm derivatives $\partial\mathsf{n}/\partial\rr$ are computed
analytically, and the sensor Jacobians $\partial\rr/\partial\x$,
$\partial\rr/\partial\uu$ come ``for free'' from the FD computation.

%-------------------------------------------------------------
\section{Gradient Descent Planner}

For the gradient descent planner, the control is parameterised by
spline knots $\bm{\theta}$.  The total cost gradient is:
\begin{equation}
\frac{\partial J}{\partial\bm{\theta}}
=\frac{\partial J}{\partial\Pi}\,\frac{\partial\Pi}{\partial\bm{\theta}},
\end{equation}
where $\Pi=\uu_{0:T}$ is the control sequence, and
$\partial J/\partial\Pi$ is computed via the Pontryagin Maximum Principle
(co-state / adjoint method):
\begin{align}
\bm{\lambda}_{t}&=\frac{\partial c}{\partial\x_{t}}+\left(\frac{\partial\f}{\partial\x_{t}}\right)^{\!\top}\bm{\lambda}_{t+1},\\
\frac{\partial J}{\partial\uu_{t}}&=\frac{\partial c}{\partial\uu_{t}}+\left(\frac{\partial\f}{\partial\uu_{t}}\right)^{\!\top}\bm{\lambda}_{t+1}.
\end{align}
The spline Jacobian $\partial\Pi/\partial\bm{\theta}$ has a simple
analytic formula (see the MJPC paper appendix).
A parallel line search over step sizes $\alpha$ completes the update.


%=============================================================================
% PART III: CONTACT DYNAMICS
%=============================================================================
\part{Contact Dynamics for Robotics}

\chapter{The Physics of Contact}
\label{ch:contact-physics}

Contact interactions are what make robotic manipulation and locomotion
possible.  They are also what make trajectory optimisation so difficult.
This chapter develops the physical principles from first principles,
building from 1D examples to the full 3D frictional contact problem.

%-------------------------------------------------------------
\section{Why Contact Is Hard}

\begin{warningbox}[The fundamental difficulty]
Contact dynamics are \emph{non-smooth}: the equations of motion change
discontinuously when contact is made or broken.  The contact force
is not a smooth function of penetration; it is better described by a
\emph{complementarity condition}: either there is a gap (no force)
or there is force (no gap), but not both simultaneously.
This non-smoothness means:
\begin{itemize}
\item Gradients are undefined at contact transitions.
\item The contact problem is naturally a Non-linear Complementarity Problem (NCP), which is non-convex.
\item Different simulators make different approximations to recover tractability, with different physical consequences.
\end{itemize}
\end{warningbox}

%-------------------------------------------------------------
\section{Rigid Body Dynamics (Free Motion)}

The equations of motion for an articulated rigid-body system
in generalised coordinates $\q\in\R^{n_{q}}$ with velocity
$\bm{v}\in\R^{n_{v}}$:
\begin{equation}\label{eq:eom-free}
\M(\q)\dot{\bm{v}}+\bm{b}(\q,\bm{v})=\bm{\tau},
\end{equation}
where $\M$ is the joint-space inertia matrix, $\bm{b}$ accounts for
Coriolis, centrifugal, and gravitational effects, and $\bm{\tau}$
are the generalised forces (joint torques plus external forces).

Discretising with a semi-implicit Euler scheme at step $\Delta t$:
\begin{equation}\label{eq:eom-discrete}
\M\bm{v}^{t+1}=\M\bm{v}^{t}+(\bm{\tau}-\bm{b})\Delta t
\;\equiv\;\M\bm{v}^{f},
\end{equation}
where $\bm{v}^{f}$ is the \emph{free velocity} (velocity in the absence of contacts).

%-------------------------------------------------------------
\section{The Signorini Condition (Non-Penetration)}
\label{sec:signorini}

\begin{definitionbox}[Signorini condition]
For a contact point with signed distance (gap) $\phi$ and normal
contact impulse $\lambda_{N}$:
\begin{equation}\label{eq:signorini}
\boxed{
0\le\lambda_{N}\perp c_{N}-c_{N}^{*}\ge 0.
}
\end{equation}
Here $c_{N}$ is the normal component of the contact-point velocity.
In words: \emph{the gap is non-negative, the normal force is non-negative,
and they cannot both be positive simultaneously.}
\end{definitionbox}

\begin{examplebox}[1D: Block on a floor]
Consider a block at height $\phi$ above a floor.
If $\phi>0$ (not touching), then $\lambda_{N}=0$ (no contact force).
If $\phi=0$ (touching), then $\lambda_{N}\ge 0$ (floor pushes up).
The floor cannot pull ($\lambda_{N}\ge 0$) and the block cannot penetrate ($\phi\ge 0$).

This is exactly an LCP in 1D:
$0\le\lambda_{N}\perp\phi\ge 0$.

Note that this does \emph{not} define $\lambda_{N}$ as a function of $\phi$---it is a \emph{set-valued} mapping (the graph is infinitely steep at $\phi=0$).
This is the fundamental source of non-smoothness.
\end{examplebox}

Substituting the constrained equations of motion
$\M\bm{v}^{t+1}=\M\bm{v}^{f}+\J^{\!\top}\bm{\lambda}$
(where $\J$ is the contact Jacobian mapping contact forces to
generalised forces) into the velocity-level Signorini condition yields
a \emph{Linear Complementarity Problem}:
\begin{equation}\label{eq:lcp-contact}
0\le\lambda_{N}\perp(\bm{G}\bm{\lambda}+\bm{g})_{N}-c_{N}^{*}\ge 0,
\end{equation}
where $\bm{G}=\J\M^{-1}\J^{\!\top}$ is the \emph{Delassus matrix}
and $\bm{g}=\J\bm{v}^{f}$ is the free contact-point velocity.

\paragraph{Why $\J^{\!\top}$: the duality of velocities and forces.}
The contact Jacobian is defined kinematically: it maps generalised
velocities to contact-point velocities, $\bm{c}=\J\bm{v}$.
Why, then, does its \emph{transpose} map contact forces to generalised
forces? Because power must be frame-independent. The power delivered by
contact forces, computed at the contact, is
$\bm{\lambda}^{\!\top}\bm{c}=\bm{\lambda}^{\!\top}\J\bm{v}$;
computed in joint space with generalised force $\bm{\tau}_{c}$, it is
$\bm{\tau}_{c}^{\!\top}\bm{v}$. Equating for all $\bm{v}$ gives
$\bm{\tau}_{c}=\J^{\!\top}\bm{\lambda}$. This \emph{virtual work}
argument is worth internalising---it is the same reason the grasp map
and its transpose appear in Chapter~\ref{ch:grasping}, and the same
primal/dual pairing that runs through Section~\ref{sec:duality}.

\begin{stepbox}[Deriving the Delassus matrix, step by step]
The Delassus matrix answers one question: \emph{if I apply a unit
impulse at one contact, how much does the velocity change at every
contact?} It converts the coupled rigid-body problem into a compact
problem living purely in contact coordinates.

\textbf{Step 1.} Start from the impulse-momentum equation over one step,
including contact impulses:
$\M\bm{v}^{+}=\M\bm{v}^{f}+\J^{\!\top}\bm{\lambda}$.

\textbf{Step 2.} Solve for the post-impact generalised velocity:
$\bm{v}^{+}=\bm{v}^{f}+\M^{-1}\J^{\!\top}\bm{\lambda}$.

\textbf{Step 3.} Map to contact coordinates by applying $\J$:
\begin{equation*}
\bm{c}^{+}=\J\bm{v}^{+}
=\underbrace{\J\bm{v}^{f}}_{\bm{g}}
+\underbrace{\J\M^{-1}\J^{\!\top}}_{\bm{G}}\bm{\lambda}.
\end{equation*}
The contact-space dynamics are \emph{affine}:
$\bm{c}^{+}=\bm{G}\bm{\lambda}+\bm{g}$. All the articulated-body
complexity has been compressed into one matrix.

\textbf{Step 4.} Read off the physics in $\bm{G}$.
Entry-wise, $G_{ij}$ is the velocity change at contact $i$ per unit
impulse at contact $j$:
\begin{itemize}
\item Diagonal entries $G_{ii}>0$ are \emph{inverse effective masses}:
      a light fingertip link has a large $G_{ii}$ (one Newton-second
      moves it a lot); a heavy torso has a small one.
\item Off-diagonals encode \emph{mechanical coupling}: pressing at one
      fingertip changes velocities at the others through the shared
      kinematic chain. For two contacts on bodies with no joint path
      between them, $G_{ij}=0$.
\item $\bm{G}\succeq 0$ always (it is a Gram-like matrix in the
      $\M^{-1}$ metric), but it is \emph{singular} whenever contact
      constraints outnumber effective DOFs (hyperstatic grasps, flat
      box on a plane with 4 contact points)---the conditioning villain
      of Section~\ref{sec:quadratic-forms}, and the reason solvers add
      compliance or proximal regularisation.
\end{itemize}
For a single point mass $m$ with one normal contact, $\J=1$ and
$\bm{G}=1/m$: the ball-on-floor LCP of Worked Example 1.4 is the
$1\times 1$ Delassus problem.
\end{stepbox}

%-------------------------------------------------------------
\section{Coulomb's Law of Friction}
\label{sec:coulomb}

\begin{definitionbox}[Coulomb friction cone]
The contact force $\bm{\lambda}^{(i)}=(\lambda_{N}^{(i)},\bm{\lambda}_{T}^{(i)})$
at contact $i$ must lie inside the friction cone:
\begin{equation}\label{eq:friction-cone}
\bm{\lambda}^{(i)}\in\Kc_{\mu^{(i)}}
=\bigl\{\bm{\lambda}\;\big|\;\lambda_{N}\ge 0,\;
\|\bm{\lambda}_{T}\|_{2}\le\mu^{(i)}\lambda_{N}\bigr\},
\end{equation}
where $\mu^{(i)}$ is the coefficient of friction.
This is a second-order cone (Section~\ref{sec:cones}).
\end{definitionbox}

\begin{definitionbox}[Maximum Dissipation Principle (Moreau)]
When sliding occurs, friction forces must maximise the dissipated power:
\begin{equation}\label{eq:mdp}
\bm{\lambda}_{T}^{(i)}=\argmin_{\|\bm{\gamma}_{T}\|\le\mu^{(i)}\lambda_{N}^{(i)}}
\bm{\gamma}_{T}^{\!\top}\bm{c}_{T}^{(i)}.
\end{equation}
For sliding ($\|\bm{c}_{T}\|>0$), this gives:
$\bm{\lambda}_{T}^{(i)}=-\mu^{(i)}\lambda_{N}^{(i)}\frac{\bm{c}_{T}^{(i)}}{\|\bm{c}_{T}^{(i)}\|}.$
\end{definitionbox}

%-------------------------------------------------------------
\section{The Full Contact Problem: NCP}
\label{sec:ncp}

Combining Signorini + Coulomb + MDP via the de~Saxc\'e bipotential:
\begin{equation}\label{eq:ncp}
\boxed{
\forall\,i:\quad
\Kc_{\mu^{(i)}}\ni\bm{\lambda}^{(i)}
\perp
\bm{c}^{(i)}+\Gamma(\bm{c}^{(i)},\mu^{(i)})
\in\Kc_{\mu^{(i)}}^{*},
}
\end{equation}
where $\Gamma(\bm{c},\mu)=[0,0,\mu\|\bm{c}_{T}\|_{2}]^{\!\top}$
is the de~Saxc\'e correction.

This is the \emph{Non-linear Complementarity Problem (NCP)}.
It is \textbf{non-convex} and \textbf{non-smooth}, making it
the hardest class of contact problem.


%=============================================================================
\chapter{Contact Solvers: Relaxations and Algorithms}
\label{ch:contact}

No simulator solves the full NCP \eqref{eq:ncp} exactly.
Every simulator makes approximations.  Understanding \emph{which}
approximations each simulator makes, and their consequences,
is essential for anyone doing contact-rich trajectory optimisation.

%-------------------------------------------------------------
\section{Summary of Contact Models}

\begin{center}
\begin{tabular}{lcccc}
\toprule
\textbf{Model} & \textbf{Signorini} & \textbf{Coulomb} & \textbf{MDP} & \textbf{Iterative solve?}\\
\midrule
LCP (pyramidal) & \checkmark & (linearised) & (linearised) & Yes\\
CCP (MuJoCo/Drake) & (relaxed) & \checkmark & \checkmark & Yes\\
RaiSim & \checkmark & \checkmark & (violated) & Yes\\
NCP (full) & \checkmark & \checkmark & \checkmark & Yes\\
ComFree (Jin) & (relaxed$^{*}$) & \checkmark & \checkmark & \textbf{No (closed-form)}\\
\bottomrule
\end{tabular}
\end{center}
{\footnotesize $^{*}$ComFree relaxes the dual-cone constraints (not directly Signorini), which has the side-effect of \emph{mitigating} the CCP's Signorini relaxation artifacts.}

%-------------------------------------------------------------
\section{LCP: Linearising the Friction Cone}
\label{sec:lcp-solver}

The simplest approach: approximate the circular Coulomb cone
$\Kc_{\mu}$ with a \emph{pyramid} (typically 4 facets):
\begin{equation}
\tilde\Kc_{\mu}=\bigl\{\bm{\lambda}\;\big|\;\lambda_{N}\ge 0,\;
\|\bm{\lambda}_{T}\|_{\infty}\le\mu\lambda_{N}\bigr\}.
\end{equation}

This reduces the NCP to an LCP, which is a well-studied problem.
Solvers include Lemke's algorithm (ODE) and Dantzig's algorithm (Bullet).

\begin{warningbox}[LCP artifacts]
Linearising the cone \textbf{breaks friction isotropy}.
Due to the MDP, friction forces are biased toward the pyramid
corners.  A cube sliding on a flat surface will curve instead of
travelling straight---a physically incorrect result.
Increasing the number of facets reduces but does not eliminate this bias.
\end{warningbox}

\paragraph{Projected Gauss--Seidel (PGS).}
The most common iterative solver for LCPs, used by Bullet, PhysX, ODE,
and DART.  PGS loops over contact points, updating each contact force
independently.  It is fast per iteration but converges slowly and can
produce ``internal forces'' that stretch or compress objects in
underdetermined (hyperstatic) scenarios.

Mechanically, PGS is coordinate descent plus projection: for contact
$i$, freeze all other forces, solve the \emph{scalar} (or per-contact
block) problem
\begin{equation}
\lambda_{i}\leftarrow
\proj_{\ge 0}\Bigl(\lambda_{i}
-\frac{1}{G_{ii}}\bigl(\bm{G}\bm{\lambda}+\bm{g}\bigr)_{i}\Bigr),
\end{equation}
and sweep over $i$ repeatedly. Each update is the exact solution of
contact $i$'s own 1D complementarity problem given everyone else's
current forces---so PGS can be read as contacts \emph{negotiating}, one
at a time, until consensus. The convergence rate degrades with coupling
(large off-diagonal $G_{ij}$) and with contact count: long kinematic
chains and grasps with many contacts need many sweeps, and PGS's
sequential structure resists GPU parallelisation
(Jacobi-style variants parallelise but converge even more slowly).
This per-contact locality is also the source of the internal-force
artifact: nothing in the sweep coordinates the global force
distribution when the solution is non-unique.

%-------------------------------------------------------------
\section{CCP: Relaxing the Signorini Condition}
\label{sec:ccp}

The \emph{Cone Complementarity Problem} drops the de~Saxc\'e correction:
\begin{equation}\label{eq:ccp}
\Kc_{\mu}\ni\bm{\lambda}\perp\bm{c}\in\Kc_{\mu}^{*}.
\end{equation}

This preserves the full circular friction cone and the MDP, but
\textbf{relaxes the Signorini condition}.  Rewriting the complementarity:
\begin{equation}
\lambda_{N}^{(i)}\perp\bigl(c_{N}^{(i)}-\mu\|\bm{c}_{T}^{(i)}\|\bigr),
\end{equation}
which allows simultaneous normal force and normal velocity when sliding.
Objects can ``interact at a distance'' with an artifact proportional to
$\Delta t\cdot\mu\cdot\|\bm{c}_{T}\|$.

The CCP is equivalent to a convex QCQP:
\begin{equation}\label{eq:ccp-qp}
\min_{\bm{\lambda}\in\Kc_{\mu}}\;\tfrac12\bm{\lambda}^{\!\top}\bm{G}\bm{\lambda}
+\bm{g}^{\!\top}\bm{\lambda}.
\end{equation}
This is what makes the CCP tractable---standard convex optimisation
machinery applies.

\begin{insightbox}[Two faces of the same problem: primal vs.\ dual CCP]
The CCP can be attacked from either end of the duality of
Section~\ref{sec:duality}, and the two ends correspond to the two
families of solvers in this chapter:
\begin{itemize}
\item \textbf{Primal (velocities).} Minimise over post-contact
      velocities:
      $\min_{\bm{v}}\tfrac12\|\bm{v}-\bm{v}^{f}\|_{\M}^{2}$
      subject to the contact-velocity cone constraints. The objective
      is the kinetic energy of the velocity \emph{change}---this is
      Moreau's least-action principle: \emph{contact alters the motion
      as little as the constraints allow}, measured in the inertia
      metric. MuJoCo and Drake live here.
\item \textbf{Dual (forces).} Form the Lagrangian of the primal,
      minimise the velocities out analytically (they enter
      quadratically, so this is a Schur-complement step,
      Section~\ref{sec:schur}), and maximise what remains over the
      multipliers $\bm{\lambda}\in\Kc_{\mu}$. What remains is exactly
      \eqref{eq:ccp-qp} (up to sign), with the Delassus matrix
      appearing automatically as
      $\bm{G}=\J\M^{-1}\J^{\!\top}$.
      The decision variables are now the \emph{contact forces}
      themselves, constrained to a cone with a cheap projection
      \eqref{eq:soc-projection}. ADMM and the complementarity-free
      model live here---when Section~\ref{sec:comfree-derivation}
      opens with ``recall the CCP dual,'' it means this object.
\end{itemize}
Strong duality holds (the problem is convex), so both routes find the
same physics; the choice is computational. Primal dimension $=n_{v}$
(DOFs); dual dimension $=3n_{c}$ (contacts). Few contacts on a big
robot $\Rightarrow$ dual is smaller; dense contact on a small object
$\Rightarrow$ primal is smaller.
\end{insightbox}

\subsection{MuJoCo's Approach}
\label{sec:mujoco-contact}

MuJoCo works on the \emph{primal} (velocity) formulation of the CCP:
\begin{equation}\label{eq:mujoco-primal}
\min_{\bm{v},\bm{y}}\;\tfrac12\|\bm{v}-\bm{v}^{f}\|_{\M}^{2}
+\tfrac12\|\bm{y}-\bm{c}^{*}\|_{\bm{R}^{-1}}^{2}
\quad\text{s.t.}\quad\J\bm{v}-\bm{y}\in\Kc_{\mu}^{*},
\end{equation}
where $\bm{R}$ is a \emph{compliance} matrix.  This can be reformulated
as the unconstrained problem:
\begin{equation}\label{eq:mujoco-unconst}
\min_{\bm{v}}\;\tfrac12\|\bm{v}-\bm{v}^{f}\|_{\M}^{2}
+\tfrac12\|\proj_{\Kc_{\mu}}^{\bm{R}}(\bm{y}(\J\bm{v}))\|_{\bm{R}}^{2},
\end{equation}
solved via Newton's method.

\begin{warningbox}[MuJoCo's compliance has no physical basis]
The compliance $\bm{R}$ is needed to make the problem strictly convex
(ensuring a unique solution even in hyperstatic cases).
MuJoCo sets $\bm{R}=\alpha\cdot\diag(\bm{G})$ where $\alpha\in[0,1]$.
This has \textbf{no physical justification}---it is a numerical
regularisation that shifts the solution.
It affects both normal \emph{and} tangential components, relaxing
both Signorini and Coulomb.
The net effect is to convert the hard-contact complementarity into
a regularised soft-contact model.
\end{warningbox}

\subsection{Deep Dive: Connecting \texttt{solref} and \texttt{solimp} to the CCP Math}
\label{sec:solref-solimp}

The user-facing MuJoCo parameters \texttt{solref} and \texttt{solimp} control
two distinct quantities in the CCP formulation~\eqref{eq:mujoco-primal}:
the \emph{reference velocity} $c^{*}$ and the \emph{compliance} $\bm{R}$.
Understanding this mapping is essential for tuning contact behaviour.

\subsubsection{Reference velocity: \texttt{solref} $=[d,\;b]$}

The reference velocity $c^{*}_{N}$ for the normal component at each contact
is computed as a spring-damper response to penetration:
\begin{equation}\label{eq:solref-math}
c^{*}_{N}=-\frac{1}{\Delta t}\Bigl(
\underbrace{k_{\text{ref}}\,\Delta t\;\phi(\q)}_{\text{spring (restoring)}}
+\underbrace{b_{\text{ref}}\;\dot\phi}_{\text{damper}}
\Bigr),
\end{equation}
where $\phi<0$ is the penetration depth (negative when penetrating),
and $\dot\phi$ is its rate of change.

The stiffness $k_{\text{ref}}$ is derived from the time constant $d_{\text{ref}}$:
\begin{equation}
k_{\text{ref}}=\frac{2}{d_{\text{ref}}^{2}\,(1+b_{\text{ref}})}.
\end{equation}

\textbf{Interpretation of default values} $[d_{\text{ref}},\;b_{\text{ref}}]=[0.02,\;1.0]$:
\begin{itemize}
\item $d_{\text{ref}}=0.02$\,s is the time constant of the virtual spring.
      This corresponds to stiffness $k=2/(0.02^{2}\cdot 2)=2500$\,N/m per unit mass.
      Smaller $d_{\text{ref}}$ $\Rightarrow$ stiffer spring $\Rightarrow$ less penetration,
      but harder numerics (stiffer ODE).
\item $b_{\text{ref}}=1.0$ is the damping ratio.
      $b=1$ is critical damping: the contact settles without bouncing.
      $b<1$ allows elastic bouncing (under-damped).
      $b>1$ over-damps (slower settling).
\end{itemize}

\begin{examplebox}[Worked Example: Effect of \texttt{solref} on a dropped ball]
A ball dropped from $1$\,m onto a flat plane. Time step $\Delta t=0.002$\,s.

\textbf{Case 1:} \texttt{solref}$\;=[0.02, 1.0]$ (default, stiff, critically damped).
Ball penetrates ${\sim}0.1$\,mm. Settles in one bounce. Looks ``hard.''

\textbf{Case 2:} \texttt{solref}$\;=[0.1, 1.0]$ (soft, critically damped).
$k=2/(0.1^{2}\cdot 2)=100$\,N/m.
Ball penetrates ${\sim}5$\,mm. Settles slowly. Looks ``squishy.''

\textbf{Case 3:} \texttt{solref}$\;=[0.02, 0.3]$ (stiff, under-damped).
Ball bounces several times before settling (like a rubber ball).
\end{examplebox}

\subsubsection{Impedance: \texttt{solimp} $=[d_{\min},\;d_{\max},\;w,\;\ldots]$}

The impedance (inverse compliance, $R^{-1}$) controls how the contact
stiffness varies with penetration depth:
\begin{equation}\label{eq:solimp-math}
R_{N}^{-1}(\phi)=d_{\min}+(d_{\max}-d_{\min})\cdot\sigma\!\left(\frac{\phi}{w}\right),
\end{equation}
where $\sigma(\cdot)$ is a smooth sigmoid (step) function and $w$ controls the
transition width.

\begin{itemize}
\item At zero penetration ($\phi\approx 0$): impedance $\approx(d_{\min}+d_{\max})/2$.
\item Deep penetration ($\phi\gg w$): $R_{N}^{-1}\to d_{\max}$ (high stiffness, hard wall).
\item No penetration ($\phi\ll 0$, separating): $R_{N}^{-1}\to d_{\min}$ (low stiffness, easy release).
\end{itemize}

Default values: $[d_{\min},d_{\max},w]=[0.9,\;0.95,\;0.001]$.
These are close to 1.0, meaning the impedance is nearly constant and close to unity.

\begin{center}
\begin{tikzpicture}
\begin{axis}[
  width=10cm, height=5.5cm,
  xlabel={Penetration depth $\phi$ (m)}, ylabel={Impedance $R_{N}^{-1}$},
  grid=major, grid style={gray!30},
  domain=-0.01:0.01, samples=100,
  title={Impedance profile from \texttt{solimp}},
  ymin=0.85, ymax=1.0,
  legend pos=south east
]
\addplot[blue, very thick]{0.9 + 0.05/(1+exp(-x/0.001))};
\addlegendentry{$d_{\min}=0.9$, $d_{\max}=0.95$, $w=0.001$}
\addplot[red, thick, dashed]{0.5 + 0.45/(1+exp(-x/0.003))};
\addlegendentry{$d_{\min}=0.5$, $d_{\max}=0.95$, $w=0.003$}
\draw[dashed, gray, thin] (axis cs:-0.01,0.9) -- (axis cs:0.01,0.9);
\draw[dashed, gray, thin] (axis cs:-0.01,0.95) -- (axis cs:0.01,0.95);
\end{axis}
\end{tikzpicture}
\end{center}

\subsubsection{How compliance enters the CCP}

For each contact point $i$, the compliance matrix is:
\begin{equation}
\bm{R}^{(i)}=\diag\bigl(R_{N}^{(i)},\;R_{T_{x}}^{(i)},\;R_{T_{y}}^{(i)}\bigr),
\end{equation}
where $R_{N}$ is determined by \texttt{solimp} and $R_{T}$ is also set to a non-zero value by MuJoCo.

This enters the Delassus matrix via:
\begin{equation}
\tilde{\bm{G}}=\bm{G}+\bm{R}=\J\M^{-1}\J^{\!\top}+\bm{R}.
\end{equation}
Adding $\bm{R}$ ensures $\tilde{\bm{G}}$ is positive definite even when $\bm{G}$
is singular (hyperstatic contacts), guaranteeing a unique solution.

\begin{warningbox}[The tangential compliance controversy]
MuJoCo also sets $R_{T}\ne 0$, which \textbf{relaxes Coulomb's law}.
With tangential compliance, friction becomes \emph{viscous} (proportional to
sliding velocity) rather than \emph{dry} (set-valued).
Consequences:
\begin{itemize}
\item Objects on slopes may slowly \emph{creep} instead of remaining stationary.
\item Friction forces become velocity-dependent, deviating from the maximum dissipation principle.
\item The sliding-cube diagnostic test (Le~Lidec et al.) shows measurably different energy dissipation.
\end{itemize}

\textbf{Le~Lidec et al.'s position:} ``The compliance added in MuJoCo has no physical purpose and should be considered a numerical trick designed to circumvent issues due to hyper-staticity or ill-conditioning, at the cost of impairing the simulation.''

\textbf{MuJoCo's counter-argument:} The compliance makes the problem well-conditioned, enables Newton's method to converge reliably, and the artifacts are small relative to other sources of modelling error for typical robotics time steps ($\Delta t\lesssim 5$\,ms).

\textbf{Our assessment:} Both positions have merit.  For \emph{trajectory optimisation}, MuJoCo's compliance provides the smooth gradients that iLQR/DDP requires, at a known and controllable cost in physical accuracy.  For \emph{high-fidelity contact analysis}, ADMM on the true CCP (or NCP) is preferable.  For \emph{contact-rich MPC}, sampling-based methods sidestep the issue entirely by not requiring gradients.
\end{warningbox}

\subsubsection{Comparison: MuJoCo vs.\ Drake vs.\ hard-contact ADMM}

\begin{center}
\begin{tabular}{lccc}
\toprule
\textbf{Property} & \textbf{MuJoCo} & \textbf{Drake} & \textbf{ADMM (hard)}\\
\midrule
Compliance $\bm{R}$ & $\alpha\diag(\bm{G})$ (numerical) & material stiffness & $\bm{R}=\0$ (proximal $\rho$)\\
Signorini & Relaxed (spring) & Relaxed (hydroelastic) & Exact\\
Coulomb & Relaxed (viscous $R_{T}$) & Relaxed (less so) & Exact\\
Hyperstatic & Always unique & Always unique & Unique via $\rho$\\
Gradients available & Yes (smooth) & Yes (smooth) & No (discontinuous)\\
Hard contacts & No & No & Yes\\
Internal forces & No & No & No\\
\bottomrule
\end{tabular}
\end{center}

\subsection{Drake's Approach}
\label{sec:drake-contact}

Drake (Castro, Permenter, Han, 2022) also solves the primal CCP
formulation but provides a more principled way of setting $\bm{R}$.
Drake's \emph{hydroelastic contact} model computes contact patches
(not point contacts) between geometries, and derives compliance from
material properties.

The key differences from MuJoCo:
\begin{enumerate}
\item Compliance $\bm{R}$ is set based on \emph{material stiffness}, not
      as a numerical hack based on $\diag(\bm{G})$.
\item Uses a Newton solver with an advanced line search.
\item The hydroelastic model computes a pressure field over the
      contact surface, giving physically meaningful ``force at a distance''
      for nearly-touching objects.
\end{enumerate}

\subsection{ADMM for the CCP}

An alternative to Newton's method is the Alternating Direction Method
of Multipliers (ADMM), which solves the \emph{dual} CCP \eqref{eq:ccp-qp}:

\begin{algorithm}[H]
\caption{ADMM for CCP}\label{alg:admm}
$\tilde{\bm{G}}^{-1}\leftarrow(\bm{G}+\rho\I)^{-1}$\;
\For{$k=1,\ldots,n_{\text{iter}}$}{
  $\bm{\lambda}\leftarrow-\tilde{\bm{G}}^{-1}(\bm{g}-\rho\bm{z}+\bm{\gamma})$\;
  $\bm{z}\leftarrow\proj_{\Kc_{\mu}}(\bm{\lambda}+\bm{\gamma}/\rho)$\;
  $\bm{\gamma}\leftarrow\bm{\gamma}+\rho(\bm{\lambda}-\bm{z})$\;
}
\end{algorithm}

ADMM has several advantages:
\begin{itemize}
\item Handles hard contacts ($\bm{R}=\0$) naturally via the proximal
      regularisation $\rho$.
\item Produces no internal forces in underdetermined problems.
\item Robust to ill-conditioned Delassus matrices.
\item Benefits from efficient Cholesky factorisation of $\bm{G}+\rho\I$.
\end{itemize}

\paragraph{Reading the three steps with Chapter 1 eyes.}
ADMM splits the variable into two copies---$\bm{\lambda}$
(unconstrained, handles the quadratic) and $\bm{z}$ (constrained,
handles the cone)---tied together by the consensus constraint
$\bm{\lambda}=\bm{z}$, and runs an augmented-Lagrangian
(Section~\ref{sec:solving-constrained}) loop on the split problem:
\begin{enumerate}
\item The $\bm{\lambda}$-step minimises the augmented Lagrangian over
      the unconstrained copy: a \emph{linear solve} with the fixed
      matrix $\bm{G}+\rho\I$---factor once with Cholesky
      (Section~\ref{ch:linalg}), back-substitute every iteration.
      Note $\rho\I$ is the identity-shift regulariser yet again:
      even a singular $\bm{G}$ (hyperstatic contact) yields a
      well-posed step, which is precisely why ADMM produces no
      internal forces where PGS does.
\item The $\bm{z}$-step minimises over the constrained copy: by
      construction this is a pure \emph{projection} onto the friction
      cone, evaluated per-contact by the closed-form
      \eqref{eq:soc-projection}.
\item The $\bm{\gamma}$-step is the multiplier update of the method of
      multipliers: the running disagreement
      $\bm{\lambda}-\bm{z}$ is absorbed into the dual variable, tilting
      the next subproblems toward consensus.
\end{enumerate}
Nothing in the loop is new---it is the augmented Lagrangian plus the
cone projection, alternated. This compositional view is also the
practical debugging view: slow convergence almost always traces to the
penalty parameter $\rho$ (too small: the projection dominates and
consensus crawls; too large: the linear solve dominates and the cone is
ignored for many iterations), and adaptive-$\rho$ schemes monitor the
primal residual $\|\bm{\lambda}-\bm{z}\|$ against the dual residual
$\rho\|\bm{z}^{k+1}-\bm{z}^{k}\|$ to rebalance.

%-------------------------------------------------------------
\section{Solving the Full NCP}

\paragraph{PGS for NCP.}
Adapts the LCP PGS by changing the projection step to use the
de~Saxc\'e correction.  No convergence guarantees.

\paragraph{Staggered projections.}
Rewrites the NCP as two interleaved convex problems (one for normal
forces, one for tangential) solved in alternation.  Typically converges
in ${\sim}5$ iterations.  Used in the differentiable simulator of
Le~Lidec et al.

%-------------------------------------------------------------
\section{Complementarity-Free Contact Modelling}
\label{sec:comfree}

All methods above---LCP, CCP, NCP---share a common structure: they formulate
contact as a \emph{complementarity problem} that must be solved (iteratively) at every
simulation time step. Jin~(2025) proposed a fundamentally different approach:
\emph{dispense with complementarity entirely} by deriving a closed-form contact
force that automatically satisfies Coulomb's friction law.

\subsection{Motivation: Why Bypass Complementarity?}

Complementarity-based solvers create three difficulties for trajectory optimisation:
\begin{enumerate}
\item \textbf{Non-smoothness.} The mapping from state to contact force is
      set-valued (or at best piecewise-smooth), breaking gradient-based planners.
\item \textbf{Combinatorial complexity.} With $n_{c}$ contacts, there are $3^{n_{c}}$
      possible contact modes (stick/slide/separate per contact).  Complementarity
      solvers must implicitly or explicitly enumerate these modes.
\item \textbf{Computational cost.} Iterative solves (PGS, ADMM, Newton) scale
      super-linearly with the number of contacts.
\end{enumerate}

The complementarity-free model addresses all three: it provides a
\emph{closed-form, differentiable, Coulomb-consistent} contact force computation
whose cost scales \emph{linearly} with the number of contacts.

\subsection{From the CCP Dual to Closed-Form Forces}
\label{sec:comfree-derivation}

Recall the CCP dual problem~\eqref{eq:ccp-qp}:
\begin{equation}\label{eq:ccp-dual-recap}
\max_{\bm{\beta}\ge\0}\;-\tfrac{1}{2h^{2}}\bm{\beta}^{\!\top}
\bigl(\tilde{\J}\bm{Q}^{-1}\tilde{\J}^{\!\top}+\bm{R}\bigr)\bm{\beta}
-\tfrac{1}{h}\bigl(\tilde{\J}\bm{Q}^{-1}\bm{b}+\tilde{\bm{\phi}}\bigr)^{\!\top}\bm{\beta}
+\text{const},
\end{equation}
where $\bm{Q}$ encodes the system inertia, $\bm{b}$ the non-contact forces,
$\tilde{\J}$ stacks the linearised dual-cone Jacobians (normal minus $\mu$ times
tangential for each polyhedral facet), $\tilde{\bm{\phi}}$ stacks the gap distances,
and $\bm{\beta}\ge\0$ are the dual variables (contact impulses in the polyhedral
dual-cone parameterisation).

The KKT optimality condition for \eqref{eq:ccp-dual-recap} is a
\emph{complementarity condition} in the dual:
\begin{equation}\label{eq:dual-comp}
0\le\bm{\beta}\perp
\frac{1}{h}\bigl(\tilde{\J}\bm{Q}^{-1}\tilde{\J}^{\!\top}+\bm{R}\bigr)\bm{\beta}
+\bigl(\tilde{\J}\bm{Q}^{-1}\bm{b}+\tilde{\bm{\phi}}\bigr)\ge\0.
\end{equation}

\begin{stepbox}[Key step: Replacing the coupled matrix with a diagonal]
The matrix $(\tilde{\J}\bm{Q}^{-1}\tilde{\J}^{\!\top}+\bm{R})^{-1}$
couples all contacts.  Jin replaces it with a positive-definite
\emph{diagonal} matrix $\bm{K}(\q)$:
\begin{equation}\label{eq:K-replacement}
\bm{K}(\q)\approx\bigl(\tilde{\J}\bm{Q}^{-1}\tilde{\J}^{\!\top}+\bm{R}\bigr)^{-1}.
\end{equation}
With this diagonal approximation, the complementarity~\eqref{eq:dual-comp}
\emph{decouples across contacts} and admits a closed-form solution
using the identity $0\le x\perp(x+\lambda)\ge 0\;\Leftrightarrow\;x=\max(-\lambda,0)$:
\begin{equation}\label{eq:beta-closed}
\boxed{
\bm{\beta}^{+}=\max\!\Bigl(-h\,\bm{K}(\q)\bigl(\tilde{\J}\bm{Q}^{-1}\bm{b}+\tilde{\bm{\phi}}\bigr),\;\0\Bigr).
}
\end{equation}
\end{stepbox}

The resulting post-contact velocity is:
\begin{equation}\label{eq:comfree-velocity}
\bm{v}^{+}=\frac{1}{h}\bm{Q}^{-1}\bm{b}
+\frac{1}{h}\bm{Q}^{-1}\tilde{\J}^{\!\top}
\max\!\Bigl(-\bm{K}(\q)\bigl(\tilde{\J}\bm{Q}^{-1}\bm{b}+\tilde{\bm{\phi}}\bigr),\;\0\Bigr).
\end{equation}

\subsection{Physical Interpretation: A Spring in the Dual Cone}
\label{sec:comfree-spring}

Equation~\eqref{eq:comfree-velocity} has a beautiful physical interpretation.
Rewriting with explicit labels:
\begin{equation}\label{eq:comfree-spring}
h\bm{v}^{+}=\bm{Q}^{-1}\!\Bigl(
\underbrace{\bm{b}}_{\substack{\text{non-contact}\\\text{forces}}}
+\;\tilde{\J}^{\!\top}\underbrace{\max\!\bigl(-\bm{K}\cdot
\overbrace{\min(\tilde{\J}\bm{Q}^{-1}\bm{b}+\tilde{\bm{\phi}},\;\0)}^{
\text{penetration into dual cone}},\;\0\bigr)}_{\bm{\lambda}_{\text{contact}}\;\text{(spring force in dual cone)}}
\Bigr).
\end{equation}

\begin{insightbox}[The dual-cone spring analogy]
\begin{enumerate}
\item First, compute the displacement $\bm{Q}^{-1}\bm{b}$ that would occur
      under non-contact forces alone (gravity, actuation).
\item Check whether this displacement \emph{violates} the dual-cone constraints
      $\tilde{\J}\bm{v}\,h+\tilde{\bm{\phi}}\ge\0$.
      The ``penetration depth'' into the dual cone is
      $\min(\tilde{\J}\bm{Q}^{-1}\bm{b}+\tilde{\bm{\phi}},\;\0)\le\0$.
\item Apply a \emph{spring-like restoring force} proportional to this penetration,
      with stiffness $\bm{K}(\q)$.  This force lives in the dual cone and is
      mapped to generalised coordinates via $\tilde{\J}^{\!\top}$.
\end{enumerate}
The $\max(\cdot,\0)$ ensures that contact forces are only generated when
the dual cone is actually violated---otherwise (separation), no force is applied.
This is illustrated in Figure~\ref{fig:comfree-spring}.
\end{insightbox}

\begin{figure}[ht]
\centering
\begin{tikzpicture}[scale=0.85]
% Panel (a): dual cone
\begin{scope}
\fill[red!10] (0,0) -- (-1.2,2.5) -- (1.2,2.5) -- cycle;
\draw[red, thick] (0,0) -- (-1.2,2.5);
\draw[red, thick] (0,0) -- (1.2,2.5);
\node[red, above] at (0,2.6) {\footnotesize dual cone};
% Predicted motion
\draw[-{Stealth}, green!50!black, very thick] (0,1.5) -- (0.8,0.3);
\node[green!50!black, right] at (0.8,0.3) {\footnotesize $\bm{Q}^{-1}\bm{b}$};
\node[below] at (0,-0.3) {\footnotesize (a) Prediction};
\end{scope}
% Panel (b): penetration
\begin{scope}[xshift=5cm]
\fill[red!10] (0,0) -- (-1.2,2.5) -- (1.2,2.5) -- cycle;
\draw[red, thick] (0,0) -- (-1.2,2.5);
\draw[red, thick] (0,0) -- (1.2,2.5);
\draw[-{Stealth}, green!50!black, very thick] (0,1.5) -- (0.8,0.3);
% Penetration arrow
\draw[{Stealth}-, thick, black] (0.55,1.07) -- (0.8,0.3);
\node[black, right] at (0.9,0.7) {\footnotesize penetration};
\node[below] at (0,-0.3) {\footnotesize (b) Violation};
\end{scope}
% Panel (c): spring correction
\begin{scope}[xshift=10cm]
\fill[red!10] (0,0) -- (-1.2,2.5) -- (1.2,2.5) -- cycle;
\draw[red, thick] (0,0) -- (-1.2,2.5);
\draw[red, thick] (0,0) -- (1.2,2.5);
\draw[-{Stealth}, green!50!black, very thick] (0,1.5) -- (0.8,0.3);
% Spring force
\draw[-{Stealth}, blue, very thick] (0.8,0.3) -- (0.3,1.3);
\node[blue, left] at (0.2,1.0) {\footnotesize $\bm{K}\cdot$pen};
% Corrected
\draw[-{Stealth}, violet, very thick, dashed] (0,1.5) -- (0.3,1.3);
\node[violet, above left] at (0.2,1.4) {\footnotesize $\bm{v}^{+}$};
\node[below] at (0,-0.3) {\footnotesize (c) Correction};
\end{scope}
\end{tikzpicture}
\caption{Physical interpretation of complementarity-free contact.
(a)~The dual cone (red) and the predicted motion under non-contact forces (green).
(b)~The predicted motion violates the dual cone; the violation depth is computed.
(c)~A spring force (blue) proportional to the violation pushes the motion back,
yielding the corrected velocity $\bm{v}^{+}$ (violet).}
\label{fig:comfree-spring}
\end{figure}

\subsection{Automatic Coulomb Satisfaction}

A key property: the contact force $\bm{\lambda}_{\text{contact}}$ in~\eqref{eq:comfree-spring}
automatically satisfies Coulomb's friction law.  This follows because
$\bm{\beta}^{+}\ge\0$ by construction (it is the output of $\max(\cdot,0)$),
and the mapping from $\bm{\beta}$ back to normal and tangential forces
via the polyhedral parameterisation guarantees
$\mu_{i}\|\bm{\lambda}_{T}^{(i)}\|\le\|\bm{\lambda}_{N}^{(i)}\|$
by the triangle inequality. No separate friction projection is needed.

\begin{examplebox}[Worked Example: ComFree on a 1D block on a floor]
Consider a unit-mass block at height $\phi$ above a floor with no tangential
motion (frictionless, 1D normal contact only).

The dual-cone constraint reduces to $v_{N}^{+}+\phi/h\ge 0$ (non-penetration
at velocity level).  With $\bm{Q}=m/h^{2}=1/h^{2}$, $\bm{b}=mv/h+\tau$
(where $\tau=-mg$), and $\tilde{J}=1$ (normal Jacobian):

$\tilde{J}Q^{-1}b+\tilde{\phi}=h^{2}(v/h-g)+\phi=hv-h^{2}g+\phi.$

For a block at rest ($v=0$) sitting on the floor ($\phi=0$) with $g=10$, $h=0.01$:

$\tilde{J}Q^{-1}b+\phi=0-0.001+0=-0.001<0.$

The dual cone is violated (gravity would push the block into the floor).
The spring force:
$\beta^{+}=\max(-K\cdot(-0.001),0)=0.001K.$

With $K\approx (JQ^{-1}J^{\!\top})^{-1}=Q=1/h^{2}=10{,}000$:
$\beta^{+}=10$, which gives a contact force $\lambda_{N}=\beta/h=1000$\,N$\cdot$s$^{-1}$...
but wait, we need to be careful with units.
In the impulse formulation: $\lambda_{N}\,h=\beta^{+}=10$, so $\lambda_{N}=10/0.01=1000$.
This is exactly $mg=1\times 10/0.01$ in the impulse-per-step units---the floor
supports the block against gravity.  \checkmark

If instead $\phi=0.01$ (block is $1$\,cm \emph{above} the floor):
$\tilde{J}Q^{-1}b+\phi=-0.001+0.01=0.009>0$.
No violation, so $\beta^{+}=\max(-K\cdot 0.009,0)=0$.  No contact force.  \checkmark
\end{examplebox}

\subsection{Extension: Adding Damping for Dynamic Stability}

In dynamic settings (as opposed to quasi-static manipulation), the spring-only
model~\eqref{eq:comfree-velocity} can exhibit oscillations at contact.
Jin adds a velocity-dependent \emph{damping} term:
\begin{equation}\label{eq:comfree-damped}
\bm{\lambda}_{\text{contact}}^{\text{ext}}=\max\!\Bigl(
-\bm{K}(\q)\bigl(\tilde{\J}\bm{Q}^{-1}\bm{b}+\tilde{\bm{\phi}}\bigr)
-\bm{D}(\q)\tilde{\J}\cdot\tfrac{1}{h}\bm{Q}^{-1}\bm{b},\;\0\Bigr),
\end{equation}
where $\bm{D}(\q)$ is a diagonal damping matrix.  The term
$\tilde{\J}\cdot\frac{1}{h}\bm{Q}^{-1}\bm{b}$ is the contact-frame velocity
caused by non-contact forces.

This gives a classical \emph{spring-damper} model in the dual cone, directly
analogous to MuJoCo's \texttt{solref} spring-damper in the primal Signorini space.

\subsection{ComFree-Sim: GPU Parallelisation}

Because the closed-form update~\eqref{eq:beta-closed} \emph{decouples across contacts}
(each $\beta_{i}^{+}$ depends only on local quantities), the computation maps
naturally to GPU kernels:
\begin{enumerate}
\item \textbf{Kernel I:} Compute smooth predicted velocity $\bm{v}_{\text{smooth}}^{+}=\bm{v}+\M^{-1}(\bm{\tau}-\bm{c})\,h$.
\item \textbf{Kernel II:} Per-contact, per-cone-facet: compute dual-cone impulse via~\eqref{eq:beta-closed}.
\item \textbf{Kernel III:} Accumulate generalised impulses $\bm{p}=\sum\tilde{\J}^{\!\top}\bm{\lambda}$.
\item \textbf{Kernel IV:} Apply velocity correction $\bm{v}^{+}=\bm{v}_{\text{smooth}}^{+}+\M^{-1}\bm{p}\,h$.
\end{enumerate}

This yields \emph{near-linear} runtime scaling with contact count, compared to
the super-linear scaling of iterative solvers.  Benchmarks show ComFree-Sim
achieves $2$--$3\times$ higher throughput than MuJoCo Warp in dense-contact scenes.

\subsection{The Dual-Cone Impedance: Connecting to MuJoCo and Drake}
\label{sec:impedance-comparison}

The complementarity-free model has two tunable parameter matrices: the stiffness
$\bm{K}(\q)$ and damping $\bm{D}(\q)$.  Understanding how these relate to
MuJoCo's \texttt{solref}/\texttt{solimp} and Drake's hydroelastic compliance
is essential for building intuition across simulators.

\begin{definitionbox}[Three impedance-based contact models compared]
All three approaches convert the hard complementarity into a soft
impedance model.  They differ in \emph{where} the impedance acts and
\emph{how} it is parameterised:

\renewcommand{\arraystretch}{1.3}
\begin{center}
\begin{tabular}{p{2.5cm}p{3.5cm}p{3.5cm}p{3.5cm}}
\toprule
& \textbf{MuJoCo} & \textbf{Drake} & \textbf{ComFree}\\
\midrule
\textbf{Where impedance acts}
& Primal space (Signorini: gap $\phi$ vs.\ force $\lambda_{N}$)
& Primal space (hydroelastic pressure field)
& \textbf{Dual-cone space} (dual-cone violation vs.\ $\bm{\beta}$)\\
\textbf{Parameters}
& \texttt{solref}$=[d,b]$, \texttt{solimp}$=[d_{\min},d_{\max},w]$
& Material stiffness, dissipation, hunt-crossley params
& $\bm{K}(\q)$ (stiffness), $\bm{D}(\q)$ (damping)\\
\textbf{Physical basis}
& Numerical (no material basis for $\bm{R}=\alpha\diag(\bm{G})$)
& Material properties (Young's modulus, Poisson ratio)
& Approximation of $({\tilde\J}\bm{Q}^{-1}{\tilde\J}^{\!\top}+\bm{R})^{-1}$\\
\textbf{Tangential compliance}
& Yes (relaxes Coulomb)
& Yes (less so)
& No (Coulomb satisfied by construction)\\
\textbf{Solver}
& Newton (iterative)
& Newton (iterative)
& \textbf{Closed-form (no iteration)}\\
\textbf{Signorini}
& Relaxed (allows penetration)
& Relaxed (pressure field)
& Dual-cone relaxed (mitigates CCP artifacts)\\
\bottomrule
\end{tabular}
\end{center}
\end{definitionbox}

\begin{insightbox}[Why ``impedance in the dual cone'' is different from ``impedance on the gap'']
MuJoCo's \texttt{solref} defines a spring-damper relating the \emph{gap} $\phi$
to the \emph{normal force} $\lambda_{N}$: deeper penetration $\Rightarrow$ larger force.
This operates in the \emph{primal} (physical) space.

ComFree's $\bm{K}$ defines a spring relating the \emph{dual-cone violation}
(a combined normal-and-tangential quantity) to the \emph{dual impulse} $\bm{\beta}$.
This operates in the \emph{dual} space of the optimisation problem.

The crucial consequence: ComFree's spring automatically respects the
Coulomb friction cone by construction (because $\bm{\beta}$ parameterises
forces inside the polyhedral cone), whereas MuJoCo must add separate
tangential compliance that \emph{relaxes} Coulomb.

Drake's hydroelastic model is also primal-space but is more principled
than MuJoCo: it derives compliance from material properties rather than
from the Delassus matrix, and computes pressure over an entire contact patch
rather than at a single point.
\end{insightbox}

\paragraph{Practical heuristic for $\bm{K}$ in ComFree-Sim.}
ComFree-Sim (Borse, Xie, Huang, Jin 2026) provides a practical parameterisation:
\begin{equation}
\bm{K}(\q)=k_{\text{user}}\cdot\bm{M}_{\phi}/h,\qquad
\bm{D}(\q)=d_{\text{user}}\cdot\bm{M}_{\phi}/h,
\end{equation}
where $k_{\text{user}},d_{\text{user}}$ are global scalar parameters
and $\bm{M}_{\phi}$ is a gap-dependent constraint-space inertia
using the same sigmoid profile as MuJoCo's \texttt{solimp}
(parameters $[r_{\min},r_{\max},w,m,p]$).
This ensures a MuJoCo-compatible interface while preserving the
analytical closed-form structure.

\begin{examplebox}[Worked Example: Sliding cube --- ComFree vs.\ CCP vs.\ NCP]
A $0.01$\,kg cube slides on a frictional plane ($\mu=0.5$) with initial
horizontal velocity $2.0$\,m/s, $\Delta t=2$\,ms.

\textbf{NCP (ground truth):} Cube decelerates at $a=-\mu g=-4.9$\,m/s$^{2}$,
stops after $t\approx 0.41$\,s.  Travels in a straight line.  No vertical motion.

\textbf{CCP (QP-based):} Dual-cone constraints force artificial vertical motion
(``boundary layer'' artifact). The cube bounces vertically while sliding.

\textbf{ComFree ($\bm{K}=\bm{I}$, $\bm{D}=0.3\bm{I}$):}
Because ComFree \emph{relaxes} the dual-cone constraints (allowing small violations),
the boundary-layer artifact is \textbf{mitigated}---the cube shows significantly
less vertical bouncing than CCP.  The horizontal velocity profile is smoother
than both CCP and MuJoCo.

This is a key advantage: by operating in the dual cone with a soft spring
rather than a hard constraint, ComFree trades strict dual-cone feasibility
for physically more realistic behaviour in the primal (position) space.
\end{examplebox}

\begin{examplebox}[Worked Example: Sphere between two planes (``sheath effect'')]
A sphere between two parallel planes is accelerated horizontally by
a $1$\,N force.  Friction coefficient $\mu=0.3$.

\textbf{CCP/Convex models:} The dual-cone constraints create vertical motion,
causing the sphere to contact the \emph{upper} plane.  This generates
opposing friction forces from both planes, producing a \emph{velocity plateau}
(``sheath effect'') where the sphere cannot accelerate past a limiting speed.

\textbf{NCP:} Correctly shows only lower-plane contact; sphere accelerates
continuously.

\textbf{ComFree:} Because dual-cone constraints are relaxed, the upper-plane
artifact is mitigated.  The plateau velocity is significantly higher than
CCP/Convex, approaching the NCP solution.
This demonstrates that relaxing dual-cone constraints can be
\emph{more physically accurate} than enforcing them.
\end{examplebox}

\subsection{Differentiability via SoftPlus}

The $\max(\cdot,0)$ in~\eqref{eq:beta-closed} is non-differentiable at zero.
For gradient-based trajectory optimisation, it is replaced with a smooth approximation:
\begin{equation}
\text{SoftPlus}(x)=\frac{\ln(1+e^{\gamma x})}{\gamma},\qquad\gamma>0,
\end{equation}
where $\gamma$ controls accuracy ($\gamma\to\infty$ recovers $\max(x,0)$).
Typical setting: $\gamma=100$.

With this substitution, the entire contact model becomes a smooth, differentiable
function $\bm{v}^{+}=\bm{f}(\q,\bm{v},\uu)$, directly usable in gradient-based MPC.
This is fundamentally different from MuJoCo's approach (which achieves differentiability
via compliance regularisation of an iterative solver) and from Drake's approach
(Newton solver with smooth compliance).

\subsection{Performance: Contact-Implicit MPC at 50--100 Hz}

Jin demonstrates that the complementarity-free model, integrated into
CasADi/IPOPT-based MPC, achieves:
\begin{itemize}
\item \textbf{96.5\% average success rate} across fingertip in-air manipulation,
      TriFinger in-hand manipulation, and Allegro hand on-palm reorientation
      with 17 diverse objects.
\item \textbf{MPC solving at 50--100\,Hz} for all tasks, compared to
      ${\sim}20$\,Hz for complementarity-based (``Implicit'') MPC.
\item \textbf{Average position error 7.8\,mm}, orientation error $11^{\circ}$.
\end{itemize}

ComFree-Sim (the GPU engine) further enables MPPI-based sampling MPC
on hardware, achieving $35$--$72$\,Hz closed-loop control on a LEAP Hand with
$2$--$3\times$ lower latency than MuJoCo Warp.

%-------------------------------------------------------------
\section{Comparison and Practical Implications}
\label{sec:contact-comparison}

\begin{center}
\begin{tabular}{lccccc}
\toprule
& \textbf{Hard contacts} & \textbf{No int.\ forces} & \textbf{Robust} & \textbf{Convergence} & \textbf{Closed-form}\\
\midrule
LCP/PGS & \checkmark & & & &\\
CCP/PGS & \checkmark & & & \checkmark &\\
CCP/MuJoCo & & \checkmark & \checkmark & \checkmark &\\
CCP/ADMM & \checkmark & \checkmark & \checkmark & \checkmark &\\
CCP/Drake & & \checkmark & \checkmark & \checkmark &\\
RaiSim & \checkmark & & & &\\
NCP/PGS & \checkmark & & & &\\
NCP/StagProj & \checkmark & \checkmark & \checkmark & &\\
\textbf{ComFree} & & \checkmark & \checkmark & \checkmark & \checkmark\\
\bottomrule
\end{tabular}
\end{center}

\begin{insightbox}[Impact on trajectory optimisation]
The choice of contact model directly impacts MPC and trajectory
optimisation:
\begin{enumerate}
\item On flat, non-slippery terrain, most solvers give similar results.
\item On bumpy, slippery terrain, the differences become pronounced:
      CCP models deviate from NCP, and PGS-based solvers fail to converge.
\item For differentiable simulation (needed for gradient-based trajopt
      through contact), the artificial compliance in MuJoCo produces
      gradients that differ from the ``true'' contact physics.
\item Complementarity-free models (Section~\ref{sec:comfree}) offer a
      compelling alternative: closed-form, differentiable, Coulomb-consistent
      contact with linear scaling---enabling real-time contact-implicit MPC
      at 50--100\,Hz for dexterous manipulation.
\item For problems where no differentiable contact model suffices,
      \emph{sampling-based} planners (Chapter~\ref{ch:sampling}) remain
      the most robust option: they require only forward simulation and
      are agnostic to the contact solver's internals.
\end{enumerate}
\end{insightbox}


%=============================================================================
%=============================================================================
\chapter{Grasping and Manipulation: Force Closure and Contact Wrenches}
\label{ch:grasping}

%-------------------------------------------------------------
\section{Motivation: The Physics of Grasping}

Up to this point, we have developed a unified view of contact as a set of forces constrained by friction cones, complementarity, and rigid-body dynamics. In manipulation, these same primitives appear again as the fundamental building blocks of \emph{grasping}. A grasp is simply a structured collection of contacts whose forces combine to produce object-level effects.

The central question of grasping is: \emph{What net forces and torques (wrenches) can a set of contacts generate on an object?}

%-------------------------------------------------------------
\section{From Contact Forces to Wrenches}

A single frictional contact point can exert a wrench (force + torque) on a grasped object. Consider a rigid body with a contact at point $\bm{p}_{i}$. A contact force $\bm{f}_{i} \in \R^{3}$ applied at that point produces both a linear force and a torque about the object's origin.

We define the contact wrench $\bm{w}_{i} \in \R^{6}$:
\begin{equation}
\bm{w}_{i} =
\begin{bmatrix}
\bm{f}_{i} \\
(\bm{p}_{i} - \bm{p}_{o}) \times \bm{f}_{i}
\end{bmatrix}
\end{equation}

This relationship is a linear map from force to wrench:
\begin{equation}
\bm{w}_{i} = \bm{G}_{i} \bm{f}_{i}
\end{equation}
where $\bm{G}_{i}$ is the grasp map for contact $i$:
\begin{equation}
\bm{G}_{i} =
\begin{bmatrix}
\I \\
[\bm{p}_{i} - \bm{p}_{o}]_{\times}
\end{bmatrix}
\end{equation}

Stacking all contacts together, the total wrench on the object is the sum of the individual wrenches:
\begin{equation}
\bm{w} = \sum_{i} \bm{w}_{i} = \sum_{i} \bm{G}_{i} \bm{f}_{i} = \bm{G} \bm{f}
\end{equation}
This is structurally identical to the force aggregation used in contact dynamics, with the addition of explicitly tracked torques.

%-------------------------------------------------------------
\section{Friction Cones and the Grasp Wrench Space (GWS)}

Each contact force is not arbitrary; it must lie within its respective friction cone: $\bm{f}_{i} \in \Kc_{\mu^{(i)}}$.

In the object frame, the set of wrenches achievable by contact $i$ is:
\begin{equation}
\mathcal{W}_{i}=\bigl\{[\bm{f}^{\!\top},\;\bm{\tau}^{\!\top}]^{\!\top}
\;\big|\;\bm{f}\in\Kc_{\mu^{(i)}},\;
\bm{\tau}=\bm{p}^{(i)}\times\bm{f}\bigr\},
\end{equation}
where $\bm{p}^{(i)}$ is the contact location relative to the object's centre of mass.

The \emph{grasp wrench space} (GWS) is the Minkowski sum of all contact wrench sets:
\begin{equation}
\text{GWS}=\mathcal{W}_{1}\oplus\mathcal{W}_{2}\oplus\cdots\oplus\mathcal{W}_{n_{c}}.
\end{equation}

\begin{insightbox}[Geometric Interpretation]
Because the GWS is the image of a product of cones under a linear map, it is itself a \emph{convex cone}. All results follow directly from the same friction and contact models used in dynamics.
\end{insightbox}

\begin{center}
\begin{tikzpicture}
  % Axes
  \draw[-{Stealth}, gray, thick] (-3.5,0) -- (3.5,0) node[right] {\footnotesize $w_x$};
  \draw[-{Stealth}, gray, thick] (0,-3) -- (0,3) node[above] {\footnotesize $w_y$};

  % GWS Polytope (Convex Hull)
  \coordinate (A) at (2.5, 1.5);
  \coordinate (B) at (1.0, 2.5);
  \coordinate (C) at (-2.0, 1.8);
  \coordinate (D) at (-2.5, -0.5);
  \coordinate (E) at (-1.0, -2.2);
  \coordinate (F) at (2.0, -1.8);

  \fill[orange, opacity=0.15] (A) -- (B) -- (C) -- (D) -- (E) -- (F) -- cycle;
  \draw[orange, very thick] (A) -- (B) -- (C) -- (D) -- (E) -- (F) -- cycle;
  
  \node[orange!80!black, above right] at (C) {\footnotesize $\text{GWS}$};

  % Origin
  \fill[black] (0,0) circle (2pt) node[below right] {\footnotesize $\0$};

  % Ferrari-Canny Metric Ball (Maximal inscribed circle)
  % Radius is the shortest distance from origin to the closest edge (approx 1.2 here)
  \draw[blue, very thick, dashed] (0,0) circle (1.2);
  \fill[blue, opacity=0.05] (0,0) circle (1.2);
  
  % Radius arrow for Q_FC
  \draw[-{Stealth}, blue, thick] (0,0) -- (-0.85, 0.85) node[midway, below left, inner sep=1pt] {\footnotesize $Q_{\text{FC}}$};

  % Note indicating force closure
  \node[black, right] at (1.5, 2.8) {\footnotesize $\0 \in \text{int}(\text{GWS}) \implies$ Force Closure};
\end{tikzpicture}
\end{center}

%-------------------------------------------------------------
\section{Force Closure vs.\ Form Closure}

\begin{definitionbox}[Force closure]
A grasp is in \emph{force closure} if the GWS contains the origin in its interior: $\0\in\text{int}(\text{GWS})$.
Equivalently, the grasp can resist any external wrench (up to some magnitude) through appropriate choice of contact forces within the friction cones.
\end{definitionbox}

\begin{itemize}
    \item \textbf{Force Closure:} Relies on friction to expand the achievable wrench space. It is the most common paradigm in practical robotic manipulation.
    \item \textbf{Form Closure:} A purely geometric condition where no motion is possible even with zero friction. Form closure strictly implies force closure, but not vice versa.
\end{itemize}

\paragraph{Positive span: the linear algebra of closure.}
Force closure is, at bottom, a statement about \emph{positive} linear
combinations. A set of vectors $\{\bm{w}_{1},\ldots,\bm{w}_{k}\}$
\emph{positively spans} $\R^{d}$ if every vector in $\R^{d}$ can be
written as $\sum_{j}\alpha_{j}\bm{w}_{j}$ with all $\alpha_{j}\ge 0$.
The sign restriction is the whole story: contacts can only \emph{push}
($\alpha_{j}\ge 0$ are forces along admissible directions), never pull,
so ordinary span is not enough. Two standard facts give quick sanity
checks:
\begin{itemize}
\item Positively spanning $\R^{d}$ requires at least $d+1$ vectors
      (with only $d$, their positive combinations fill at most a
      half-space---there is always an ``escape direction'' with negative
      inner product against all of them). Hence in the plane
      ($d=3$ wrench dimensions: $w_{x},w_{y},\tau$) at least 4 extreme
      wrench directions are needed, which two frictional contacts
      supply (2 cone edges each); in 3D ($d=6$) at least 7, which is
      why frictionless form closure needs 7 contacts but three
      frictional contacts (3--4 cone generators each) can suffice.
\item A test: $\{\bm{w}_{j}\}$ positively spans $\R^{d}$ iff the origin
      is in the \emph{interior} of their convex hull of directions---%
      which is exactly the $\0\in\text{int}(\text{GWS})$ condition.
      Equivalently, there is no separating hyperplane: no wrench
      direction $\bm{d}$ with $\bm{d}^{\!\top}\bm{w}_{j}\le 0$ for all
      $j$. A separating direction, if it exists, is precisely a
      disturbance the grasp cannot resist---the certificate of failure
      is itself physically meaningful.
\end{itemize}

\begin{definitionbox}[Nguyen's condition for planar antipodal grasps]
For two contacts in the plane, force closure holds iff the line segment
connecting the contact points lies \emph{strictly inside both friction
cones}. Each cone has half-angle $\arctan\mu$ about its inward normal,
so the condition is a simple angle check: the connecting line must make
an angle less than $\arctan\mu$ with each contact normal.
\end{definitionbox}

\begin{examplebox}[Worked Example: Nguyen's condition, numerically]
Grasp a disk of radius $r=4$\,cm with two fingers, with $\mu=0.4$, so
the friction half-angle is $\arctan 0.4\approx 21.8^{\circ}$.

\textbf{Case 1: perfectly antipodal contacts.} The connecting segment is
a diameter, aligned with both inward normals: angle $0^{\circ}<21.8^{\circ}$
at both contacts. Force closure. \checkmark

\textbf{Case 2: misplaced finger.} Slide one contact along the rim by
$30^{\circ}$ of arc. For a disk, the inward normal at each contact still
points at the centre, and a chord subtending a central angle of
$180^{\circ}-30^{\circ}=150^{\circ}$ makes an angle of
$(180^{\circ}-150^{\circ})/2=15^{\circ}$ with each radius.
Since $15^{\circ}<21.8^{\circ}$: still force closure---friction buys
substantial placement tolerance. \checkmark

\textbf{Case 3: how much misplacement is too much?} The chord angle
equals the half-cone angle when the arc offset reaches
$2\times 21.8^{\circ}=43.6^{\circ}$. Beyond that, the connecting line
exits one cone: squeezing harder no longer helps---the squeeze itself
generates a net tangential component that the other contact cannot
cancel, and the disk squirts out (every novice chopstick user has run
this experiment). With $\mu=0.1$ ($\arctan\mu\approx 5.7^{\circ}$) the
tolerance shrinks to ${\sim}11^{\circ}$ of arc: low-friction objects
demand precise antipodal placement.
\end{examplebox}

%-------------------------------------------------------------
\section{Quality Metrics and Computation}

The quality of a force-closure grasp can be quantified by the \emph{Ferrari--Canny metric}: the radius of the largest ball centred at the origin that fits inside the GWS.
\begin{equation}
Q_{\text{FC}}=\min_{\|\bm{w}\|=1}\max_{\bm{f}\in\prod\Kc_{\mu^{(i)}}}\bm{w}^{\!\top}\bm{G}_{\text{grasp}}\bm{f},
\end{equation}
where $\bm{G}_{\text{grasp}}$ is the grasp map relating contact forces to object wrenches. A larger $Q_{\text{FC}}$ means a more robust grasp.

For discrete (pyramidal) friction approximations, the GWS becomes a \emph{polytope} that can be computed via convex hull algorithms. This connects directly to the LCP linearisation (Section~\ref{sec:lcp-solver}): the pyramid friction model produces a polyhedral GWS, while the true Coulomb model produces a curved boundary.

%-------------------------------------------------------------
\section{Worked Examples}

\begin{examplebox}[Worked Example 1: Planar Pinch Grasp]
Consider a 2D object grasped by two opposing contacts. Each contact provides a friction cone (a wedge in 2D).
\begin{itemize}
    \item Without friction ($\mu = 0$), there is no force closure.
    \item With sufficient friction, the tangential forces allow the contact cones to expand the wrench space until it spans the full $\R^{3}$ planar wrench space.
\end{itemize}
\textbf{Insight:} Friction expands each contact cone, which expands the wrench space until it spans the full space.
\end{examplebox}

\begin{center}
\begin{tikzpicture}
  % Object
  \draw[fill=gray!20, thick, rounded corners=3mm] (-2,-1.2) rectangle (2,1.2);
  \fill[black] (0,0) circle (2pt) node[below right] {\footnotesize CoM ($\bm{p}_o$)};

  % Left contact (p1)
  \coordinate (P1) at (-2,0);
  % Friction cone 1
  \fill[blue, opacity=0.15] (P1) -- (-0.5, 1) -- (-0.5, -1) -- cycle;
  \draw[blue, thick] (P1) -- (-0.5, 1) node[above] {\footnotesize $\Kc_{\mu^{(1)}}$};
  \draw[blue, thick] (P1) -- (-0.5, -1);
  % Force vector 1
  \draw[-{Stealth}, very thick, red] (P1) -- (-0.8, 0.5) node[left] {$\bm{f}_1$};
  \fill[black] (P1) circle (2pt) node[left] {\footnotesize $\bm{p}_1$};

  % Right contact (p2)
  \coordinate (P2) at (2,0);
  % Friction cone 2
  \fill[blue, opacity=0.15] (P2) -- (0.5, 1) -- (0.5, -1) -- cycle;
  \draw[blue, thick] (P2) -- (0.5, 1) node[above] {\footnotesize $\Kc_{\mu^{(2)}}$};
  \draw[blue, thick] (P2) -- (0.5, -1);
  % Force vector 2
  \draw[-{Stealth}, very thick, red] (P2) -- (1.1, -0.4) node[right] {$\bm{f}_2$};
  \fill[black] (P2) circle (2pt) node[right] {\footnotesize $\bm{p}_2$};

  % Moment arms (dashed)
  \draw[dashed, gray!80, thick] (0,0) -- (P1);
  \draw[dashed, gray!80, thick] (0,0) -- (P2);
\end{tikzpicture}
\end{center}

\begin{examplebox}[Worked Example 2: Three-Point 3D Grasp]
Three contacts provide three friction cones in $\R^{3}$.
A common failure mode occurs if the contacts lie on a single line or plane, causing the torque directions to become linearly dependent. To achieve force closure, the set of mapped contact forces must positively span $\R^{6}$.
\end{examplebox}

%=============================================================================
% PART IV: PUTTING IT ALL TOGETHER
%=============================================================================
\part{Synthesis}

\chapter{Trajectory Optimisation Through Contact}
\label{ch:trajopt-contact}

%-------------------------------------------------------------
\section{The Challenge}

Contact-rich trajectory optimisation must simultaneously:
\begin{enumerate}
\item Plan motions that exploit contact (pushing, pivoting, grasping).
\item Handle non-smooth dynamics (making/breaking contact).
\item Solve for contact forces that are physically consistent.
\end{enumerate}

Derivative-based methods (iLQR/DDP) struggle because:
\begin{itemize}
\item Contact transitions create discontinuities in the dynamics Jacobian.
\item The contact solver's output (forces) is not a smooth function of the state, especially with hard contacts.
\item MuJoCo's compliance regularisation produces gradients, but they may not reflect the true physics.
\end{itemize}

Sampling-based methods are more robust because they only need
\emph{forward simulations}, not derivatives.  This is a major reason
MuJoCo MPC includes Predictive Sampling alongside iLQG.

%-------------------------------------------------------------
\section{Strategies for Contact-Rich Trajectory Optimisation}

\paragraph{1. Sampling-based MPC.}
Use MPPI, CEM, or annealed sampling (DIAL-MPC) with a physics
simulator in the loop.  No derivatives needed.
Works well for tasks where the contact sequence is not known a priori.

\paragraph{2. Smoothed contact models.}
Use compliant contact (as in MuJoCo) to smooth the dynamics, then
apply derivative-based methods.  The compliance acts as a regulariser
that removes the non-smoothness, at the cost of physical accuracy.

\paragraph{3. Through-contact trajectory optimisation.}
Treat contact forces as decision variables with complementarity
constraints.  This is a mathematical programming approach
(see Posa, Tedrake, etc.) that can find optimal contact sequences
but requires solving difficult complementarity-constrained programs.

\paragraph{4. Complementarity-free contact-implicit MPC.}
Use a closed-form contact model (Section~\ref{sec:comfree}) that bypasses
complementarity entirely.  The resulting MPC problem is a standard NLP
(no nested complementarity), solvable by off-the-shelf tools (CasADi/IPOPT)
at 50--100\,Hz.  Coulomb friction is satisfied by construction.
This approach achieves state-of-the-art results in dexterous manipulation
with diverse objects (Jin, 2025; Borse et al., 2026).

\paragraph{5. Hybrid approaches.}
Use a high-level planner (possibly RL-based) to determine contact
sequences, and a low-level MPC (MuJoCo MPC) to execute each phase.
This is the ``high-level agent'' concept from the MJPC discussion.

%-------------------------------------------------------------
\section{The MuJoCo MPC Approach in Practice}

For contact-rich tasks, MuJoCo MPC typically uses:
\begin{itemize}
\item Predictive Sampling (or iLQG with compliant contacts).
\item Spline-parameterised controls to reduce the search space.
\item Rich sensor-based cost functions with proximity/contact sensors.
\item Risk-sensitive costs to handle uncertainty from contact.
\item Simulation slow-down to give the planner more ``thinking time.''
\end{itemize}

The ``lean task'' (bracing against a table to reach a distant object)
and ``avoid task'' (collision avoidance using capacitive skin sensors)
from the MJPC extensions demonstrate how sensor-based costs,
computed ``for free'' (Section~\ref{sec:free-sensors}), enable
complex contact-aware behaviours.


%=============================================================================
\appendix
\chapter{Notation Reference}

\begin{longtable}{ll}
\toprule
\textbf{Symbol} & \textbf{Meaning}\\
\midrule
$\x\in\R^{n}$ & State vector\\
$\uu\in\R^{m}$ & Control vector\\
$\f(\x,\uu)$ & Discrete dynamics\\
$c(\x,\uu)$ & Running cost\\
$c_{T}(\x)$ & Terminal cost\\
$J$ & Total cost / objective\\
$V_{t}(\x)$ & Value function (cost-to-go)\\
$Q_{t}(\x,\uu)$ & State-action value function\\
$T$ & Horizon length\\
$\Delta t$, $h$ & Time step\\
$\q$ & Generalised coordinates\\
$\bm{v}$ & Generalised velocity\\
$\M$ & Joint-space inertia matrix\\
$\J$ & Contact Jacobian\\
$\bm{G}=\J\M^{-1}\J^{\!\top}$ & Delassus matrix\\
$\bm{\lambda}$ & Contact impulse / force\\
$\lambda_{N}$ & Normal component\\
$\bm{\lambda}_{T}$ & Tangential component\\
$\bm{c}$ & Contact point velocity\\
$\mu$ & Friction coefficient\\
$\Kc_{\mu}$ & Coulomb friction cone\\
$\bm{R}$ & Compliance matrix\\
$\rr$ & Residual (sensor) vector\\
$\rho(\cdot;R)$ & Risk-sensitive transformation\\
$\bm{\theta}$ & Spline parameters\\
$N$ & Number of rollouts / samples\\
$\sigma$ & Noise standard deviation\\
\bottomrule
\end{longtable}


\chapter{Further Reading}

\begin{enumerate}
\item Howell et al., ``Predictive Sampling: Real-time Behaviour Synthesis with MuJoCo,'' DeepMind, 2022.
\item Le~Lidec et al., ``Contact Models in Robotics: a Comparative Analysis,'' Inria, 2024.
\item Kelly, ``An Introduction to Trajectory Optimization,'' SIAM Review, 2017.
\item Tassa et al., ``Synthesis and stabilization of complex behaviors through online trajectory optimization,'' IROS, 2012.
\item Tassa et al., ``Control-limited differential dynamic programming,'' ICRA, 2014.
\item Williams et al., ``Model predictive path integral control,'' JGCD, 2017.
\item Castro, Permenter, Han, ``An unconstrained convex formulation of compliant contact,'' IEEE T-RO, 2022.
\item Todorov, Erez, Tassa, ``MuJoCo: A physics engine for model-based control,'' IROS, 2012.
\item Jin, ``Complementarity-Free Multi-Contact Modeling and Optimization for Dexterous Manipulation,'' RSS, 2025.
\item Borse, Xie, Huang, Jin, ``ComFree-Sim: A GPU-Parallelized Analytical Contact Physics Engine,'' arXiv:2603.12185, 2026.
\item Anitescu, ``Optimization-based simulation of nonsmooth rigid multibody dynamics,'' Math.\ Prog., 2006.
\end{enumerate}

\end{document}
